\chapter{学科与院校差异}
\label{ch20}
\ist{学科与院校差异}
\begin{keybox}
本章回答一个问题：不同学科之间的推免规则差在哪里。答案不是“医学重规培、工科重机试”这类概括，而是——差异几乎全部由\textbf{院系自己设定}。学校层面的办法通常只写通用条件，真正卡人的排名门槛、考核构成、专业白名单，写在学院、学部、附属医院那一级当年发布的细则里。

本章每一条都点名具体学校、具体学院与文件年份，并把可对照的条款做成表格。读的时候注意两件事。

第一，\textbf{以下差异均为院系自定，不可跨校套用}。同一年、同一座城市，两所学校对“医学影像学本科能否报临床专硕”的答复可能相反；同一所学校内部，文学院与计算机学院的复试构成也是两套。本章出现的比例与门槛，只对注明的那所学院、那个年度有效。

第二，三类内容分开写。全国统一规则（如推免名额自 2014 年起不再区分学术学位与专业学位）写在正文并标明“全国统一”；院校特定规则点明学校与年份；经验性建议放进“说明（经验性，非官方规则）”一栏。

部分方向本次未找到可对照的一手文件，已在相应条目如实标注并给出低置信度。证据不足处，本章说明边界，不补数字。
\end{keybox}

\section{院校类型-新型高校}
\begin{qa}{E0375}{西湖大学、上海科技大学、南方科技大学这类学校的推免怎么报？}
\textbf{先分清两种身份：西湖大学是 2025 年才拿到推免资格的“新晋推免高校”，上科大、南科大则未被本批证据覆盖。}

\par\vspace{2pt}\nobreak\begin{center}\footnotesize
\begin{tabularx}{\textwidth}{>{\raggedright\arraybackslash}X>{\raggedright\arraybackslash}X>{\raggedright\arraybackslash}X}
\toprule
\textbf{学校} & \textbf{推免身份} & \textbf{依据} \\ \midrule
西湖大学 & 2025 年新增推免资格高校 & 教学厅函〔2025〕20 号备案名单 \\
深圳大学 & 早有推免资格，正常收推免 & 医学部 2027 级公告、法学院 2027 预报名 \\
上海科技大学 & 未找到一手文件 & — \\
南方科技大学 & 未找到一手文件 & — \\
\bottomrule
\end{tabularx}
\end{center}\vspace{2pt}\par

西湖大学的“新”在于：它是 2025 年被教育部办公厅备案的 67 所新增推免资格高校之一，也就是说从 2025 年起它本校学生才具备往外推的资格；它与上科大、南科大一样属于研究型大学序列。你若是这三校学生，走的是同一套流程——先在本校拿到推免资格，再通过目标院校的夏令营或预推免、最后在教育部推免服务系统里确认待录取。
\nref{S0050}
\nref{S0436}
\nref{S0640}
\nref{S0265}
\nref{S0241}
\end{qa}

\iin{中华人民共和国教育部}
\ith{新型大学}
\ist{院校与项目差异}
\begin{qa}{E0376}{宁波诺丁汉、西交利物浦这类学校的学生能推免到国内高校吗？}
\textbf{能，但只有已获推免资格的那几所可以，且要走内地推免系统。}

关键分水岭是 2025 年：教育部办公厅以教学厅函〔2025〕20 号对 67 所新增推免资格高校予以备案，其中\textbf{宁波诺丁汉大学、香港中文大学（深圳）}在列。这意味着这两所中外合作办学高校从 2025 年起进入内地推免体系，本校学生可按常规路径申请内地高校。

需要注意两点。第一，资格是“学校有”不等于“你一定能拿到”，仍要在校内竞争名额，规则按各校当年办法执行。第二，跨校接收时对方是否认可，由接收单位自己的条件决定——多数内地院系在申请条件里写“经教育部批准具有推免资格的高校应届本科毕业生”（如中南大学法学院 2027 年通知），并不排斥中外合作办学高校，但具体仍以对方当年通知为准。

西交利物浦大学、上海纽约大学未出现在本批备案名单中，本手册也未找到其推免一手文件，不能据名单外推。
\nref{S0050}
\nref{S0436}
\nref{S0640}
\nref{S0237}
\end{qa}

\iin{中华人民共和国教育部}
\ith{中外合办}
\ist{院校与项目差异}
\begin{qa}{E0391}{2025 年新增的那批推免资格高校，学生能保研到哪些学校？}
\textbf{能报的学校和普通推免生完全一样，新增的是“你有了资格”，不是“你能去哪”。}

依据是教育部办公厅教学厅函〔2025〕20 号《关于对新增推免资格高校予以备案的通知》（2025 年 8 月 15 日）：经高校自主申报、所在地省级教育行政部门确认推荐并公示，对北京电子科技学院等 \textbf{67 所}新增推免资格高校予以备案。名单里既有西湖大学、宁波诺丁汉大学、香港中文大学（深圳）这类新型高校，也有江苏、浙江、山东、广东一批地方本科院校（如常州大学、苏州科技大学、浙江财经大学、广东财经大学、成都大学等）。

拿到资格后，你的去向由\textbf{接收方}的条件决定，而不是由“是否新增”决定。观察 2027 年度各校通知可以看到接收条件写法差异很大：

\par\vspace{2pt}\nobreak\begin{center}\footnotesize
\begin{tabularx}{\textwidth}{>{\raggedright\arraybackslash}X>{\raggedright\arraybackslash}X}
\toprule
\textbf{接收方} & \textbf{2027 年条件写法} \\ \midrule
武汉大学计算机学院 & 凡取得所在学校推免资格、身心健康即可接收 \\
西电计算机学院 & 仅要求获得本科就读高校推荐资格、无学术不端 \\
山东大学控制学院 & 其他有推免资格高校需专业排名前 8\% \\
\bottomrule
\end{tabularx}
\end{center}\vspace{2pt}\par

所以新增高校的学生既可能被“零门槛”院系接收，也会在分档院系那里吃亏。
\nref{S0050}
\nref{S0436}
\nref{S0640}
\nref{S0119}
\nref{S0120}
\nref{S0102}
\end{qa}

\iin{中华人民共和国教育部}
\ith{新增资格}
\ist{院校与项目差异}
\begin{qa}{S0306}{警校（如中国人民公安大学）推免需要满足哪些特殊条件？}
\textbf{未找到该方向的一手文件。}

本手册所依据的资料中没有中国人民公安大学或其他公安类院校的推免文件，无法给出经得起核查的特殊条件（体测、政审、入警政策衔接等）。为避免编造数字，此处不做推断。

可以确认的相邻事实有两条：一是“军队院校招收”类政策本手册只找到 2017 年威海市教育局关于军队院校招收普通高中毕业生的政策解读，属于高考招生而非推免，不能外推；二是公安技术类专业在体检上受限——教育部、卫生部、中国残疾人联合会《普通高等学校招生体检工作指导意见》明确“轻度色觉异常（色弱）不能录取的专业”包含公安技术类各专业，这对推免阶段的体检要求有参考意义，但并非公安院校推免的专门规定。

建议直接以中国人民公安大学研究生院当年发布的推免接收办法为准，并向招生办确认入警与推免的衔接安排。
\nref{S0039}
\nref{S0572}
\end{qa}

\iin{中华人民共和国教育部}
\ith{院校类型-警校}
\ist{院校与项目差异}
\begin{qa}{S0311}{双非院校中获得推免资格的比例有多低，怎么争取？}
\textbf{国家层面的比例低到 2\%，但“双非”内部差异极大，关键是先看你们学校属于哪一档。}

教育部教学〔2006〕14 号《全国普通高等学校推荐优秀应届本科毕业生免试攻读硕士学位研究生工作管理办法（试行）》第十二条按学校层次定名额：

\par\vspace{2pt}\nobreak\begin{center}\footnotesize
\begin{tabularx}{\textwidth}{>{\raggedright\arraybackslash}X>{\raggedright\arraybackslash}X}
\toprule
\textbf{学校类型} & \textbf{推免名额比例} \\ \midrule
教育部批准设立研究生院的高校 & 约应届本科毕业生数的 15\% \\
未设立研究生院的“211 工程”高校 & 约 5\% \\
其他高校 & 约 2\%（初次开展工作的前 3 年约 1\%） \\
人文、理科等人才培养基地班 & 按基地班招生人数 50\% 左右单独增加 \\
\bottomrule
\end{tabularx}
\end{center}\vspace{2pt}\par

也就是说，一所没有研究生院的普通本科院校，一个专业往往只有个位数名额，甚至隔年才有。校内竞争因此从大一开始。

争取路径要看你们学校当年的校内办法，门槛写法各不相同：成都大学商学院 2026 年细则要求本专业必修课程加权平均分的专业排名在前 20\%（含）；安徽师范大学 2027 年方案要求推免课程加权平均成绩排名前 35\% 且综合测评前 35\%；北航计算机学院 2027 年方案则要求学业综合成绩排名专业前 70\%。先拿到本校细则，再谈加分。
\nref{S0454}
\nref{S0079}
\nref{S0646}
\nref{S0579}
\nref{S0594}
\end{qa}

\iin{教育部学生服务与素质发展中心（学信网/研招网）}
\ith{院校类型-双非}
\ist{院校与项目差异}
\begin{qa}{S0324}{中外合作大学（宁诺、西浦、上纽）有推免资格吗？}
\textbf{有的有，有的没有——判断标准只有一个：是否出现在教育部 2025 年的备案名单里。}

教育部办公厅教学厅函〔2025〕20 号对 67 所新增推免资格高校予以备案，名单中明确包含\textbf{宁波诺丁汉大学}和\textbf{香港中文大学（深圳）}。这是中外合作办学高校首次成规模进入内地推免体系，两校学生自此可按常规路径在本校竞争名额、对外申请。

对照来看，你提到的三所里：

\par\vspace{2pt}\nobreak\begin{center}\footnotesize
\begin{tabularx}{\textwidth}{>{\raggedright\arraybackslash}X>{\raggedright\arraybackslash}X>{\raggedright\arraybackslash}X}
\toprule
\textbf{学校} & \textbf{2025 年备案名单中} & \textbf{说明} \\ \midrule
宁波诺丁汉大学 & 在列 & 已获推免资格 \\
香港中文大学（深圳） & 在列 & 已获推免资格 \\
上海纽约大学 & 不在列 & 本手册未收录其推免文件 \\
西交利物浦大学 & 不在列 & 本手册未收录其推免文件 \\
\bottomrule
\end{tabularx}
\end{center}\vspace{2pt}\par

需要提醒的是，“有资格”是学校层面的资质，不等于个人能拿到名额，也不等于内地高校一定接收。接收与否由对方院系的当年申请条件决定。名单之外的学校不能反推。
\nref{S0050}
\nref{S0436}
\nref{S0640}
\end{qa}

\iin{中华人民共和国教育部}
\ith{院校类型-中外合作}
\ist{院校与项目差异}
\begin{qa}{S0326}{独立学院、民办本科有推免资格吗？}
\textbf{独立学院学生基本没有推免资格，这不是院校政策而是全国统一规定的排除项；民办本科则要看该校是否被批准开展推免工作。}

教育部教学〔2006〕14 号第十四条在界定遴选对象时写得很明确：纳入国家普通本科招生计划录取的应届毕业生，\textbf{不含专升本、第二学士学位、独立学院学生}。这不是某一所学校自设的门槛，而是国家层面对“谁能被推荐”的范围界定。

各高校的校内办法都会照抄这一条。例如内蒙古大学实施办法第二条写明不含辅修学士学位学生、独立学院学生；广东外语外贸大学 2026 年实施办法第二条列明不含辅修学位、专升本、第二学士学位学生、本硕项目国际班学生和独立学院学生；湘潭大学 2026 年资格审查通知同样写“不含专升本、第二学士学位、职高对口、独立学院学生”。

至于民办本科：推免资格是按学校逐一审批公布的（第十二条之前的申请与公示程序），学校在名单内，其学生就有机会；不在名单内则完全没有。
\nref{S0454}
\nref{S0544}
\nref{S0134}
\nref{S0567}
\end{qa}

\iin{教育部学生服务与素质发展中心（学信网/研招网）}
\ith{院校类型-民办}
\ist{院校与项目差异}
\begin{qa}{S0331}{上海交通大学电子信息与电气工程学院的推免流程有什么特点？}
\textbf{未找到上海交通大学电子信息与电气工程学院的一手文件。}

本手册所依据的资料中有上海交通大学的多份推免材料，但没有该学院的推免通知，因此无法给出它的流程特点。

可以确证的是上海交通大学同层次院系的写法，供你校准预期：

\begin{itemize}
  \item \textbf{上海交通大学 2027 年招收推免生（含直博）办法}（研究生招生网）为全校层面的总办法；
  \item \textbf{上海交通大学计算机学院 2026 年夏令营通知}：面向“双一流”建设高校及建设学科的理工科专业本科三年级在校生，要求本科前五或前六学期专业排名\textbf{前 20\%}，C9 高校或教育部基础学科拔尖学生培养计划 2.0 基地拔尖班同学可适当放宽；英语原则上 CET6≥425 或 TOEFL≥90 或 IELTS≥6.0；营员分直博、直硕两类，直博报名时需选专业方向并填报意向导师；
  \item \textbf{船舶海洋与建筑工程学院 2027 年预推免办法}、\textbf{人工智能学院 2026 年夏令营通知}、\textbf{人文学院 2026 年优才夏令营通知}等为各学院自设。
\end{itemize}

这些材料能说明上交的院系普遍采取“校级总办法＋学院自设细则”两层结构，但电院的具体批次与考核方式仍需以电院当年通知为准。
\nref{S0515}
\nref{S0358}
\nref{S0575}
\nref{S0514}
\nref{S0319}
\end{qa}

\iin{上海交通大学}
\ith{院校特定-上交}
\ist{院校与项目差异}
\begin{qa}{S0332}{浙江大学计算机学院的推免考核环节和偏好是什么？}
\textbf{未找到浙江大学计算机学院的一手考核文件。}

所查资料中与浙大计算机学院直接相关的只有浙江大学 2026 年各学院（系）夏令营通知汇总页，它列出了《浙江大学计算机学院招收 2026 年暑期实习生计划》的链接，但正文未收录，因此无法还原它的考核环节与偏好。

可对照的是浙江大学同校其他院系的公开写法，用于理解浙大的院系级差异：

\par\vspace{2pt}\nobreak\begin{center}\footnotesize
\begin{tabularx}{\textwidth}{>{\raggedright\arraybackslash}X>{\raggedright\arraybackslash}X>{\raggedright\arraybackslash}X}
\toprule
\textbf{院系（年份）} & \textbf{排名门槛} & \textbf{英语门槛} \\ \midrule
网络空间安全学院（2026） & 综合成绩排名原则上本专业前 10\%，有出色学术成果者可放宽 & 六级 460 分及以上 \\
航空航天学院（2026） & 欢迎预期可获推免资格者 & 直博：六级 460 或托福 80 或雅思 5.5 \\
集成电路学院（2026） & 预期可获推免资格 & 直博：六级 460 或托福 80 或雅思 5.5 \\
医学院（2027） & — & 直博：六级 480，可放宽至 460 \\
\bottomrule
\end{tabularx}
\end{center}\vspace{2pt}\par

共同点是浙大直博普遍设 460 分以上的六级线，且托福雅思明确\textbf{不认可家考成绩与拼分成绩}。计算机学院是否沿用同一口径，需看其当年正式通知。
\nref{S0031}
\nref{S0216}
\nref{S0017}
\nref{S0215}
\nref{S0029}
\end{qa}

\iin{浙江大学}
\ith{院校特定-浙大}
\ist{院校与项目差异}
\begin{qa}{S0337}{北京航空航天大学、北京理工大学的推免偏好是什么？}
\textbf{北航有明确的一手文件，北理未找到。}

\textbf{北京航空航天大学计算机学院}（2027 年）把“校内遴选”和“对外接收”分成两份文件，很容易混淆：

\begin{itemize}
  \item \textbf{《计算机学院 2027 年推荐优秀应届本科毕业生免试攻读研究生工作实施方案》}：这是本校学生的推荐方案，基本条件之一是学业综合成绩排名在\textbf{专业前 70\%}；综合成绩＝（加权平均分＋综合加分）÷1.1；加分项单独设清单，学院成立专家组审核科研创新成果与竞赛奖项的真实性、原创性；在校期间参军入伍服兵役者经审核可酌情加分。
  \item \textbf{《北航计算机学院接收推荐免试攻读 2027 年研究生（含博士）工作方案》}：这是接收外校/本校学生的复试方案，附件含“C 语言考试成绩折算协议书”。
\end{itemize}

注意：专业前 70\% 是北航\textbf{给自己学生}定的推荐门槛，不是它\textbf{接收外校学生}的门槛，两者不要套用。

北京理工大学：本手册所依据的资料中未找到北理工的推免文件，无法给出其偏好。北航其他学院（如仪器科学与光电工程学院 2027 年）另发工作方案，说明北航是院系各自出细则。
\nref{S0579}
\nref{S0580}
\nref{S0421}
\nref{S0502}
\end{qa}

\iin{北京航空航天大学}
\ith{院校特定-北航北理}
\ist{院校与项目差异}
\begin{qa}{S0338}{华中科技大学、西安交通大学的保研友好度如何？}
\textbf{看“友好度”要分开看两件事：校内名额多不多、对外门槛高不高。}

\textbf{华中科技大学计算机科学与技术学院}：发布《2027 年推荐免试攻读硕士研究生工作方案及加分细则》，材料以附件形式给出（奖励加分细则、工作流程），公开页面本身不含条件正文，故现有公开材料不足以给出其具体门槛。

\textbf{西安交通大学}：以“校级汇总＋各学院细则”的形式发布——《西安交通大学各学院（部、中心）2027 年招收推免生实施细则》集中在研究生招生信息网公布，各学院的预报名通知则分散在各学院站点（如能动学院 2026 年环境科学与工程夏令营、经济与金融学院 2026 年夏令营、外国语学院 2026 年夏令营）。

所以两校的共同点是：\textbf{不存在一份统一的“保研友好度”标准}，你在华科要去看计算机学院的加分细则，在西交要逐学院去查细则页。

需要更正一个常见误解：本手册所依据的高校推免工作通知中，多份明确写着“除规定的专项计划外，各单位不得专门组织遴选推免的考试（包括笔试、面试）”（如中南民族大学 2027 年通知），这是\textbf{校内推荐环节}的规则，与目标院校的接收考核是两回事。
\nref{S0110}
\nref{S0545}
\nref{S0546}
\nref{S0147}
\nref{S0312}
\nref{S0317}
\nref{S0594}
\end{qa}

\iin{华中科技大学}
\ith{院校特定-华科西交}
\ist{院校与项目差异}
\begin{qa}{S0339}{电子科技大学、西安电子科技大学为什么被认为对双非友好？}
\textbf{这个说法可以从它们的公开条件里找到依据：两电的接收条件里没有写本科学校层次门槛，也没有排名百分比。}

先看西电：

\begin{itemize}
  \item 《西安电子科技大学 2027 年推荐免试研究生接收工作实施办法》：接收条件共 6 条，为中华人民共和国公民、拥护党的领导、身体健康、\textbf{获得本科就读高校推荐资格}、无学术不端与不良诚信记录、符合报考学院规定的其他条件；复试考核“以综合面试为主”，时间一般不少于 20 分钟，全程录音录像。
  \item 《西安电子科技大学 2027 年接收推荐免试研究生预报名通知》：基本要求与上述一致，未设排名或本科出身门槛。
  \item 计算机科学与技术学院 2026 年预报名通知：报名条件同样为 5 条基础条件，综合面试含外国语水平、专业知识与知识结构、科研创新意识、综合素质，不少于 20 分钟。
\end{itemize}

再看电子科技大学 2027 年接收通知：基本条件第 4 条只要求“获得目前就读学校推荐推荐免试攻读研究生资格”，考核为外语测试满分 100 分＋综合测试满分 200 分，总成绩达满分 60\% 为合格。

\textbf{这就是“对双非友好”的实证含义：门槛写在“有没有推免资格”，而不是“本科是哪一档”。}
\nref{S0157}
\nref{S0156}
\nref{S0120}
\nref{S0121}
\nref{S0495}
\end{qa}

\iin{西安电子科技大学}
\ith{院校特定-两电}
\ist{院校与项目差异}
\begin{qa}{S0340}{北京邮电大学、南京邮电大学的推免特点是什么？}
\textbf{未找到北京邮电大学、南京邮电大学的一手推免文件。}

本手册所依据的资料中没有北邮、南邮的推免通知或接收办法，无法给出这两校的推免特点，也不做推测。

可以确认的相关事实只有两条，都不足以支撑结论：

\begin{itemize}
  \item 教育部办公厅教学厅函〔2025〕20 号备案的 67 所新增推免资格高校名单中，包含\textbf{西安邮电大学}，但不含北邮、南邮（两校属于早已具有推免资格的院校序列，未出现在“新增”名单中是正常的，不能用名单外推其资格状态）。
  \item 所查资料中与“两电”（电子科技大学、西安电子科技大学）相关的文件丰富，但那两校的写法不能移植到两邮。
\end{itemize}

如果你要报北邮或南邮，建议直接查两校研究生院当年发布的推免接收办法与各学院实施细则，重点确认三点：是否设本科学校层次分档、是否设排名百分比、复试是纯面试还是含机试。
\nref{S0640}
\nref{S0050}
\end{qa}

\iin{中国政府网}
\ith{院校特定-两邮}
\ist{院校与项目差异}
\begin{qa}{S0341}{东南大学、同济大学的推免难度和偏好如何？}
\textbf{两校都不是“一个门槛”，而是学院各自出细则；本手册收录的材料足以看出这个结构。}

\textbf{同济大学（2027）}：校级层面发布《同济大学 2027 年接收推荐免试研究生（含直接攻博）预报名通知》，各学院另发细则，例如环境科学与工程学院 2027 年预报名通知、法学院 2027 年预报名通知；此外还有多个专项通道单独设通知，包括“上海期智研究院未来学者计划”专项接收推免直博生、“同济大学—北京中关村学院、上海创智学院、深圳河套学院”专项接收推荐免试直博生、“国优计划”推荐免试研究生。

\textbf{东南大学计算机科学与工程学院、软件学院（2025 年夏令营通知）}：申请条件为能在本科学校取得推免资格、研究兴趣浓厚、诚实守信、身心健康、英语水平良好，另有“学院其他要求”（外校学生需获得一位东南大学计算机科学相关老师推荐）。

两校的差异在于：同济把专项计划做成独立入口，东大则把“老师推荐”写进了学院附加条件。

需要说明的是，“推免难度”不是本手册所依据的指标，也无法从文件里客观得出，这里只给条件原文，不做难度排序。
\nref{S0484}
\nref{S0313}
\nref{S0242}
\nref{S0487}
\nref{S0488}
\nref{S0486}
\nref{S0126}
\end{qa}

\iin{同济大学}
\ith{院校特定-东南同济}
\ist{院校与项目差异}
\begin{qa}{S0344}{西安交通大学、西北工业大学的西部保研格局如何？}
\textbf{未找到可支撑“西部保研格局”判断的一手文件，只找到两校的常规推免材料。}

\textbf{西安交通大学}：以校级汇总页形式发布《西安交通大学各学院（部、中心）2027 年招收推免生实施细则》，另有 2024 年的预报名通知汇总；各学院自设通知分散发布，如经济与金融学院 2026 年（第十届）夏令营、能源与动力工程学院 2026 年（第十五届）夏令营、外国语学院 2026 年（第十届）夏令营。这些通知的共同结构是“学院介绍＋活动安排＋申请条件”。

\textbf{西北工业大学}：发布《西北工业大学 2027 年接收推荐免试研究生预报名的通知》（2026 年 9 月 4 日），申请基本条件为中华人民共和国公民、2027 届普通全日制本科毕业生且有望取得推免资格、具有科研潜力和创新意识、诚实守信学风端正、在校期间未受过任何纪律处分。

需要指出的是，所查资料中“西部”相关的唯一硬性条款来自国家文件——教学〔2006〕14 号第十七条规定设有研究生院的高校接收本校推免生不得超过本校推免生总数的 65\%，其中地处西部省份或军工、矿业、石油、地质、农林等特殊类型的高校可放宽但不得超过 75\%。这条是国家统一规定，不是西交或西工大自设的政策。
\nref{S0545}
\nref{S0546}
\nref{S0185}
\nref{S0312}
\nref{S0147}
\nref{S0317}
\nref{S0507}
\nref{S0454}
\end{qa}

\iin{西安交通大学}
\ith{院校特定-西北}
\ist{院校与项目差异}
\begin{qa}{S0348}{深圳大学、南方科技大学等新型高校有推免资格吗？}
\textbf{深圳大学有，南方科技大学未找到。}

\textbf{深圳大学}：本手册收录多份院系级推免接收文件，说明它长期具备推免资格并正常接收推免生——医学部《接收 2027 级推免研究生公告》（2026 年 9 月 2 日）、法学院 2027 年接收推荐免试研究生预报名通知、经济学院 2026 年接收推免生预报名公告。深大还自建了“深圳大学推免生预报名系统”（ehall.szu.edu.cn）作为报名入口。

\textbf{南方科技大学}：所查资料中仅在同济大学 2027 年环境科学与工程学院预报名通知里被提及，未找到南科大自身的推免办法或资格文件，因此无法确认其推免资格状态与流程。

\textbf{西湖大学}：属于另一个问题给出的信息——它是教学厅函〔2025〕20 号备案的 67 所新增推免资格高校之一，2025 年起获得资格。

需要提醒：深大是“早就具备推免资格的高校”，西湖大学是“2025 年新增”，二者在资格上不是同一类；南科大既不在新增名单中也未找到一手文件，不做判断。
\nref{S0265}
\nref{S0241}
\nref{S0136}
\nref{S0313}
\nref{S0436}
\nref{S0050}
\end{qa}

\iin{深圳大学}
\ith{院校类型-新型高校}
\ist{院校与项目差异}
\begin{qa}{S0349}{上海科技大学、西湖大学这类新型研究型大学的推免路径是什么？}
\textbf{两校材料很不均衡：西湖大学有资格层面的确证，上海科技大学没有一手文件。}

\textbf{西湖大学}：教学厅函〔2025〕20 号《教育部办公厅关于对新增推免资格高校予以备案的通知》（2025 年 8 月 15 日）将其列入 67 所新增推免资格高校备案名单。这意味着从 2025 年起西湖大学学生可以走常规推免路径。至于西湖大学\textbf{接收}外校推免生的路径，本手册未收录其接收办法。

\textbf{上海科技大学}：本手册所依据的资料中没有任何上科大的推免文件（既无资格文件也无接收办法），无法给出其推免路径，不做推测。

作为可参照的“新型研究型大学”对照，深圳大学医学部 2027 年公告展示了这类学校的典型做法：自建推免生预报名系统，申请人填报后由学部组织面试，考核环节以材料审核加面试为主，具体安排通过短信、邮箱、QQ 通知。

如果这两所是你的目标，最可靠的做法是查其研究生招生网当年的“接收推荐免试研究生”栏目，而不是参考同类学校的做法。
\nref{S0050}
\nref{S0436}
\nref{S0640}
\nref{S0265}
\end{qa}

\iin{中华人民共和国教育部}
\ith{院校类型-新型高校}
\ist{院校与项目差异}
\section{院校类型-师范}
\begin{qa}{E0379}{农林类院校的推免有什么特殊专项？}
\textbf{农林类有一条只写在国家文件里的特殊政策，本手册未收录农林院校自设的推免专项。}

\textbf{国家统一规定}：教育部教学〔2006〕14 号第十七条——设有研究生院的高等学校接收本校推免生的人数，不得超过本校推免生总数的 65\%；\textbf{其中地处西部省份或军工、矿业、石油、地质、农林等特殊类型的高等学校，上述比例可适当放宽，但不得超过 75\%}。这是“农林”出现在推免制度里的唯一一处硬性条款，且它管的是\textbf{本校接收本校推免生的上限}，不是给学生的加分或专项名额。

\textbf{农林院校的实践}：本手册收录华南农业大学经济管理学院 2027 年接收推荐免试硕士研究生预通知、北京林业大学外语学院 2026 年夏令营招生简章、福建农林大学 2026 年硕士研究生招生简章、甘肃农业大学 2027 年遴选推荐工作管理办法等，均为常规推免材料，其中未见“农林专项”的独立通道。

与农林真正相关的国家专项是研究生支教团、硕师计划（面向农村学校）、退役大学生士兵计划等，这些并非农林院校专属。
\nref{S0454}
\nref{S0087}
\nref{S0023}
\nref{S0413}
\nref{S0059}
\end{qa}

\iin{教育部学生服务与素质发展中心（学信网/研招网）}
\ith{农林专项}
\ist{院校与项目差异}
\begin{qa}{S0294}{教育学、师范类保研的考核重点和其他学科有什么不同？}
\textbf{差别不在“考不考教育学知识”，而在考核形式与专项通道：师范类保研大量走“计划型”通道，而不是单纯比排名。}

一手文件中能看到三个层次：

\textbf{一、专项计划通道（国家层面）}。国办发〔2024〕27 号确立部属师范大学本研衔接师范生公费教育（本科 4 年＋教育硕士 2 年）；教师〔2021〕4 号设优师计划；教师〔2009〕5 号设硕师计划。这些通道的录取不只看排名，还看从教意愿、教师资格证、履约地选择。

\textbf{二、院校层面的名额倾斜}。浙江师范大学 2027 年推荐工作通知给出具体结构：总名额 460 名，其中直接攻博 22 名、\textbf{卓越中学教师培养计划 80 名}（学科教学“地理科学＋思想政治教育”40 名、“化学＋生物科学”40 名，且该计划仅限培养高中师资的师范类专业学生报名）、其它 358 名；另有支教团 29 名、硕师计划 37 名。这说明师范院校的推免名额被专项计划切走了一大块。

\textbf{三、考核形式}。东北师范大学 2027 年招生录取工作办法把“国优计划”作为独立专项计划在报名系统中单独选择；华中师范大学 2027 年公费师范生免试攻读教育硕士招生办法按学科教学（思政、语文、数学、物理、化学、生物、英语、历史、地理、音乐等）逐专业对口设置本科专业要求。
\nref{S0355}
\nref{S0456}
\nref{S0543}
\nref{S0525}
\nref{S0569}
\nref{S0607}
\nref{S0559}
\end{qa}

\iin{教育部学生服务与素质发展中心（学信网/研招网）}
\ith{学科差异-教育}
\ist{院校与项目差异}
\begin{qa}{S0304}{核工程、兵器、航空航天类保研有什么特殊通道？}
\textbf{线索集中在两个方向：军队系统的“参军入伍生”通道，以及高校的国防特色学院。}

\textbf{国防科技大学 2027 年接收方案}（接收地方院校推免生，含直博生、参军入伍生）给出最明确的一条：

\begin{itemize}
  \item 招生类型分推免硕士、直博生、参军入伍生；
  \item \textbf{参军入伍生}须符合军队直接选拔招录军官条件，本科所在院校为“双一流”建设高校，本科专业为本科所在院校法学、文学、理学、工学、管理学门类“双一流”建设学科或 A 类学科的对应专业；
  \item 直博生要求同硕士生一致并具备博士培养潜质，\textbf{招生人数一般不超过当年博士招生计划的 30\%}；另设“重大工程联合培养”专项计划，不受博士招生计划比例限制；
  \item 复试按“优秀”“合格”“不合格”三类等次确定结果，各等次人数比例由各招生学院自行确定。
\end{itemize}

\textbf{高校端}：浙江大学航空航天学院 2026 年夏令营、北航仪器科学与光电工程学院 2027 年推免工作方案等，属于相关学科的常规推免材料，未见“核工程/兵器”专属通道。

结论：最实质的特殊通道是\textbf{参军入伍生}，它把“双一流学科/A 类学科”写成了硬条件。
\nref{S0650}
\nref{S0017}
\nref{S0421}
\end{qa}

\iin{国防科技大学}
\ith{学科差异-国防}
\ist{院校与项目差异}
\begin{qa}{S0309}{211院校的保研名额比例为什么普遍低于985？}
\textbf{不是“为什么”，是文件明文写死的两档比例。}

教育部教学〔2006〕14 号第十二条把名额比例按学校层次分成三档：

\par\vspace{2pt}\nobreak\begin{center}\footnotesize
\begin{tabularx}{\textwidth}{>{\raggedright\arraybackslash}X>{\raggedright\arraybackslash}X}
\toprule
\textbf{学校类型} & \textbf{名额比例} \\ \midrule
教育部批准设立研究生院的高等学校 & 一般按应届本科毕业生数的 \textbf{15\% 左右} \\
未设立研究生院的“211 工程”建设高等学校 & 一般按应届本科毕业生数的 \textbf{5\% 左右} \\
其他高等学校 & 一般按 \textbf{2\%}（初次开展工作者前 3 年每年约 1\%） \\
\bottomrule
\end{tabularx}
\end{center}\vspace{2pt}\par

关键在“设立研究生院”这五个字。很多 211 高校并未设立研究生院，因此落在 5\% 这一档；而设立研究生院的高校名单里既有 985 也有少量 211，所以真正的分界线是\textbf{有没有研究生院}，不是单纯 985/211 的标签。

另外两条容易被忽略：经教育部确定的人文、理科等人才培养基地的高校，按基地班招生人数的 50\% 左右单独增加名额；对发展急需专业适当增加名额。且第十二条末尾写明“教育部可根据研究生教育改革与发展的形势，对上述比例做适当调整”——这是 2006 年文件的口径，最新年度以当年下达的名额为准。
\nref{S0454}
\end{qa}

\iin{教育部学生服务与素质发展中心（学信网/研招网）}
\ith{院校类型-211}
\ist{院校与项目差异}
\begin{qa}{S0310}{「双一流」新入选高校的保研认可度提升了吗？}
\textbf{本手册所依据的文件不支持“认可度”这类主观判断，但可以用“是否写进申请条件”来观察。}

一个可用的替代指标是：目标院系会不会把学校层次写进硬条件。对照 2026 年度两份材料：

\par\vspace{2pt}\nobreak\begin{center}\footnotesize
\begin{tabularx}{\textwidth}{>{\raggedright\arraybackslash}X>{\raggedright\arraybackslash}X}
\toprule
\textbf{院系（年份）} & \textbf{对本科出身的要求} \\ \midrule
上海财经大学计算机与人工智能学院（2026） & “双一流”建设高校专业成绩名列前茅，\textbf{或其他高校}专业成绩或综合成绩在所学专业名列前 10\% \\
武汉大学计算机学院（2027） & 双非无差别：“凡取得所在学校推免生资格、身心健康的考生，均可接收” \\
\bottomrule
\end{tabularx}
\end{center}\vspace{2pt}\par

前者说明：写“双一流”时往往同时给出“其他高校”的替代档（这里是前 10\%），并非一刀切排除。后者说明：有的院系根本不写本科出身。

另外，教育部办公厅教学厅函〔2025〕20 号在 2025 年一次性备案了 67 所新增推免资格高校，说明推免资格盘子本身在扩大，但这属于资格准入变化，与“认可度提升”是两件事，不宜混为一谈。

本手册未收录任何院校对“双一流新入选高校”的差别化评价条款，故不做认可度断言。
\nref{S0302}
\nref{S0119}
\nref{S0050}
\nref{S0436}
\end{qa}

\iin{计算机与人工智能学院}
\ith{院校类型-双一流}
\ist{院校与项目差异}
\begin{qa}{S0317}{师范类院校（六所部属）保研的政策共性和差异是什么？}
\textbf{共性来自一份国家文件，差异体现在各校自定的接收办法与专业对口规则上。}

\textbf{共性}：国办发〔2024〕27 号转发的《教育部直属师范大学本研衔接师范生公费教育实施办法》是六所部属师范大学共同的制度基础——从 2024 年起实施本研衔接师范生公费教育（本科 4 年、教育硕士 2 年），支持符合条件的公费师范生\textbf{免试攻读本校全日制教育硕士研究生}再履约任教；重点为中西部省会城市之外的地（市、州、盟）及以下区域培养研究生层次中小学教师；研究生一年级课程结束后按本科以来综合考核结果排序，在公布的本省履约任教地范围内选择。

\textbf{过渡安排}：教育部办公厅 2025 年通知明确，《实施办法》印发前已签协议的 2025—2027 届本科层次公费师范生，本科三年级结束后可申请接受本研衔接，须\textbf{经学校推免综合考核、取得学校推免资格并达到学校硕士录取标准}，再报生源省份省级教育行政部门同意并签署补充协议。

\textbf{差异}：华中师范大学 2027 年公费师范生免试攻读教育硕士招生办法按学科教学各领域逐一对口本科专业并指定招生单位（如学科教学思政对应思想政治教育、由马克思主义学院培养）；东北师范大学 2027 年办法则是公费师范毕业生申请免试攻读\textbf{非全日制}教育硕士，要求履约任教满一学期、经任教学校考核合格并批准。

即入口统一、出口各校自定。
\nref{S0355}
\nref{S0456}
\nref{S0044}
\nref{S0559}
\nref{S0609}
\end{qa}

\iin{教育部学生服务与素质发展中心（学信网/研招网）}
\ith{院校类型-师范}
\ist{院校与项目差异}
\begin{qa}{S0318}{地方师范院校的保研名额和去向有什么特点？}
\textbf{地方师范院校的推免名额被专项计划切得很细，去向集中在省内基础教育。}

以浙江师范大学 2027 年推荐工作通知为例，名额结构一目了然：

\par\vspace{2pt}\nobreak\begin{center}\footnotesize
\begin{tabularx}{\textwidth}{>{\raggedright\arraybackslash}X>{\raggedright\arraybackslash}X}
\toprule
\textbf{计划类型} & \textbf{名额} \\ \midrule
总名额 & 460 名 \\
其中直接攻博 & 22 名 \\
其中卓越中学教师培养计划 & 80 名（地理＋思政 40、化学＋生物 40） \\
其中其它 & 358 名 \\
研究生支教团专项 & 29 名 \\
硕师计划 & 37 名 \\
\bottomrule
\end{tabularx}
\end{center}\vspace{2pt}\par

\textbf{去向特点}：卓越中学教师培养计划\textbf{仅限培养高中师资的师范类专业学生报名}，毕业时获教育硕士专业学位证书；硕师计划则指向县镇及以下农村学校。

再看两条更基层的实例：吉林师范大学 2027 年共获推免指标 80 个，其中常规硕师计划 76 个、优师计划专项 4 个，且从 2027 年起省级主管部门将优师计划与硕师计划结合实施；云南师范大学 2025 年硕师计划推荐 23 名、按云南省师资培养需求接收 14 名，推荐条件为专业学习成绩总排名前 60\%（无专业课不及格）且毕业时必须取得教师资格证。

共同点是：名额总量小、专项占比高、履约与定向色彩强。
\nref{S0569}
\nref{S0593}
\nref{S0376}
\nref{S0543}
\end{qa}

\iin{浙江师范大学}
\ith{院校类型-师范}
\ist{院校与项目差异}
\begin{qa}{S0320}{财经类院校保研的内部竞争和内卷程度如何？}
\textbf{财经类的“内卷”可以从两处硬条件看出来：接收端对本科出身的分档，和招生端对硕博贯通的高度倾斜。}

\textbf{接收端分档}：西南财经大学国际商学院 2027 级工作细则规定，考生须为“双一流”建设高校应届毕业生，\textbf{或本科专业（不含辅修专业）所属一级学科在教育部第四轮学科评估排名为“B”以上（含）或前 30\%（含）}；并限定“非经济类或管理类专业毕业生接收比例一般不超过总接收人数的 20\%”。

\textbf{招生结构}：上海财经大学金融学院 2027 年预报名通知写明，本次选拔针对直接攻博、学术学位硕士（\textbf{仅招硕博连读}）和硕士专业学位等项目。也就是说，想读学硕只有硕博连读一条路，“只读两年硕士”的空间被压缩——这是财经类内卷的结构性来源之一。

\textbf{名额稀缺的极端例子}：西南财经大学国际商学院 2027 年接收直博生的名额为国际贸易学 1 人、国际商务 1 人，且报名前须与博导联系并获得同意与推荐。

再看规模，上海财经大学经济学院 2027 年通知、广东财经大学 2027 年预报名公告（该校为 2025 年新增推免资格高校）等，均以此类推。
\nref{S0320}
\nref{S0330}
\nref{S0483}
\nref{S0138}
\nref{S0503}
\nref{S0436}
\end{qa}

\iin{西南财经大学国际商学院}
\ith{院校类型-财经}
\ist{院校与项目差异}
\begin{qa}{S0321}{政法类院校保研的特点是什么？}
\textbf{政法类最鲜明的特点是“院校出身”被写进申请资格，这在其他学科里很少见。}

\textbf{上海交通大学凯原法学院 2026 年第十六届“卓越法律人夏令营”}的申请资格写得很直接：法学硕士营、法律硕士（法学）营、本科直博营申请人\textbf{须为全国双一流建设高校及建设学科、传统优势政法院校的 2027 届法学专业本科毕业生}；法律硕士（非法学）营申请人须为全国双一流建设高校及建设学科 2027 届\textbf{非法学专业}本科毕业生。英语要求为 CET6≥425 或 TOEFL≥90 或 IELTS≥6.0，特别优秀者第 3、4 条可适当放宽。

这条“传统优势政法院校”的表述，等于给五院四系类院校的学生开了一个与双一流并列的入口——这是政法类独有的写法。

\textbf{对比其他院校}：中南大学法学院 2027 年考核录取通知的申请条件只写“经教育部批准具有推荐免试攻读硕士学位研究生资格高校的 2027 届优秀本科毕业生”，不设出身分档；湖南大学法学院 2027 年预报名通知亦同。

\textbf{专业方向设置}：中南大学法学院 2027 年招收法学学术学位（含法学理论、宪法学与行政法学、刑法学、民商法学等）与法律（非法学）、法律（法学）专业学位，非法学背景只能报非法学法硕。
\nref{S0240}
\nref{S0237}
\nref{S0238}
\end{qa}

\iin{上海交通大学}
\ith{院校类型-政法}
\ist{院校与项目差异}
\begin{qa}{S0322}{民族类院校保研有哪些特殊政策和限制？}
\textbf{民族类的“特殊政策”主要在国家专项，院校层面反而更强调程序规范。}

\textbf{国家统一的专项}是“少数民族高层次骨干人才计划”。教育部《2026 年全国硕士研究生招生工作管理规定》（教学〔2025〕2 号）第四十八条明确，该计划与少数民族照顾政策、工商管理等管理类专业学位及有关照顾专业\textbf{单独划线}；中山大学 2027 年接收推免通知进一步说明，推免生招生类别一般为非定向就业，但录取为“少数民族高层次骨干人才计划”等国家专项计划的为\textbf{定向就业}。

\textbf{院校层面的民族类实例}：中南民族大学 2027 年推荐工作通知（依据教育部办公厅教学厅〔2026〕1 号《全国高等学校推荐优秀应届本科毕业生免试攻读研究生工作管理规定》与民大教学〔2024〕49 号管理办法）要求各学院在推荐过程中全面考查、以德为先；学院须成立\textbf{至少 5 名同行专家（副教授以上职称）}组成的材料审核专家小组，对竞赛获奖、学术论文及发明专利进行审核鉴定，并\textbf{组织公开答辩、全程录音录像}；同时明确“除规定的专项计划外，各单位不得专门组织遴选推免的考试（包括笔试、面试）”。

内蒙古大学的实施办法则在加分上做民族地区适配：科研论文、竞赛获奖的学分绩点加分分别设 0.15、0.12、0.1 的上限。
\nref{S0051}
\nref{S0166}
\nref{S0594}
\nref{S0544}
\end{qa}

\iin{中华人民共和国教育部}
\ith{院校类型-民族}
\ist{院校与项目差异}
\begin{qa}{S0323}{艺术类院校保研名额少、跨校难度大，具体难在哪？}
\textbf{难处不在名额数字，而在两条：专业目录窄，且外语门槛被单独下调说明基础普遍偏弱。}

\textbf{名额少}：这一点可以从 2025 年新增备案名单侧面看出——67 所新增推免资格高校中，艺术类院校集中出现，包括天津美术学院、沈阳音乐学院、广州美术学院、四川美术学院、四川音乐学院、沈阳体育学院等（教学厅函〔2025〕20 号）。一次性放进来这么多艺术类院校，恰恰说明此前拥有推免资格的艺术院校数量极少。

\textbf{外语门槛单列}：多所院校对艺术类学生单独降低英语要求，这是其他学科没有的写法：

\par\vspace{2pt}\nobreak\begin{center}\footnotesize
\begin{tabularx}{\textwidth}{>{\raggedright\arraybackslash}X>{\raggedright\arraybackslash}X>{\raggedright\arraybackslash}X}
\toprule
\textbf{学校（年份）} & \textbf{普通专业} & \textbf{艺术/体育类专业} \\ \midrule
安徽工业大学（2027 支教团） & 四级达报考六级要求分数 & 艺术设计类专业\textbf{不作硬性要求} \\
安徽师范大学（2027） & 四级不低于 425 分 & 四级不低于 355 分，小语种不低于 50 分 \\
湘潭大学（2026） & 四级 425 分以上 & 艺术类专业 400 分以上 \\
\bottomrule
\end{tabularx}
\end{center}\vspace{2pt}\par

\textbf{加分通道}：湖南大学 2027 年对艺术类竞赛国家级及以上奖励单独公示加分名单，创新创业赛事国家级特等奖排名第 1 加 6 分、三等奖排名第 1 加 2 分。

跨校难度则与专业目录直接相关——艺术类学科在综合类大学里硕士点本就稀少。
\nref{S0436}
\nref{S0050}
\nref{S0589}
\nref{S0646}
\nref{S0567}
\nref{S0541}
\end{qa}

\iin{教育部学生服务与素质发展中心（学信网/研招网）}
\ith{院校类型-艺术}
\ist{院校与项目差异}
\begin{qa}{S0342}{中山大学、华南理工大学的保研地域优势体现在哪？}
\textbf{地域优势的证据在招生结构与批次安排上，而不在口号式的“地缘”。}

\textbf{中山大学 2027 年推荐优秀应届本科毕业生招生录取工作办法}：所有专业均招收硕士推免生；\textbf{所有学术学位专业，以及除 1051 临床医学和 1052 口腔医学之外的专业学位，均招收直博生}。直博生学制 5 年，全日制硕士学制一般 2\textasciitilde{}3 年。可跨学科申请攻读硕士、博士学位研究生。报考条件里明确引用了《全国高等学校推荐优秀应届本科毕业生免试攻读研究生工作管理规定》（教学厅〔2026〕1 号）。

\textbf{华南理工大学 2027 年接收推免生预报名通知}：全校 2027 年拟招收攻读硕士学位研究生约 6800 人、博士学位研究生约 1900 人（最终以教育部下达计划为准），并自建“华南理工大学接收推荐免试研究生预报名系统”作为入口。

\textbf{批次安排}：华南理工大学《关于公布各院（系）2027 年研究生招生预推免接收工作复试方案的通知》说明，各院（系）分批审核，\textbf{第一批复试名单筛选截止报名时间为 9 月 15 日}，各院系视第一批接收情况决定是否进行第二批复试。

这两点合起来构成“地域优势”的可核对内容：招生规模大、直博口径宽、批次多、报名入口独立。
\nref{S0167}
\nref{S0166}
\nref{S0641}
\nref{S0479}
\nref{S0480}
\end{qa}

\iin{中山大学}
\ith{院校特定-中山华工}
\ist{院校与项目差异}
\begin{qa}{S0345}{中国海洋大学、山东大学的特色学科保研有什么特点？}
\textbf{两校材料的分布很不均衡：山东大学找到的是“学院各自出通知”，中国海洋大学找到的是“学校统一口径”。}

\textbf{山东大学（2026）}：以学院为单元发布夏令营通知，本手册收录多份，覆盖不同学科——数学学院 2026 年全国优秀大学生夏令营、物理学院 2026 年“物理学优秀大学生暑期夏令营”、软件学院 2026 年夏令营、密码科学与工程学院 2026 年夏令营、公共卫生学院 2026 年夏令营、控制科学与工程学院 2026 年夏令营，以及经济研究院 2026 年暑期夏令营。其中控制学院的通知把排名门槛按本科学校层次分成三档（985 前 20\%／211 前 10\%／其他有推免资格高校前 8\%），是所查文件中最完整分档实例。

\textbf{中国海洋大学（2027）}：以学校层面统一发布——《2027 年接收推荐免试研究生预报名通知》说明“凡按规定可接收应届本科毕业生报考的学科（类别）、专业（领域）均可接收推免生”，且\textbf{各一级学科博士点（含专业学位博士点）下的专业均可接收本科直博生}；另有《关于 2027 年推免综合考核结果查询及报考提示》《2027 年部分单位推荐免试研究生第二轮报名通知》（公布海洋生命学院、工程学院、经济学院等单位的第二轮通知链接）。

结构差异本身就是可对照的信息：山大看学院，海大看学校统一页。
\nref{S0037}
\nref{S0385}
\nref{S0387}
\nref{S0125}
\nref{S0335}
\nref{S0102}
\nref{S0026}
\nref{S0758}
\nref{S0759}
\nref{S0760}
\end{qa}

\iin{山东大学}
\ith{院校特定-山东}
\ist{院校与项目差异}
\begin{qa}{S0347}{苏州大学、郑州大学等地方双一流的保研情况如何？}
\textbf{两校的共同点是学院各自出细则，差异在跨专业政策与保留入学资格的态度。}

\textbf{苏州大学 2027 年接收推荐免试研究生（含直博生）工作办法}：

\begin{itemize}
  \item 接收范围最宽——\textbf{所有全日制硕士研究生招生专业均可接收推免生}，部分博士招生专业招收本科直博生；
  \item 申请条件只有 4 条，要求取得所在高校推荐资格、诚实守信、学风端正、学术研究兴趣浓厚；
  \item 申请本科直博者须在复试前提供\textbf{两名所报考学科专业领域内教授（或相当专业技术职务专家）的书面推荐意见}；
  \item 面试为复试考核的必要环节，各单位可多批次组织复试，综合考核成绩百分制、60 分为合格。
\end{itemize}

\textbf{郑州大学 2027 年接收推荐免试研究生（含直博生）工作办法}：

\begin{itemize}
  \item 接收推免硕士生与直博生两类；
  \item 直博生应具有扎实基础、较强自学与实践创新能力，具有博士生培养潜质，\textbf{录取类别应为全日制非定向}；
  \item 明确“\textbf{除特殊推荐类型外，原则上不接收需要保留入学资格的推免生}”；
  \item 跨专业推免生所学专业必须与接收专业相近或相关。
\end{itemize}

差别很具体：苏大给直博加“两名教授推荐信”这道程序，郑大则卡住“保留入学资格”和“跨专业相近性”。
\nref{S0599}
\nref{S0635}
\end{qa}

\iin{苏州大学}
\ith{院校特定-地方双一流}
\ist{院校与项目差异}
\section{学科差异-计算机}
\begin{qa}{E0360}{双非学生保 985 有多难？哪些学校对双非比较友好？}
\textbf{难点可以量化：一部分院系把双非写进了分档，另一部分完全不看本科出身。}

\par\vspace{2pt}\nobreak\begin{center}\footnotesize
\begin{tabularx}{\textwidth}{>{\raggedright\arraybackslash}X>{\raggedright\arraybackslash}X}
\toprule
\textbf{会卡出身的院系（2026—2027）} & \textbf{双非需达到的条件} \\ \midrule
山东大学控制科学与工程学院（2026 夏令营） & 其他有推免资格高校，专业排名前 8\% \\
天津大学计算机科学与技术学院（2026 夏令营） & 未设研究生院高校或国家重点学科院校，排名前 10\% \\
上海财经大学计算机与人工智能学院（2026） & 其他高校专业成绩或综合成绩在所学专业前 10\% \\
中国药科大学（2027 直博） & 非国家级一流本科专业考生，排名前 8\% \\
\bottomrule
\end{tabularx}
\end{center}\vspace{2pt}\par

\textbf{不看出身的院系}：武汉大学计算机学院 2027 年实施细则的基本条件只有“凡取得所在学校推免生资格、身心健康的考生，均可接收”；西安电子科技大学 2027 年接收实施办法只要求“获得本科就读高校推荐资格”；电子科技大学 2027 年接收通知第 4 条同样只要求获得推免资格。

所以“双非保 985 有多难”取决于报哪一所，而不是一个统一答案。
\nref{S0102}
\nref{S0118}
\nref{S0302}
\nref{S0417}
\nref{S0119}
\nref{S0157}
\nref{S0495}
\end{qa}

\iin{山东大学}
\ith{双非友好}
\ist{院校与项目差异}
\begin{qa}{E0362}{哪些学校有接收双非学生的传统？}
\textbf{“传统”这种说法无法从文件里证实，但可以找出条件写法上确实不看本科出身的院系。}

以下是本手册收录原文、且申请条件中未设本科学校层次或排名分档的院系（均为 2026—2027 年度）：

\par\vspace{2pt}\nobreak\begin{center}\footnotesize
\begin{tabularx}{\textwidth}{>{\raggedright\arraybackslash}X>{\raggedright\arraybackslash}X}
\toprule
\textbf{院系} & \textbf{条件写法} \\ \midrule
武汉大学计算机学院（2027） & 凡取得所在学校推免生资格、身心健康的考生均可接收 \\
西安电子科技大学（2027 校级与计算机学院） & 仅要求获得本科就读高校推荐资格、无学术不端 \\
湖南大学计算机学院（2027） & 2027 届相关普通全日制本科毕业生，有望取得推免资格 \\
苏州大学（2027） & 取得所在高校推荐资格、诚实守信、学风端正 \\
郑州大学（2027） & 获得所在学校当年推荐免试资格的优秀本科毕业生 \\
深圳大学医学部（2027 级） & 公告未设本科出身条款 \\
\bottomrule
\end{tabularx}
\end{center}\vspace{2pt}\par

与之相对，同一年度里明确分档的院系也存在（如山东大学控制学院 2026 年要求“其他有推免资格高校”排名前 8\%）。

需要注意的是：\textbf{“不设出身门槛”与“录取概率高”是两件事}。前者是文件事实，后者取决于报考人数和导师意愿，本手册未收录报录比数据，因此不做判断。
\nref{S0119}
\nref{S0157}
\nref{S0120}
\nref{S0130}
\nref{S0129}
\nref{S0599}
\nref{S0635}
\nref{S0265}
\nref{S0102}
\end{qa}

\iin{武汉大学}
\ith{生源偏好}
\ist{院校与项目差异}
\begin{qa}{E0366}{哪些学校保研要求过六级或者雅思托福？}
\textbf{分布很宽：从“四级 425”到“六级 480”，还有院系干脆不给分数线。}

\par\vspace{2pt}\nobreak\begin{center}\footnotesize
\begin{tabularx}{\textwidth}{>{\raggedright\arraybackslash}X>{\raggedright\arraybackslash}X}
\toprule
\textbf{学校/院系（年份）} & \textbf{英语门槛写法} \\ \midrule
内蒙古大学 & 四级 425 分以上（艺术、体育类 355 分） \\
安徽师范大学（2027） & 四级不低于 425 分（艺术、体育类 355 分） \\
北大环境与能源学院（2027） & 四级 500 分以上，或六级 425 分以上 \\
天津大学计算机学院（2026） & 六级≥425，\textbf{或四级≥550} \\
南开大学法学院（2027） & 六级 425 分及以上，\textbf{或四级 550 分及以上} \\
北大计算机学院（2026） & 四级 525 分以上，或六级 450 分以上 \\
浙江大学网络空间安全学院（2026） & 六级 460 分及以上 \\
浙江大学医学院（2027 直博） & 六级 480 分及以上（可放宽至 460） \\
\bottomrule
\end{tabularx}
\end{center}\vspace{2pt}\par

\textbf{三点规律}：一是“四级 550 替代六级 425”的写法大量存在；二是同一所学校不同院系档次不同（北大计算机 525/450，北大环境与能源 500/425）；三是部分院系不设分数线——北京大学软件与微电子学院 2026 年夏令营只写“英语水平优良（需提供相关证明）”。

另需注意，浙江大学多个院系明确\textbf{不认可家考成绩与拼分成绩}，北京大学多个学院写明托福“只接受当次考核成绩，不接受拼分成绩”。
\nref{S0544}
\nref{S0567}
\nref{S0646}
\nref{S0310}
\nref{S0114}
\nref{S0118}
\nref{S0239}
\nref{S0216}
\nref{S0017}
\nref{S0029}
\nref{S0341}
\nref{S0213}
\end{qa}

\iin{内蒙古大学}
\ith{语言门槛}
\ist{院校与项目差异}
\begin{qa}{E0367}{哪些学校九推复试是纯面试没有笔试？}
\textbf{确实有一批院系是纯面试，规律是：计算机类以机试为主，文科与部分综合类以纯面试为主。}

\par\vspace{2pt}\nobreak\begin{center}\footnotesize
\begin{tabularx}{\textwidth}{>{\raggedright\arraybackslash}X>{\raggedright\arraybackslash}X}
\toprule
\textbf{院系（年份）} & \textbf{考核形式原文} \\ \midrule
湖南大学计算机学院（2027） & 面试由外语能力测试、专业素质和能力测试、综合素质测试组成；\textbf{综合考核总成绩＝面试成绩} \\
华中师范大学计算机学院（2026） & 线下面试，综合素质面试 80\%＋英语能力面试 20\% \\
湖北工业大学马克思主义学院（2027） & 考核方式为综合面试，总分 100 分 \\
华北电力大学（2027） & 硕士复试总分 150＝综合面试 100＋外语听力及口语 50 \\
苏州大学（2027） & 面试为复试考核的必要环节，综合考核百分制、60 分合格 \\
\bottomrule
\end{tabularx}
\end{center}\vspace{2pt}\par

\textbf{两点提醒}：一是校级办法常写“具体考核形式以学院通知为准”（华北电力大学 2027 年即如此），只有学院细则才能最终确定有没有机试；二是反例就在同一年度——武汉大学计算机学院 2027 年上机测试占 40\%，北京师范大学人工智能学院 2027 年考机试（程序设计、数据结构与算法）。
\nref{S0130}
\nref{S0105}
\nref{S0014}
\nref{S0305}
\nref{S0415}
\nref{S0599}
\nref{S0119}
\nref{S0072}
\end{qa}

\iin{湖南大学}
\ith{考核形式}
\ist{院校与项目差异}
\begin{qa}{E0386}{哪些学校是出了名的“鸽王”？}
\textbf{未找到任何能支撑“鸽王”名单的证据，本手册不列这类名单。}

“鸽王”通常指学校在推免中给出意向或口头承诺后不兑现。要证明这一点，需要的是录取名单变动、待录取通知撤回、补录通知之类的记录，而本手册所依据的公开文件里没有这类材料。

恰恰相反，多家高校在公开文件中写的是\textbf{反方向的承诺}：

\begin{itemize}
  \item 西安电子科技大学 2027 年推免接收实施办法与计算机学院预报名通知：“考生同意待录取后，我校一律不再取消”；
  \item 电子科技大学 2027 年接收通知：“申请人同意接受我校待录取通知后，学校均不再取消待录取”；
  \item 湖南大学计算机学院 2027 年实施细则：明确推免服务系统各环节的时间节点，并说明接受或拒绝操作“一旦操作不可更改”。
\end{itemize}

这些条款约束的其实是\textbf{双方}：你接受了待录取就不能反悔，学校也不能单方面撤销。所以真正可控的做法是重视待录取通知的确认时限——多所院校写明了逾时未确认即视为自动放弃。

本手册不作院校声誉类判断。
\nref{S0157}
\nref{S0120}
\nref{S0495}
\nref{S0130}
\end{qa}

\iin{西安电子科技大学}
\ith{毁约名声}
\ist{院校与项目差异}
\begin{qa}{E0387}{哪些学校对推免生特别友好、给钱多？}
\textbf{“给钱多”可以在公开的奖助条款里查到一部分，本手册列出可核对的原文，不做“友好度”排序。}

\textbf{电子科技大学 2027 年接收通知}写得很具体：学业奖学金——

\par\vspace{2pt}\nobreak\begin{center}\footnotesize
\begin{tabularx}{\textwidth}{>{\raggedright\arraybackslash}X>{\raggedright\arraybackslash}X}
\toprule
\textbf{类别} & \textbf{学业奖学金} \\ \midrule
本科院校为第二轮“双一流”建设高校的推免硕士生（非定向） & 可获评硕士一年级学业奖学金一等奖 1 万元/生，或二等奖 0.8 万元/生 \\
直博生（非定向） & 可获评博士一年级学业奖学金 1 万元/生 \\
\bottomrule
\end{tabularx}
\end{center}\vspace{2pt}\par

注意这条的限定条件是“本科院校为第二轮'双一流'建设高校”——本科出身会直接影响到手金额。

\textbf{西南交通大学 2026 年接收实施办法}给出的是非现金类倾斜：被录取的推免生可优先自主选择导师、优先进入导师课题组和实验室、优先推荐参与国际交流项目并享受相关资助、优先推荐硕博连读，优秀本科生可直接攻读博士研究生并享受博士生奖助计划。

\textbf{武汉大学计算机学院 2027 年实施细则}提到推免生奖励措施，包括研究生优秀新生奖学金、国家奖学金、学业奖学金、国家助学金以及“三助”工作岗位与研究生专项奖学金，具体依据学校各项评定文件执行。

国家层面的保底是研究生国家助学金（财政部、教育部《研究生国家助学金管理暂行办法》）。
\nref{S0495}
\nref{S0699}
\nref{S0119}
\nref{S0203}
\end{qa}

\iin{电子科技大学}
\ith{待遇口碑}
\ist{院校与项目差异}
\begin{qa}{S0266}{计算机类保研和其他工科相比，最特殊的规则差异是什么？}
\textbf{最特殊的不是门槛更高，而是考核里多了其他工科几乎没有的环节——机试／上机测试，且往往占大头。}

\par\vspace{2pt}\nobreak\begin{center}\footnotesize
\begin{tabularx}{\textwidth}{>{\raggedright\arraybackslash}X>{\raggedright\arraybackslash}X>{\raggedright\arraybackslash}X}
\toprule
\textbf{院系（年份）} & \textbf{计算机化考核环节} & \textbf{权重} \\ \midrule
武汉大学计算机学院（2027） & 上机测试（程序实现能力与算法思维） & 上机 40\% \\
贵州大学计算机科学与技术学院·贵州保密学院（2027） & 上机考试（限 C、C++、Java） & 上机 100 分／总分 200 分 \\
北京师范大学人工智能学院（2027） & 上机测试（程序设计、数据结构与算法） & 综合考核＝机试＋面试 \\
北航计算机学院（2027 接收方案） & 附件含“C 语言考试成绩折算协议书” & 折算制 \\
\bottomrule
\end{tabularx}
\end{center}\vspace{2pt}\par

同校对比更说明问题：湖南大学计算机学院 2027 年综合考核总成绩\textbf{就等于面试成绩}，而武汉大学计算机学院同年上机占 40\%——同为计算机类院系，是否上机、上机占多少，差别极大。

另一层是竞赛成绩被量化进评价：湖南大学计算机学院 2027 届推免名单公示中，CCF CSP 认证 350 分加 0.02、ICPC 国际大学生程序设计竞赛三等奖加 0.05、蓝桥杯省二等奖加 0.02，均公开可查。数量化指标是计算机类比其他工科多出来的一层选拔机制。
\nref{S0119}
\nref{S0002}
\nref{S0072}
\nref{S0580}
\nref{S0130}
\nref{S0131}
\end{qa}

\iin{武汉大学}
\ith{学科差异-计算机}
\ist{院校与项目差异}
\begin{qa}{S0267}{计算机保研中机试成绩的权重有多大，各校差别如何？}
\textbf{各校差别极大，从 0\% 到 50\% 都有实例，不存在统一权重。}

\par\vspace{2pt}\nobreak\begin{center}\footnotesize
\begin{tabularx}{\textwidth}{>{\raggedright\arraybackslash}X>{\raggedright\arraybackslash}X}
\toprule
\textbf{院系（年份）} & \textbf{机试权重} \\ \midrule
湖南大学计算机学院（2027） & 无上机环节，\textbf{综合考核总成绩＝面试成绩} \\
华中师范大学计算机学院（2026） & 无上机环节，综合素质面试 80\%＋英语能力面试 20\% \\
北京师范大学人工智能学院（2027） & 综合考核成绩＝机试成绩＋面试成绩 \\
武汉大学计算机学院（2027） & 素质评估 20\%＋\textbf{上机测试 40\%}＋综合面试 40\% \\
贵州大学计算机科学与技术学院（2027） & 上机 100 分／复试总分 200 分，即 \textbf{50\%} \\
\bottomrule
\end{tabularx}
\end{center}\vspace{2pt}\par

三个规律：一是上机测试通常单独计分且占比等于或超过面试；二是面试普遍设 60 分及格线（湖南大学计算机学院明确“面试成绩低于 60 分不予录取”）；三是机试内容集中在编程与算法——贵州大学限用 C、C++、Java 并限定开发环境（VC++6.0、codeblock、dev C++、C-free、Eclipse），北京师范大学明确考“程序设计、数据结构与算法”。

所以报计算机类院系前，必须先确认该院当年细则里有没有上机测试这一项。
\nref{S0130}
\nref{S0105}
\nref{S0072}
\nref{S0002}
\nref{S0119}
\end{qa}

\iin{湖南大学}
\ith{学科差异-计算机}
\ist{院校与项目差异}
\begin{qa}{S0268}{计算机保研为什么很多顶尖院校只招直博，想去读硕怎么办？}
\textbf{证据显示是结构性倾斜，而非一刀切；但确有院校把学硕并入了硕博贯通。}

\par\vspace{2pt}\nobreak\begin{center}\footnotesize
\begin{tabularx}{\textwidth}{>{\raggedright\arraybackslash}X>{\raggedright\arraybackslash}X}
\toprule
\textbf{类型} & \textbf{实例（年份）} \\ \midrule
直博口径宽 & 中山大学 2027 年：所有学术学位专业及除临床医学、口腔医学外的专业学位均招收直博生，直博学制 5 年 \\
直博名额明确单列 & 武汉大学计算机学院 2027 年：拟招硕士推免生约 101 人（学术 53、专业 48），另选拔招收直博生约 50 人 \\
学硕被并入硕博贯通 & 上海财经大学金融学院 2027 年：选拔针对直接攻博、学术学位硕士（\textbf{仅招硕博连读}）和硕士专业学位 \\
直博设比例上限 & 国防科技大学 2027 年：所有招博学科均可接收直博生，一般不超过当年博士招生计划的 30\%，“重大工程联合培养”专项不受限制 \\
\bottomrule
\end{tabularx}
\end{center}\vspace{2pt}\par

想去读硕，两条路：一是报专业学位（武汉大学计算机学院的计算机技术 085404、软件工程 085405、电子信息 085400 为专业学位，接收推免硕士）；二是选“硕士＋可硕博连读”的院系过渡，如武汉大学计算机学院学术学位硕士同时接收推免硕士与直博两类。

本手册未收录“顶尖院校只招直博”的普遍性规定，该说法无法证伪，也不做断言。
\nref{S0167}
\nref{S0650}
\nref{S0330}
\nref{S0119}
\end{qa}

\iin{中山大学}
\ith{学科差异-计算机}
\ist{院校与项目差异}
\begin{qa}{S0269}{计算机保研中顶会论文的作用比其他学科大多少？}
\textbf{论文在计算机类的作用是可以被量化的，而其他学科的推免文件里几乎不出现这类量化条款。}

最直接的实例来自\textbf{中国人民大学 2027 年工程硕博士培养改革专项推免研究生招生报名通知}：申请材料中可提交体现自身学术水平的代表性学术论文、获奖证书、专利证明、\textbf{CSP 成绩}等；其中写明“申请高瓴人工智能学院的学生，如进入复试考核，且 \textbf{CCF CSP 成绩在 300 分及以上，笔试成绩可增加 5 分，总分不超过 75 分}”。

这是计算机类独有的写法——把一个具体的程序设计能力考试成绩直接折算成分数，而非常见的“同等条件下优先”。

另一类证据是\textbf{校内推免遴选中的竞赛与认证加分}：湖南大学计算机学院 2027 届推免名单公示里，逐人列出加分依据与分值，如“第 41 次 CCF CSP 认证 350 分 ＋0.02”“2026 年 ICPC 国际大学生程序设计竞赛三等奖 ＋0.05”“第十六届蓝桥杯全国软件和信息技术专业人才大赛省二等奖 ＋0.02”。这些加分精确到小数点后两位并公开公示。

需要说明的是：本手册未收录“顶会论文在其他学科作用如何”的可比数据，因此不能断言计算机类的作用“比其他学科大多少”，只能说计算机类存在把论文、竞赛、认证折算为分数的成熟机制，且该机制公开可查。
\nref{S0286}
\nref{S0131}
\nref{S0110}
\end{qa}

\iin{中国人民大学}
\ith{学科差异-计算机}
\ist{院校与项目差异}
\begin{qa}{S0270}{双非计算机学生保研到985，哪些学校对双非相对友好？}
\textbf{结论：先找出身条件里不设分档、且考核以面试为主的计算机院系，这与学校名气关系不大。}

\textbf{两类可直接核对的院系（2027 年）}

\par\vspace{2pt}\nobreak\begin{center}\footnotesize
\begin{tabularx}{\textwidth}{>{\raggedright\arraybackslash}X>{\raggedright\arraybackslash}X>{\raggedright\arraybackslash}X}
\toprule
\textbf{院系} & \textbf{申请条件怎么写} & \textbf{考核形式} \\ \midrule
武汉大学计算机学院 & 凡取得所在学校推免生资格、身心健康的考生均可接收；双一流只是“优先考虑” & 上机测试占 40\% \\
西电计算机科学与技术学院 & 接收基本要求 6 条，只要求“获得本科就读高校推荐资格”、无学术不端 & 学院统一综合面试，≥20 分钟 \\
湖南大学计算机学院 & 2027 届相关普通全日制本科毕业生、有望取得推免资格 & \textbf{综合考核总成绩＝面试成绩}，无机试 \\
\bottomrule
\end{tabularx}
\end{center}\vspace{2pt}\par

\textbf{反例同样存在}：山东大学控制科学与工程学院 2026 年夏令营按本科学校分层——985 前 20\%、211 前 10\%、其他有推免资格高校前 8\%。

所以双非计算机同学的机会不在“哪所学校名声好”，而在于能否找到条件原文里不设分档、考核又不押注机试的院系。
\nref{S0119}
\nref{S0120}
\nref{S0130}
\nref{S0129}
\nref{S0102}
\end{qa}

\iin{武汉大学}
\ith{学科差异-计算机}
\ist{院校与项目差异}
\begin{qa}{S0356}{计算机保研中「机试零分直接淘汰」是普遍规则吗？}
\textbf{不是普遍规则。“机试零分直接淘汰”在本手册所依据的文件里找不到出处，但有几条相邻的真实规则容易被误传成这句话。}

\textbf{真实规则一：综合成绩不及格不录取。} 武汉大学计算机学院 2027 年实施细则写明“复试综合成绩＝素质评估×20\%＋上机测试×40\%＋综合面试×40\%，复试综合成绩、素质评估成绩及综合面试成绩均为百分制。\textbf{复试综合成绩不及格者不能录取}”——不及格的是加权后的\textbf{综合成绩}，不是单纯的机试。

\textbf{真实规则二：面试低于 60 分不录取。} 湖南大学计算机学院 2027 年实施细则：“面试成绩低于 60 分不予录取”；贵州大学计算机科学与技术学院（贵州保密学院）2027 年工作方案：“未经复试或复试成绩不及格（低于 60 分）的申请者不录取”。

\textbf{真实规则三：复试成绩按总分排序。} 贵州大学该方案明确“复试总成绩 200 分，其中上机考试满分 100 分、英语面试满分 20 分、综合面试满分 80 分”，按复试总成绩（百分制）排序择优录取。

\textbf{反例}：湖南大学计算机学院 2027 年根本不设机试，总成绩就是面试成绩——在它那里“机试零分”这个前提不存在。

结论：机试重要，但把机试零分等同于直接淘汰，是把“综合成绩不及格”和“面试不及格”两条规则压缩后形成的说法，与原文不符。
\nref{S0119}
\nref{S0130}
\nref{S0002}
\end{qa}

\iin{武汉大学}
\ith{学科差异-计算机}
\ist{院校与项目差异}
\section{科研院所}
\begin{qa}{E0357}{中科院各所的推免和高校有什么不同？}
\textbf{结论：差在学籍归属、第一年上课地点、日常管理与食宿口径这四处，而且同一批中科院研究所内部也不统一。}

\par\vspace{2pt}\nobreak\begin{center}\footnotesize
\begin{tabularx}{\textwidth}{>{\raggedright\arraybackslash}X>{\raggedright\arraybackslash}X}
\toprule
\textbf{单位（年份）} & \textbf{制度上写明的差异} \\ \midrule
中国科大天文与空间科学学院（紫金山天文台，2027） & 属于中国科大学籍；硕士与直博生第一年在中国科大本部完成主要课程，第二年起在紫台（南京）科研实践与学位论文，发中国科大证书 \nref{S0294} \\
中国科学院数学与系统科学研究院（2026） & 京外高校学生线上、京内高校学生线下参加营；营期内组织推免生面试 \nref{S0074} \\
中国科学院心理研究所（2026） & 学员自己解决住宿及往返交通，研究所发放适量住宿补贴 \nref{S0298} \\
国科大杭高院生命与健康科学学院（2026） & 报销从学校来杭的单程路费，营期提供免费餐食、住宿由院统一安排 \nref{S0007} \\
中国科大应用化学与工程学院（长春应化所，2027） & 用问卷星预报名，研究所初审后另行通知面试，未接到通知即未获面试资格 \nref{S0057} \\
\bottomrule
\end{tabularx}
\end{center}\vspace{2pt}\par

对照高校：高校的系统入口是校级研究生招生系统加教育部推免服务系统两级；上面这些研究所常用本所邮箱、问卷星作为第一道报名入口，学院那一层的话事权更大。报考前先读通知抬头的学籍与培养地点条款。
\nref{S0294}
\nref{S0074}
\nref{S0298}
\nref{S0007}
\nref{S0057}
\end{qa}

\iin{中国科学院}
\ith{科研院所}
\ist{院校与项目差异}
\begin{qa}{E0359}{985 分校、异地校区和本部含金量差多少？}
\textbf{结论：制度文件里没有含金量这个说法，只有培养地点与培养单位这个可查字段；同一所学校内部就能按校区分开招生。}

\par\vspace{2pt}\nobreak\begin{center}\footnotesize
\begin{tabularx}{\textwidth}{>{\raggedright\arraybackslash}X>{\raggedright\arraybackslash}X}
\toprule
\textbf{单位（年份）} & \textbf{校区相关原文口径} \\ \midrule
北京师范大学文理学院（2026） & 学院是北京师范大学在珠海校区设立的建制性教学科研机构，下设中文、历史、哲学、外语、数学、物理等13个文理基础学科专业系 \nref{S0149} \\
北京师范大学数学科学学院（2026） & 学术型硕士和直博生营员共200人左右；专业型硕士（珠海校区培养）营员30人左右，同一学院按校区分开招 \nref{S0257} \\
中山大学生物医学工程学院（2026） & 夏令营在中山大学深圳校区举办 \nref{S0078} \\
哈尔滨工业大学（2027） & 全日制硕士研究生可申请住宿，校本部原则上2-4人间；威海校区、深圳校区结合校区实际情况，提供住宿保障 \nref{S0067} \\
东南大学计算机学院（2025） & 苏州联合研究生院计算机技术/人工智能专业的学生由我院负责培养，但报名要关注苏州联合研究生院发布的通知 \nref{S0126} \\
\bottomrule
\end{tabularx}
\end{center}\vspace{2pt}\par

判断办法是看招生专业目录的备注栏与培养单位名称，而不是看学校牌子。
\nref{S0149}
\nref{S0257}
\nref{S0078}
\nref{S0067}
\nref{S0126}
\end{qa}

\iin{北京师范大学}
\ith{校区差异}
\ist{院校与项目差异}
\begin{qa}{E0361}{哪些学校明确不收双非推免生？怎么提前避开？}
\textbf{结论：没有一所学校会发文写不收双非，但不少院系把生源院校写进了申请条件——这是可查的量化条款，不是传闻。}

\par\vspace{2pt}\nobreak\begin{center}\footnotesize
\begin{tabularx}{\textwidth}{>{\raggedright\arraybackslash}X>{\raggedright\arraybackslash}X}
\toprule
\textbf{单位（年份）} & \textbf{写进申请条件的生源范围} \\ \midrule
华中科技大学网络空间安全学院（2026） & 申请人须所在学校为国家双一流大学，或所在学科为国家双一流学科或在最新一轮学科评估中排名靠前；其他高校须学习成绩特别优异或有突出成果或有特殊专业特长 \nref{S0003} \\
西安交通大学网络空间安全学院（2026） & ①双一流建设高校且专业前30\%；②所在学校专业为双一流建设学科，或最新一轮学科评估A-及以上，且专业前15\%；③经学院认定在科研、国家级竞赛有突出成果 \nref{S0132} \\
北京交通大学计算机学院（2026） & 分四档：本校应届前40\%；同类高校应届前40\%；具有相关优势学科的高校应届前10\%；或竞赛全国二等奖以上、第一发明人专利、第一作者A类论文 \nref{S0104} \\
中南大学直博（2027） & 限国内高水平大学优秀应届本科毕业生；本科留学生须本科就读高校同时入围QS、THE、USNews、软科前100名 \nref{S0625} \\
\bottomrule
\end{tabularx}
\end{center}\vspace{2pt}\par

这些条款一律写在申请条件第一段，不需要打听。
\nref{S0003}
\nref{S0132}
\nref{S0104}
\nref{S0625}
\end{qa}

\iin{华中科技大学}
\ith{隐性门槛}
\ist{院校与项目差异}
\begin{qa}{E0372}{招生简章里“拟招生人数含推免”该怎么理解？}
\textbf{结论：拟招生人数含推免，是说这个专业的总规模已经把推免生算进去了，推免没占满的名额才回落到统考；而目录里的数字本身是可调的预估数。}

\par\vspace{2pt}\nobreak\begin{center}\footnotesize
\begin{tabularx}{\textwidth}{>{\raggedright\arraybackslash}X>{\raggedright\arraybackslash}X}
\toprule
\textbf{单位（年份）} & \textbf{原文口径} \\ \midrule
四川大学（2027章程） & 目录招生人数为预计招生人数，实际录取人数将视申请及考核情况进行调整；2027年预计招收推免硕士3150名、直博生810名 \nref{S0476} \\
天津大学化工学院（2027） & 专业目录中拟招生人数均为根据以往生源情况所做的预估数，实际推免招生人数将根据录取情况，报经学校研究生招生委员会同意后进行调整 \nref{S0303} \\
中国社会科学院大学法学院（2027） & 各专业实际接收人数将根据生源情况进行调整，如未完成推免接收计划，剩余名额将用于统考招生 \nref{S0243} \\
西安交通大学（2027章程） & 招生目录在研招网公布，含各学院拟招生学科、专业及推免计划 \nref{S0548} \\
\bottomrule
\end{tabularx}
\end{center}\vspace{2pt}\par

对你真正有用的数字，是该院拟接收推免生人数那一栏，而不是统考目录里的总数。
\nref{S0476}
\nref{S0303}
\nref{S0243}
\nref{S0548}
\end{qa}

\iin{四川大学}
\ith{招生计划}
\ist{院校与项目差异}
\begin{qa}{E0374}{导师制的学校和企业导师制的学校有什么区别？}
\textbf{结论：制度上区分的不是哪种导师制，而是学术学位走硕博贯通、专业学位走校企协同；企业是否全程参与是分界线。}

\par\vspace{2pt}\nobreak\begin{center}\footnotesize
\begin{tabularx}{\textwidth}{>{\raggedright\arraybackslash}X>{\raggedright\arraybackslash}X}
\toprule
\textbf{单位（年份）} & \textbf{培养模式原文} \\ \midrule
西安交通大学（2027章程） & 学术学位全日制研究生一般按硕博贯通模式培养；专业学位全日制研究生一般按校企协同模式、硕博贯通模式（包括百千万卓越工程人才培养项目、协同培养育人项目）培养，行业企业全程参与到人才培养的各个环节，并在苏州、青岛、深圳等地建立联合培养、专业实践基地 \nref{S0548} \\
西安交通大学计算机学院（2027细则） & 同一段表述被学院细则逐字沿用，并补相关信息可向报考学院咨询 \nref{S0662} \\
哈尔滨工业大学（2027） & 全日制硕士可申请住宿，校本部原则上2-4人间；校企联培学生住宿依据培养模式另行安排 \nref{S0067} \\
四川大学（2027） & 另发2027年四川大学-南方电网联合培养项目推荐免试研究生招生简章，专项单独成文 \nref{S0478} \\
\bottomrule
\end{tabularx}
\end{center}\vspace{2pt}\par

同一所学校里，读学硕还是专硕，直接决定你在实验室还是企业待得多。
\nref{S0548}
\nref{S0662}
\nref{S0067}
\nref{S0478}
\end{qa}

\iin{西安交通大学}
\ith{培养模式}
\ist{院校与项目差异}
\begin{qa}{E0377}{军校、警校的推免怎么走？}
\textbf{结论：军校推免与地方高校最大的差别是多了参军入伍这一条并行通道，录取按等次分批而不是单纯按分数排序；所有类型最终仍要走教育部推免系统。}

以国防科技大学《接收2027级地方院校推荐免试硕士研究生（含直博生、参军入伍生）工作方案》（2026年6月18日发布）为例 \nref{S0650}：

\par\vspace{2pt}\nobreak\begin{center}\footnotesize
\begin{tabularx}{\textwidth}{>{\raggedright\arraybackslash}X>{\raggedright\arraybackslash}X}
\toprule
\textbf{环节} & \textbf{原文口径} \\ \midrule
预报名 & 先在国防科技大学研究生招生信息网推荐免试平台报名，7至9月；必须先通过我校报名系统预报名后方可获得复试资格，如未进行预报名将被视为无效申请 \\
复试与体检 & 7至9月参加复试、体检；提交学生证、学籍认证报告、申请表、自述（1000字以内）、成绩单等（2026级方案同类结构见 \nref{S0420}） \\
结果等次 & 根据考生复试成绩从高到低依次按优秀、合格、不合格三类等次确定复试结果，不合格等次不得录取 \\
分批发待录取 & 为优秀等次生源发放待录取通知、录取不分先后；为合格生源按复试成绩从高到低依次发送待录取通知，直至计划录满为止 \\
参军入伍生 & 须符合军队直接选拔招录军官条件，本科院校为双一流建设高校，本科专业为双一流建设学科或A类学科的对应专业；入学不参军的是地方研究生 \\
\bottomrule
\end{tabularx}
\end{center}\vspace{2pt}\par

体检不合格但符合地方硕士录取条件者可申请转录。警校方向本轮未找到一手推免文件。
\nref{S0650}
\nref{S0420}
\end{qa}

\iin{国防科技大学}
\ith{军警院校}
\ist{院校与项目差异}
\begin{qa}{E0380}{强基计划高校的转段比例大概是多少？}
\textbf{判断：未找到任何单位公布统一的强基转段比例数值，也不能用某个比例去套另一所学校。转段是考核制，不是按比例直升。}

\par\vspace{2pt}\nobreak\begin{center}\footnotesize
\begin{tabularx}{\textwidth}{>{\raggedright\arraybackslash}X>{\raggedright\arraybackslash}X}
\toprule
\textbf{单位或文件（年份）} & \textbf{原文口径} \\ \midrule
教育部教学〔2020〕1号 & 只规定强基计划的招生办法，全篇未设转段比例条款 \nref{S0045} \\
华中科技大学（2026） & 已获得本研转段资格的强基生只能申请校内相关院系的硕士或直接攻读博士，每人限选1个专业 \nref{S0536} \\
中山大学生命科学学院（2026） & 未参加或未通过转段考核的强基计划本科毕业生，视为退出强基计划，不再具有申请免试攻读研究生资格 \nref{S0248} \\
中南大学（2027直博通知） & 拟安排直博生指导性招生计划400名，其中包含强基计划转段直博；二级招生单位各专业直博生一般不超过本专业招生人数的25\% \nref{S0625} \\
华东师范大学（2027预报名通知） & 强基计划转段、公费师范生本研衔接转段等项目不通过本系统报名 \nref{S0610} \\
\bottomrule
\end{tabularx}
\end{center}\vspace{2pt}\par

要看比例，只能查该校当年强基计划学生本研衔接转段工作通知里的专业范围与考核办法。
\nref{S0045}
\nref{S0536}
\nref{S0248}
\nref{S0625}
\nref{S0610}
\end{qa}

\iin{中华人民共和国教育部}
\ith{强基比例}
\ist{院校与项目差异}
\begin{qa}{S0305}{军校保研（如国防科技大学）和地方高校推免流程有什么不同？}
\textbf{结论：军校先给考生分优秀、合格、不合格等次，再按等次分批录取；地方高校多为差额复试后按成绩排序。军校还多一道军队体检与政审环节。}

\par\vspace{2pt}\nobreak\begin{center}\footnotesize
\begin{tabularx}{\textwidth}{>{\raggedright\arraybackslash}X>{\raggedright\arraybackslash}X>{\raggedright\arraybackslash}X}
\toprule
\textbf{对比项} & \textbf{国防科技大学（2027级方案）} & \textbf{地方高校（以北京大学2027办法为例）} \\ \midrule
预报名 & 必须在国防科大推荐免试平台预报名，未进行预报名将被视为无效申请 \nref{S0650} & 北大研招网申报，2026年9月1日10:00至9月8日10:00；外校生与本校跨院系生须交全部材料 \nref{S0682} \\
复试结果 & 按优秀、合格、不合格三类等次确定，不合格不得录取 \nref{S0650} & 实行差额复试、择优录取，差额比例与复试形式由各院系公布 \nref{S0682} \\
录取顺序 & 对优秀等次生源发放待录取通知、不分先后；对合格生源按成绩从高到低发至录满 \nref{S0650} & 按院系复试择优确定待录取名单，未取得推免资格者待录取无效 \nref{S0682} \\
共同点 & 最终都须在教育部推免服务系统填志愿并确认待录取，否则无法完成录取 & 同上 \nref{S0682} \\
\bottomrule
\end{tabularx}
\end{center}\vspace{2pt}\par

军校另设参军入伍生类别，须符合军队直接选拔招录军官条件。
\nref{S0650}
\nref{S0682}
\end{qa}

\iin{国防科技大学}
\ith{院校类型-军校}
\ist{学科与院校差异}
\begin{qa}{S0315}{社科院系统的推免有什么特点？}
\textbf{结论：社科院系统的特点是学院是招生主体、体量小、学硕与直博并重、复试统一为高权重面试，而且同一所大学内部各学院条件可以差很多。}

\par\vspace{2pt}\nobreak\begin{center}\footnotesize
\begin{tabularx}{\textwidth}{>{\raggedright\arraybackslash}X>{\raggedright\arraybackslash}X>{\raggedright\arraybackslash}X}
\toprule
\textbf{对比项} & \textbf{中国社会科学院大学经济学院（2027）} & \textbf{中国社会科学院大学法学院（2027）} \\ \midrule
接收类型 & 各二级学科均接收推免硕士；直博仅在经济史、西方经济学、发展经济学和数字经济学四个专业接收 \nref{S0308} & 拟接收学术学位推免硕士11人、直博生2人 \nref{S0243} \\
直博外语 & 六级425分（含）以上 \nref{S0308} & 六级525分（含）以上，且学业成绩一般在前20\%以内 \nref{S0243} \\
复试形式 & 网络远程、面试；总分300分（专业素质150加综合素质150，其中外语听说50），180分合格；原则上按不低于120\%的比例差额组织复试考核 \nref{S0308} & 网络远程、面试；总分300分，合格线180分（专业能力150加外语综合50加综合能力100）\nref{S0243} \\
待录取确认 & 通知发出后2小时内确认，逾期视为放弃 \nref{S0308} & 在学院规定时间内确认，逾期视为放弃 \nref{S0243} \\
\bottomrule
\end{tabularx}
\end{center}\vspace{2pt}\par

直博学习年限不低于五年（经济学院2027原文）。
\nref{S0308}
\nref{S0243}
\end{qa}

\iin{中国社会科学院大学经济学院}
\ith{院校类型-社科院}
\ist{学科与院校差异}
\begin{qa}{S0316}{中国农业科学院、中国医学科学院等院所的推免值得考虑吗？}
\textbf{判断：中国农业科学院、中国医学科学院的一手推免简章本轮未找到，本条不能替你做结论。可用的同类制度证据如下。}

中国医学科学院与北京协和医学院2026年夏令营公告列出18个培养单位（含医药生物技术研究所、病原生物学研究所、药用植物研究所、血液病医院、阜外医院、肿瘤医院、医学信息研究所等），并写明具体选拔活动由各培养单位视情况决定，以各培养单位发布的夏令营通知为准 \nref{S0162}。

以病原生物学研究所2026年夏令营为例 \nref{S0375}：招营员60人左右，活动为期一天，地点在北京大兴；申请条件为英语水平须达到大学英语四级425分以上或相当水平，同等条件下优先考虑达到六级430分以上者，且同等条件下优先考虑获得推荐免试研究生资格者；另需两封副教授及以上职称专家推荐信。

可见这类院所是院所合一、多培养单位的结构，规则下沉到每个所和每所医院，必须逐个看通知。
\nref{S0162}
\nref{S0375}
\end{qa}

\iin{北京协和医学院}
\ith{院校类型-科研院所}
\ist{学科与院校差异}
\begin{qa}{S0333}{中国科学技术大学的推免（含科教融合所）有什么特点？}
\textbf{结论：中科大的科教融合所是以中国科大某某学院（某所）的名义招生、学籍归中国科大、培养在研究所，报名入口分散在中科大报考服务系统和研究所自有渠道两处。}

\par\vspace{2pt}\nobreak\begin{center}\footnotesize
\begin{tabularx}{\textwidth}{>{\raggedright\arraybackslash}X>{\raggedright\arraybackslash}X}
\toprule
\textbf{单元（年份）} & \textbf{原文口径} \\ \midrule
中国科大天文与空间科学学院（紫金山天文台，2027） & 学院由中科大与紫台联合共建；录取研究生属于中国科大学籍；硕士和直博生第一年在中科大本部完成主要课程，第二年起在紫台（南京）科研实践和学位论文，毕业颁发中国科大毕业证书和学位证书；预报名走中科大研究生报考服务系统，学院代码170 \nref{S0294} \\
中国科大应用化学与工程学院（长春应化所，2027） & 以中国科大应用化学与工程学院（长春应化所）名义接收硕士与直博生；用问卷星预报名，研究所组织初审后通知面试，未接到通知即未获得面试资格 \nref{S0057} \\
中国科大与宁波东方理工大学（2027） & 单独发布招收免试直博生预推免办法，联合培养项目单列 \nref{S0362} \\
中国科大研究生招生在线（2027） & 校级设有2027年推免生接收办法、2027年工程硕博士培养改革专项接收推免研究生等分项文件 \nref{S0496} \\
\bottomrule
\end{tabularx}
\end{center}\vspace{2pt}\par

判断方法：读通知抬头与学籍、培养地点条款。
\nref{S0294}
\nref{S0057}
\nref{S0362}
\nref{S0496}
\end{qa}

\iin{中国科学院}
\ith{院校特定-中科大}
\ist{学科与院校差异}
\section{院校差异}
\begin{qa}{E0378}{部属师范大学和省属师范的推免政策差在哪？}
\textbf{结论：部属师范大学多出一条本研衔接公费教育通道，把推免资格嵌进公费师范履约里；省属师范按普通推免办法走，细则自定。}

\par\vspace{2pt}\nobreak\begin{center}\footnotesize
\begin{tabularx}{\textwidth}{>{\raggedright\arraybackslash}X>{\raggedright\arraybackslash}X>{\raggedright\arraybackslash}X}
\toprule
\textbf{对比项} & \textbf{部属师范大学} & \textbf{省属师范大学} \\ \midrule
政策依据 & 国务院办公厅转发《教育部直属师范大学本研衔接师范生公费教育实施办法》（国办发〔2024〕27号），教育部教师厅函〔2025〕3号配套落实就业 \nref{S0456}\nref{S0044} & 各校自己的推免工作方案，如安徽师范大学2027届工作方案 \nref{S0646} \\
通道形式 & 部属师范大学本科层次公费师范生在本科三年级结束后，可根据个人意愿提出接受本研衔接公费教育申请，经学校对其进行推免综合考核，取得学校推免资格并达到学校硕士研究生录取条件 \nref{S0044} & 东北师范大学2027另设公费师范毕业生申请免试攻读非全日制教育硕士专业学位研究生办法 \nref{S0609} \\
条件口径 & 由实施办法与部属校办法规定 & 云南师范大学2027：文科、理工类学生前三年专业课平均成绩在专业内排名前20\%；艺术类、体育类不设专业课排名限制，由学院自主制定标准 \nref{S0568} \\
\bottomrule
\end{tabularx}
\end{center}\vspace{2pt}\par

实操上先把公费师范生身份与普通推免生身份分清。
\nref{S0456}
\nref{S0044}
\nref{S0646}
\nref{S0609}
\nref{S0568}
\end{qa}

\iin{教育部学生服务与素质发展中心（学信网/研招网）}
\ith{师范差异}
\ist{院校与项目差异}
\begin{qa}{E0389}{同一个学校不同学院的推免难度差多少？}
\textbf{结论：同一所学校不同学院的推免难度差异，可以直接量化为同样的排名在不同院能不能过初审。}

\par\vspace{2pt}\nobreak\begin{center}\footnotesize
\begin{tabularx}{\textwidth}{>{\raggedright\arraybackslash}X>{\raggedright\arraybackslash}X>{\raggedright\arraybackslash}X}
\toprule
\textbf{对比} & \textbf{单位（年份）} & \textbf{门槛原文} \\ \midrule
同校不同院，外语门槛差100分 & 中国社会科学院大学经济学院（2027）直博 & 六级425分（含）以上 \nref{S0308} \\
同上 & 中国社会科学院大学法学院（2027）直博 & 六级525分（含）以上，且学业成绩一般在前20\%以内 \nref{S0243} \\
同校不同院，批次时间不同 & 天津大学计算机学院（2027） & 直博报名9月2日至15日，考核计划9月18日前后 \nref{S0688} \\
同上 & 天津大学建筑工程学院（2027） & 推免硕士最后一批报名截止9月11日，预计审核9月15日之前 \nref{S0701} \\
同校不同院，名额结构不同 & 中国社会科学院大学法学院（2027） & 拟接收学硕11人、直博2人 \nref{S0243} \\
\bottomrule
\end{tabularx}
\end{center}\vspace{2pt}\par

校级文件只给框架，例如华东师范大学2027年把各院系综合考核实施细则做成汇总表逐院挂链接 \nref{S0613}；条件与名额都在学院那一层。
\nref{S0308}
\nref{S0243}
\nref{S0688}
\nref{S0701}
\nref{S0613}
\end{qa}

\iin{中国社会科学院大学经济学院}
\ith{院内差异}
\ist{院校与项目差异}
\begin{qa}{S0308}{985高校内部不同学院之间的保研难度差异有多大？}
\textbf{结论：985高校内部学院间的差异主要体现在考核组织形式上——有的学院单独出细则、单独定考核内容，有的把考核挂在学院自己组织的营或面试里。}

\par\vspace{2pt}\nobreak\begin{center}\footnotesize
\begin{tabularx}{\textwidth}{>{\raggedright\arraybackslash}X>{\raggedright\arraybackslash}X>{\raggedright\arraybackslash}X}
\toprule
\textbf{角度} & \textbf{单位（年份）} & \textbf{原文口径} \\ \midrule
逐院出细则 & 华东师范大学（2027） & 发布各院系接收优秀应届本科毕业生免试攻读研究生综合考核实施细则汇总（持续更新），按院系代码逐条挂链接 \nref{S0613} \\
逐院出细则 & 西安交通大学（2027） & 连续发布多份各学院（部、中心）招收推免生实施细则，同一名称、不同学院各一份 \nref{S0545}\nref{S0546}\nref{S0547} \\
逐院出细则 & 中国农业大学（2027） & 各院系2027年接收推荐免试研究生工作实施细则（陆续更新中）\nref{S0615} \\
学院自定考核内容 & 北京大学（2027校级办法） & 实行差额复试、择优录取；复试形式包括综合面试、综合笔试、综合实践等，具体差额比例、复试形式和复试选拔办法由各院系根据其学科特点和生源状况确定后公布 \nref{S0682} \\
学院自定活动与测评 & 北京大学计算机学院（2026） & 夏令营内容含专业知识与技能测评，并明示已退出强基计划者请勿报名 \nref{S0114} \\
\bottomrule
\end{tabularx}
\end{center}\vspace{2pt}\par

同一所985里，A院考笔试加机试、B院只有面试，是很正常的事。
\nref{S0613}
\nref{S0545}
\nref{S0546}
\nref{S0547}
\nref{S0615}
\nref{S0682}
\nref{S0114}
\end{qa}

\iin{华东师范大学}
\ith{院校类型-985}
\ist{学科与院校差异}
\begin{qa}{S0330}{北京大学计算机、电子、智能等学院的推免偏好和批次是怎样的？}
\textbf{结论：北大是按校级定规则、学院各自发通知运作的，学院之间没有统一偏好；批次要按各学院自己的通知去看。}

校级底线来自《北京大学2027年接收推荐免试研究生办法（校本部）》（2026年8月31日）\nref{S0682}：

\begin{itemize}
  \item 申报时间：北大研招网2026年9月1日10:00至9月8日10:00，逾期不再接受申请；
  \item 材料：申请表、个人陈述、两封副教授及以上专家推荐信、成绩单、获奖证书、外语水平证明；本校推免生申请本院系可只交其中三项，外校及本校跨院系申请的须交全部；
  \item 复试：现场复试并签署诚信承诺书；实行差额复试，具体差额比例与形式由各院系公布；成绩按百分制、60分及格，任一复试形式不及格者不予录取；面试全程录音录像，每人一般不少于20分钟。
\end{itemize}

学院层面各自发通知：计算机学院2026夏令营 \nref{S0114}、智能学院2026 \nref{S0299}、电子学院2026 \nref{S0143}、软件与微电子学院2026 \nref{S0341}、物理学院2026 \nref{S0384}。偏好没有官方表述。
\nref{S0682}
\nref{S0114}
\nref{S0299}
\nref{S0143}
\nref{S0341}
\nref{S0384}
\end{qa}

\iin{北京大学}
\ith{院校特定-北大}
\ist{学科与院校差异}
\begin{qa}{S0334}{南京大学计算机、人工智能学院的推免偏好如何？}
\textbf{结论：南大同样是校级办法给框架、学院细则定条件，计算机方向的条件写在学院层，不在校级文件里。}

可查到的两级文件：

\begin{itemize}
  \item 校级：《南京大学2027年接收推荐免试研究生工作办法》，附件含2027年接收推荐免试研究生申请材料清单、南京大学考生诚信复试承诺书 \nref{S0505}；
  \item 院级：《南京大学计算机学院2027年接收推荐免试研究生工作实施细则》独立成文 \nref{S0111}；
  \item 对照参考：南京大学商学院2026年细则规定我单位将于9月25日9:00开始在教育部推免服务系统对申请人发出待录取通知，请收到通知的申请人于30分钟内完成确认，未按时确认视为自动放弃并启动补录 \nref{S0282}。
\end{itemize}

关于南大人工智能学院是否有独立细则、两院是否分开招生，本轮所查资料未找到该院独立通知，无法给出判断。偏好属主观判断，制度文件里没有这一项。
\nref{S0505}
\nref{S0111}
\nref{S0282}
\end{qa}

\iin{南京大学}
\ith{院校特定-南大}
\ist{学科与院校差异}
\begin{qa}{S0335}{复旦大学计算机、大数据学院的推免竞争情况如何？}
\textbf{结论：复旦的计算机方向存在明确的多轮选拔——同一年度会发第二轮通知，说明名额是分批投放的；具体竞争程度没有官方数值。}

\par\vspace{2pt}\nobreak\begin{center}\footnotesize
\begin{tabularx}{\textwidth}{>{\raggedright\arraybackslash}X>{\raggedright\arraybackslash}X}
\toprule
\textbf{文件（年份）} & \textbf{原文口径} \\ \midrule
复旦大学2027年招收优秀应届本科毕业生免试攻读研究生章程 & 校级章程，规定当年度推免接收的总规则 \nref{S0190} \\
复旦大学招收2027级推荐免试研究生报名通知 & 校级报名通知 \nref{S0195} \\
复旦大学计算与智能创新学院2027年推免生招生选拔（第二轮）工作通知 & 学院发文，2026年9月27日发布，明确为第二轮选拔 \nref{S0109} \\
复旦大学2027年工程硕博士培养改革专项推免研究生招生报名通知 & 专项计划单独发文 \nref{S0189} \\
\bottomrule
\end{tabularx}
\end{center}\vspace{2pt}\par

可见复旦的批次不是全校一盘棋，而是学院按名额投放节奏自己开轮次。要判断竞争情况，客观口径只有该院往年拟录取公示名单的人数与生源结构，而所查资料中未找到复旦计算机方向当年拟录取公示。
\nref{S0190}
\nref{S0195}
\nref{S0109}
\nref{S0189}
\end{qa}

\iin{复旦大学}
\ith{院校特定-复旦}
\ist{学科与院校差异}
\begin{qa}{S0354}{同一所学校不同学院的推免政策差异有多大？}
\textbf{结论：同一所学校不同学院的推免政策差异，本质是名额分配与细则制定权下放到学院，这在国家文件里就有依据。}

\begin{itemize}
  \item 国家层面：教育部办公厅教学厅〔2014〕5号要求做好信息公开，推免名额、推荐办法、推免生名单（含姓名、院系、综合测评成绩等）、复试录取办法、拟录取推免生名单（含姓名、复试成绩等）、咨询申诉渠道等均须公开 \nref{S0048}\nref{S0205}；
  \item 校级层面：上海交通大学2027年招收优秀应届本科毕业生免试攻读研究生（含直博）办法为全校总办法，各院系在其下自行组织 \nref{S0515}；北京大学2027年办法（校本部）明确具体差额比例、复试形式和复试选拔办法由各院系根据其学科特点和生源状况确定后公布 \nref{S0682}；
  \item 院级层面：以中国社科院大学为例，同一年同一所大学，经济学院直博要求六级425分 \nref{S0308}，法学院直博要求六级525分且成绩一般在前20\%以内 \nref{S0243}。
\end{itemize}

所以我校怎么规定的这句话，必须补上哪个学院。
\nref{S0048}
\nref{S0205}
\nref{S0515}
\nref{S0682}
\nref{S0308}
\nref{S0243}
\end{qa}

\iin{中华人民共和国教育部}
\ith{院校差异}
\ist{学科与院校差异}
\begin{qa}{S0355}{各校推免政策每年的变化大吗，往年经验还能信吗？}
\textbf{结论：逐年会变的是时间节点、专业与名额以及个别口径；不变的是教育部那套公开加系统确认的骨架。往年经验只能用来判断量级，不能替代当年文件。}

同一所学校的两年对比最直观——湖南大学推免工作通知 \nref{S0537}\nref{S0540}：

\par\vspace{2pt}\nobreak\begin{center}\footnotesize
\begin{tabularx}{\textwidth}{>{\raggedright\arraybackslash}X>{\raggedright\arraybackslash}X>{\raggedright\arraybackslash}X}
\toprule
\textbf{事项} & \textbf{2026届} & \textbf{2027届} \\ \midrule
学院学业综合成绩排名公示 & 7月15日前 & 7月21日前 \\
学院推免遴选总成绩排名并公示 & 9月4日前 & 8月31日前 \\
专项环节 & 未列 & 新增国优计划和研究生支教团遴选（9月6日前） \\
\bottomrule
\end{tabularx}
\end{center}\vspace{2pt}\par

制度层面的稳定项：教学厅〔2014〕5号要求推免名额、推荐办法、推免生名单、复试录取办法、拟录取推免生名单、咨询申诉渠道一律公开 \nref{S0048}；有的学校连目录数字都注明可调，如四川大学2027年章程写明目录招生人数为预计招生人数，实际录取人数将视申请及考核情况进行调整 \nref{S0476}。
\nref{S0537}
\nref{S0540}
\nref{S0048}
\nref{S0476}
\end{qa}

\iin{湖南大学}
\ith{院校差异}
\ist{学科与院校差异}
\begin{qa}{S0361}{各校对六级、雅思、托福的门槛差异有多大？}
\textbf{结论：外语门槛跨度很大，从四级合格到六级525分，而且同一所大学内部不同学院就能差100分。必须查目标院系当年通知，不能按学校档次推断。}

\par\vspace{2pt}\nobreak\begin{center}\footnotesize
\begin{tabularx}{\textwidth}{>{\raggedright\arraybackslash}X>{\raggedright\arraybackslash}X}
\toprule
\textbf{单位（年份）} & \textbf{外语门槛原文} \\ \midrule
山西大学哲学学院（2026） & 申请条件仅要求英语四级成绩合格以上 \nref{S0534} \\
中国医学科学院与北京协和医学院病原生物学研究所（2026） & 英语水平须达到大学英语四级425分以上或相当水平，同等条件下优先考虑六级430分以上者 \nref{S0375} \\
华中科技大学网络空间安全学院（2026） & 六级425分及以上；或专业八级；或TOEFL 90分及以上；或IELTS 6分及以上；或GRE 300分及以上；或GMAT 650分及以上 \nref{S0003} \\
中国社会科学院大学经济学院（2027直博） & 参加国家英语六级考试且成绩达到425分（含）以上 \nref{S0308} \\
中国社会科学院大学法学院（2027直博） & 国家英语六级考试且成绩达到525分（含）以上，或具备同等水平的其他种类外语国家级考试成绩 \nref{S0243} \\
中国科学院紫金山天文台（2027） & 只写英语水平良好、具有较强的英语听说读写能力，交四六级、托福、雅思成绩作证明 \nref{S0294} \\
\bottomrule
\end{tabularx}
\end{center}\vspace{2pt}\par

有的院是硬门槛，有的只是同等条件优先，措辞不同性质完全不同。
\nref{S0534}
\nref{S0375}
\nref{S0003}
\nref{S0308}
\nref{S0243}
\nref{S0294}
\end{qa}

\iin{通知公告丨山西大学哲学学院}
\ith{英语门槛}
\ist{学科与院校差异}
\section{院校类型-中科院}
\begin{qa}{E0373}{推免招生目录里的数字准不准？会不会临时改？}
\textbf{结论：会改。招生目录里的数字在招生单位自己的文件里就被定义为预计数，最终以教育部推免服务系统中的招生专业与拟录取公示为准。}

\par\vspace{2pt}\nobreak\begin{center}\footnotesize
\begin{tabularx}{\textwidth}{>{\raggedright\arraybackslash}X>{\raggedright\arraybackslash}X}
\toprule
\textbf{单位（年份）} & \textbf{原文口径} \\ \midrule
四川大学（2027章程） & 目录招生人数为预计招生人数，实际录取人数将视申请及考核情况进行调整 \nref{S0476} \\
天津大学化工学院（2027） & 目录中拟招生人数均为根据以往生源情况所做的预估数，实际推免招生人数将根据录取情况，报经学校研究生招生委员会同意后进行调整 \nref{S0303} \\
中国社会科学院大学法学院（2027） & 各专业实际接收人数将根据生源情况进行调整，如未完成推免接收计划，剩余名额将用于统考招生 \nref{S0243} \\
天津大学（2027） & 教育部推免服务系统中的招生专业和导师为我校最终招生目录，以该系统报名信息为准；未在该系统报名的考生报考信息无效 \nref{S0516} \\
\bottomrule
\end{tabularx}
\end{center}\vspace{2pt}\par

判断口径：目录用于估规模，推免系统里的专业用于确认可报，拟录取公示用于确认最终结果。
\nref{S0476}
\nref{S0303}
\nref{S0243}
\nref{S0516}
\end{qa}

\iin{四川大学}
\ith{目录准确性}
\ist{院校与项目差异}
\begin{qa}{E0381}{学校的本硕连读班和推免有什么区别？}
\textbf{结论：推免是资格加可跨校报考的招生通道；本硕连读、本研贯通是校内培养模式，两者性质不同，而且多数学校明确要求先取得推免资格、且限报本校。}

\par\vspace{2pt}\nobreak\begin{center}\footnotesize
\begin{tabularx}{\textwidth}{>{\raggedright\arraybackslash}X>{\raggedright\arraybackslash}X}
\toprule
\textbf{类型} & \textbf{制度依据与原文} \\ \midrule
推免（全国统一） & 教育部《全国普通高等学校推荐优秀应届本科毕业生免试攻读硕士学位研究生工作管理办法（试行）》：免试指应届本科毕业生不必经过全国统考初试，直接进入复试 \nref{S0454} \\
本硕贯通（校内培养项目） & 湖北工业大学马克思主义学院2027年细则：本校应届本科毕业生获得推免资格、报考本校推免研究生拟录取后，根据个人申请可选择本硕贯通，学习年限3加1加2；另引《湖北工业大学本研贯通人才培养项目实施办法（试行）》\nref{S0014} \\
硕博贯通（培养模式） & 西安交通大学2027年章程：学术学位全日制研究生一般按硕博贯通模式培养 \nref{S0548} \\
强基转段（独立通道） & 教育部教学〔2020〕1号设强基计划本研衔接；华中科技大学2026年要求获转段资格者只能申请校内相关院系的硕士或直接攻读博士 \nref{S0045}\nref{S0536} \\
校内选拔自主权 & 北京大学2027年推荐办法（校本部）：元培学院学生的资格审核和推荐工作由元培学院组织进行 \nref{S0553} \\
\bottomrule
\end{tabularx}
\end{center}\vspace{2pt}\par

一句话：连读班给你的是路径，推免给你的是重新选校的权利。
\nref{S0454}
\nref{S0014}
\nref{S0548}
\nref{S0045}
\nref{S0536}
\nref{S0553}
\end{qa}

\iin{教育部学生服务与素质发展中心（学信网/研招网）}
\ith{长学制}
\ist{院校与项目差异}
\begin{qa}{S0276}{物理类保研中理论物理和凝聚态的推免难度差别在哪？}
\textbf{判断：未找到任何把理论物理与凝聚态物理按难度排序的一手文件。制度上能查到的差别是——两个方向常被拆成不同的营、不同的考核口径。}

\par\vspace{2pt}\nobreak\begin{center}\footnotesize
\begin{tabularx}{\textwidth}{>{\raggedright\arraybackslash}X>{\raggedright\arraybackslash}X}
\toprule
\textbf{单位（年份）} & \textbf{方向组织方式原文} \\ \midrule
中国科学院大学物理科学学院（2026） & 单独举办凝聚态及原子分子物理大学生夏令营，活动含营员选拔面试；营期免费提供食宿，交通费用自理，家庭经济困难的外地同学可申请往返路费补助 \nref{S0292} \\
北京大学物理学院（2026） & 物理学院统一举办优秀大学生暑期夏令营 \nref{S0384} \\
南开大学物理科学学院（2027） & 发布接收推免（含直博）生工作的通知，按学院口径组织 \nref{S0291} \\
兰州大学物理学院（2026） & 学院统一举办暑期夏令营 \nref{S0290} \\
厦门大学物理科学与技术学院（2026） & 学院拥有凝聚态物理国家重点学科，统一举办夏令营 \nref{S0103} \\
山东大学物理学院（2026） & 物理学优秀大学生暑期夏令营，学院统一口径 \nref{S0385} \\
\bottomrule
\end{tabularx}
\end{center}\vspace{2pt}\par

要判断某方向难不难，看它当年是单独成营还是并入大营、以及是否单列直博名额。
\nref{S0292}
\nref{S0384}
\nref{S0291}
\nref{S0290}
\nref{S0103}
\nref{S0385}
\end{qa}

\iin{physics.ucas.ac.cn}
\ith{学科差异-物理}
\ist{学科与院校差异}
\begin{qa}{S0285}{金融专硕保研的竞争程度和录取偏好是怎样的？}
\textbf{判断：制度上能确认的是金融专硕是一条独立招生口径（单独发文、单独办营，有的还再分方向）；竞争程度与录取偏好没有官方定义，所查资料中也未找到可比数值。}

\par\vspace{2pt}\nobreak\begin{center}\footnotesize
\begin{tabularx}{\textwidth}{>{\raggedright\arraybackslash}X>{\raggedright\arraybackslash}X}
\toprule
\textbf{单位（年份）} & \textbf{原文口径} \\ \midrule
中国人民大学财政金融学院（2027） & 把直博项目与专业学位硕士项目拆成两份报名通知分别发布 \nref{S0018}\nref{S0019} \\
清华大学五道口金融学院（2025） & 发布金融专业硕士研究生夏令营报名说明，专硕单独办营 \nref{S0382} \\
上海财经大学金融学院（2027） & 发布接收推荐免试研究生（含直博生）预报名通知 \nref{S0330} \\
厦门大学经济学科（2026） & 设第十二届全国优秀大学生金融硕士（含人工智能金融方向）夏令营，专硕内部再分方向 \nref{S0329} \\
\bottomrule
\end{tabularx}
\end{center}\vspace{2pt}\par

可用的客观口径只有该院当年拟录取公示的人数与生源结构，以及报名条件里是否设排名档位。
\nref{S0018}
\nref{S0019}
\nref{S0382}
\nref{S0330}
\nref{S0329}
\end{qa}

\iin{中国人民大学}
\ith{学科差异-金融}
\ist{学科与院校差异}
\begin{qa}{S0312}{中科院各研究所的保研难度梯度大致是怎样的？}
\textbf{判断：未找到把中科院各研究所按难度排序的官方文件，任何梯度表都无制度依据。可核实的差异在各所自己的招生条件与组织方式上。}

\par\vspace{2pt}\nobreak\begin{center}\footnotesize
\begin{tabularx}{\textwidth}{>{\raggedright\arraybackslash}X>{\raggedright\arraybackslash}X}
\toprule
\textbf{单位（年份）} & \textbf{原文口径} \\ \midrule
中国科学院心理研究所（2026） & 优先考虑心理学、生理学、生物化学、医学、计算机科学、人工智能等专业学生；营员名单由教育处依据导师的预选意见进行审核；学员自己解决住宿及往返交通，研究所发放适量住宿补贴 \nref{S0298} \\
中国科学院数学与系统科学研究院（2026） & 招生计划130人左右；京外高校学生线上、京内高校学生线下参加；营期组织推免生面试 \nref{S0074} \\
中国科学院化学研究所（2026） & 以所为单位自行发布夏令营报名通知 \nref{S0211} \\
中国科学院微电子研究所（2026） & 举办王守武微电子大学生夏令营 \nref{S0373} \\
中国科学院大气物理研究所（2026） & 以所为单位自行发布夏令营通知 \nref{S0208} \\
中国科学院信息工程研究所（2026） & 举办网络空间安全技术全国优秀大学生夏令营 \nref{S0219} \\
中国科学院软件研究所（2026） & 发布软件与网络夏令营网报公告 \nref{S0223} \\
中国科学院沈阳自动化研究所（2026） & 举办机器人与智能制造夏令营 \nref{S0056}；西北生态环境资源研究院2026年亦独立办营 \nref{S0054} \\
\bottomrule
\end{tabularx}
\end{center}\vspace{2pt}\par

比难度更该比的是：是否导师预选、招多少营员、报名门槛写多细。
\nref{S0298}
\nref{S0074}
\nref{S0211}
\nref{S0373}
\nref{S0208}
\nref{S0219}
\nref{S0223}
\nref{S0056}
\nref{S0054}
\end{qa}

\iin{中国科学院}
\ith{院校类型-中科院}
\ist{学科与院校差异}
\begin{qa}{S0313}{中科院京内所和京外所在保研难度和待遇上差别多大？}
\textbf{结论：待遇口径各所自定，同一座城市里的两个所就可能完全不同；难度梯度没有官方依据，也不按京内京外划线。}

\par\vspace{2pt}\nobreak\begin{center}\footnotesize
\begin{tabularx}{\textwidth}{>{\raggedright\arraybackslash}X>{\raggedright\arraybackslash}X}
\toprule
\textbf{单位（年份）} & \textbf{待遇原文口径} \\ \midrule
中国科学院数学与系统科学研究院（京内，2026） & 京内高校学生采取线下活动形式，入选的京内营员将获得夏令营活动期间的免费餐食和交通补助；京外高校学生线上参加 \nref{S0074} \\
中国科学院心理研究所（京内，2026） & 学员自己解决住宿及往返交通，研究所将发放适量住宿补贴 \nref{S0298} \\
中国科学院大学杭高院生命与健康科学学院（杭州，2026） & 报销从学校来杭的单程路费；营期提供免费餐食，住宿由院统一安排 \nref{S0007} \\
中国科学院长春应用化学研究所（长春，2027） & 用问卷星预报名，研究所初审后另行通知面试，未接到通知即未获面试资格 \nref{S0057} \\
西北生态环境资源研究院（兰州，2026） & 以院所名义独立办营 \nref{S0054} \\
\bottomrule
\end{tabularx}
\end{center}\vspace{2pt}\par

可见京内所待遇好是说不通的：同在京内，数学院给交通补助，心理所只给住宿补贴。
\nref{S0074}
\nref{S0298}
\nref{S0007}
\nref{S0057}
\nref{S0054}
\end{qa}

\iin{中国科学院}
\ith{院校类型-中科院}
\ist{学科与院校差异}
\begin{qa}{S0314}{国科大统一招生和研究所自主招生的推免有什么区别？}
\textbf{结论：判断统一招生还是研究所自主招生，看通知抬头写的是谁。}

\par\vspace{2pt}\nobreak\begin{center}\footnotesize
\begin{tabularx}{\textwidth}{>{\raggedright\arraybackslash}X>{\raggedright\arraybackslash}X}
\toprule
\textbf{抬头写法（年份）} & \textbf{学籍与招生归属} \\ \midrule
中国科学技术大学天文与空间科学学院（紫金山天文台）（2027） & 录取研究生属于中国科大学籍；第一年在中国科大本部上课，第二年起在紫台（南京）科研实践与学位论文，发中国科大证书；报名走中科大研究生报考服务系统，学院代码170 \nref{S0294} \\
中国科大应用化学与工程学院（长春应化所）（2027） & 以中国科大学院名义接收硕士与直博生；研究所组织初审与面试 \nref{S0057} \\
中国科学院大学心理学系 中国科学院心理研究所（2026） & 双署名发布招生简章，走国科大报名系统 \nref{S0298} \\
国科大杭州高等研究院生命与健康科学学院（2026） & 国科大下属办学机构自行发布夏令营通知 \nref{S0007} \\
中国科学院化学研究所（2026） & 研究所属名发布夏令营通知，由研究所教育处组织 \nref{S0211} \\
\bottomrule
\end{tabularx}
\end{center}\vspace{2pt}\par

判断路径：读抬头，再读学籍与培养地点条款，最后在教育部推免系统里核对学院代码。
\nref{S0294}
\nref{S0057}
\nref{S0298}
\nref{S0007}
\nref{S0211}
\end{qa}

\iin{中国科学院}
\ith{院校类型-中科院}
\ist{学科与院校差异}
\begin{qa}{S0325}{港校在内地的研究生招生和推免有什么区别？}
\textbf{判断：这一方向未找到港校一手招生文件，无法给出制度对比。}

所查资料中唯一相关的成文事实是教育部办公厅《关于对新增推免资格高校予以备案的通知》公布的67所新增推免资格高校名单，其中包含香港中文大学（深圳）\nref{S0640}。这份名单的成员是内地高校与内地办学机构，香港高校本部（如香港大学、香港中文大学本部、香港科技大学）不在该名单之内。

由此可以确定的制度边界是：内地推免资格由教育部备案名单确定，港校本部不属于该体系；港校在内地的研究生招生按各校自主申请办法进行，与教育部推免服务系统不互通。你手里的推免资格对港校本部招生不构成制度性效力，港校推免这种说法本身不准确。

具体申请条件、时间与奖学金安排，须以各港校研究生院官网当年公告为准，本轮未找到该类文件。
\nref{S0640}
\end{qa}

\iin{中国政府网}
\ith{院校类型-港澳}
\ist{学科与院校差异}
\begin{qa}{S0343}{四川大学、重庆大学、电子科大的西南保研生态是怎样的？}
\textbf{结论：能核实的是各校自己的招生规模与校级定规则、学院分散发布的结构；重庆大学本轮未找到推免一手文件，生态这类说法无制度依据。}

四川大学：
\begin{itemize}
  \item 2027年章程：预计招收推免硕士3150名、直博生810名；招生专业目录在研招网发布，目录招生人数为预计招生人数，实际录取人数将视申请及考核情况进行调整 \nref{S0476}；
  \item 各招生院系所分别发布接收通知，形成一份汇总清单 \nref{S0475}；
  \item 学院级通知可查：人工智能学院2026夏令营 \nref{S0073}、网络空间安全学院2026夏令营 \nref{S0083}、建筑与环境学院2027预报名 \nref{S0070}；
  \item 另有四川大学-南方电网联合培养项目推免招生简章，专项单独成文 \nref{S0478}。
\end{itemize}

电子科技大学：
\begin{itemize}
  \item 《接收2027年优秀应届本科毕业生免试攻读研究生（含直博生）通知》：直博生学制5年；金融硕士学制2年，其他硕士专业3年；强基计划转段、工程硕博士培养改革等专项按上级要求另行通知 \nref{S0495}。
\end{itemize}

重庆大学：所查资料中仅见研究生招生信息网页面，无当年推免文件。
\nref{S0476}
\nref{S0475}
\nref{S0073}
\nref{S0083}
\nref{S0070}
\nref{S0478}
\nref{S0495}
\end{qa}

\iin{四川大学}
\ith{院校特定-西南}
\ist{学科与院校差异}
\section{优营效力}
\begin{qa}{E0350}{XX 大学的计算机夏令营和 XX 大学比，哪个更难进？}
\textbf{判断：不同学校夏令营难不难进，不能凭校名比，要用各院当年通知里三个可量化字段比——报名条件的成绩档位、计划营员数、是否含测评环节。}

\par\vspace{2pt}\nobreak\begin{center}\footnotesize
\begin{tabularx}{\textwidth}{>{\raggedright\arraybackslash}X>{\raggedright\arraybackslash}X}
\toprule
\textbf{单位（年份）} & \textbf{可比字段原文} \\ \midrule
北京交通大学计算机学院（2026） & 报名条件分四档：本校应届前40\%；同类高校应届前40\%；具有相关优势学科的高校应届前10\%；或竞赛全国二等奖以上、第一发明人专利、第一作者A类论文 \nref{S0104} \\
上海交通大学计算机学院（2026） & 学院2025年由原计算机科学与工程系、软件学院、网络空间安全学院、约翰·霍普克罗夫特计算机科学中心、人工智能研究院等重组而成，以学院为单位发布夏令营活动通知 \nref{S0358} \\
北京大学计算机学院（2026） & 面向计算机科学与技术学科优秀大学生，2026年7月4日至7日线下举行，地址北京大学燕园校区，活动含专业知识与技能测评 \nref{S0114} \\
天津大学计算机科学与技术学院（2026） & 以学院为单位发布暑期夏令营活动通知 \nref{S0118} \\
东南大学计算机科学与工程学院、软件学院（2025） & 以学院为单位发布夏令营报名通知，入选名单由学院官网公布 \nref{S0126} \\
湖南大学计算机学院（2027） & 发布接收推荐免试研究生（含直博生）预报名通知，与夏令营分批衔接 \nref{S0129} \\
\bottomrule
\end{tabularx}
\end{center}\vspace{2pt}\par

把这几栏抄进同一张表，你就能看出哪个院的硬条件更紧。
\nref{S0104}
\nref{S0358}
\nref{S0114}
\nref{S0118}
\nref{S0126}
\nref{S0129}
\end{qa}

\iin{计算机科学与技术学院}
\ith{院校对比}
\ist{院校与项目差异}
\begin{qa}{E0354}{有的学校夏令营不发 offer 只给优秀营员，这个有用吗？}
\textbf{结论：有用，但有用的含金量完全取决于该院当年通知怎么写。同样是优秀营员，有的院写着择优录取，有的院写着招生咨询服务对象。}

\par\vspace{2pt}\nobreak\begin{center}\footnotesize
\begin{tabularx}{\textwidth}{>{\raggedright\arraybackslash}X>{\raggedright\arraybackslash}X}
\toprule
\textbf{单位（年份）} & \textbf{优秀营员效力原文} \\ \midrule
中国科学院数学与系统科学研究院（2026） & 夏令营期间将组织推免生面试，优秀营员有机会被优先录取为免试硕士研究生或直博生 \nref{S0074} \\
上海交通大学人工智能学院（2026） & 优秀营员如获得所在学校推荐免试研究生资格，申请我院并符合我校推免生招生条件，学院将根据接收规模从高分到低分择优录取，额满为止 \nref{S0514} \\
上海交通大学上海高级金融学院（2026） & 获得高金优秀营员资格的考生必须按学院规定的具体要求在教育部系统报名，完成复试及录取信息确认等程序，否则视同放弃 \nref{S0271} \\
江南大学（2027） & 报名我校2026年暑期夏令营并获得优秀营员的学生仍须登录我校预报名系统进行预报名 \nref{S0655} \\
天津大学化工学院（2026） & 未接到我院入营通知或者未获得优秀营员的，不影响参加我院后续9月推免的招生申请 \nref{S0304} \\
山东师范大学心理学部（2026） & 对综合表现突出的营员评定为优秀营员，作为后续重点联系和招生咨询服务对象 \nref{S0296} \\
\bottomrule
\end{tabularx}
\end{center}\vspace{2pt}\par

看到择优录取、已排序的才算数；写咨询服务或须再预报名的，基本只是一张名片。
\nref{S0074}
\nref{S0514}
\nref{S0271}
\nref{S0655}
\nref{S0304}
\nref{S0296}
\end{qa}

\iin{中国科学院}
\ith{优营效力}
\ist{院校与项目差异}
\begin{qa}{E0355}{“优秀营员”到底等不等于预录取？}
\textbf{结论：不等于预录取。同一个优秀营员名称下，所查资料中至少能查到六种不同效力，务必看该单位当年简章原文。}

\par\vspace{2pt}\nobreak\begin{center}\footnotesize
\begin{tabularx}{\textwidth}{>{\raggedright\arraybackslash}X>{\raggedright\arraybackslash}X>{\raggedright\arraybackslash}X}
\toprule
\textbf{效力类型} & \textbf{实例（单位与年份）} & \textbf{原文口径} \\ \midrule
纯荣誉证书 & 四川外国语大学国际关系学院（2026） & 优秀营员可获得优秀营员荣誉证书 \nref{S0322} \\
表彰并颁发证书 & 江南大学数学与数据科学学院（2026） & 学院将对表现优秀的营员颁发优秀营员证书 \nref{S0258} \\
需再走流程 & 上海交通大学上海高级金融学院（2026） & 获得高金优秀营员资格的考生必须在教育部系统报名，完成复试及录取信息确认等程序，否则视同放弃 \nref{S0271} \\
有机会被优先录取 & 中国科学院数学与系统科学研究院（2026） & 优秀营员有机会被优先录取为免试硕士研究生或直博生 \nref{S0074} \\
仅作为下一轮门票 & 江南大学（2027） & 获得优秀营员的学生仍须登录我校预报名系统进行预报名 \nref{S0655} \\
与录取无关 & 西安交通大学网络空间安全学院（2026） & 学院不组织统一的优秀营员面试或评选活动；优秀实习生仅作为本次实习的荣誉认定，不涉及任何与推免、考研、奖学金评定或录取相关的优先政策 \nref{S0132} \\
干脆不评选 & 华东师范大学信息与电子工程学院（2026）、北京师范大学教育学部（2026） & 本次夏令营不评选优秀营员 \nref{S0217}\nref{S0401} \\
\bottomrule
\end{tabularx}
\end{center}\vspace{2pt}\par

判断方法只有一步：读该院当年通知里优秀营员那一段的原文。
\nref{S0322}
\nref{S0258}
\nref{S0271}
\nref{S0074}
\nref{S0655}
\nref{S0132}
\nref{S0217}
\nref{S0401}
\end{qa}

\iin{国际关系学院}
\ith{优营效力}
\ist{院校与项目差异}
\begin{qa}{E0370}{有的学校是不是规定必须保内？}
\textbf{结论：普通推免里必须保内是违规做法；只有特殊专项计划才在政策允许范围内限制报考本校。}

全国统一依据是教育部办公厅《关于进一步完善推荐优秀应届本科毕业生免试攻读研究生工作办法的通知》（教学厅〔2014〕5号）\nref{S0048}\nref{S0205}，其中两句话最直接：

\begin{itemize}
  \item 所有推免生均享有依据招生政策自主选择报考招生单位和专业的权利，所有推免名额（除有特殊政策要求的专项计划外）均可向其他招生单位推荐；推荐高校要充分尊重并维护考生自主选择志愿的权利，不得将报考本校作为遴选推免生的条件，也不得以任何其他形式限制推免生自主报考；
  \item 2014年起，教育部下达推免名额时不再区分学术学位和专业学位，不再设置留校限额；推荐高校也不得对本校推免名额限制学术学位与专业学位报考类型，不得自行设置留校限额或名额。
\end{itemize}

校级文件也有重申：南开大学推免工作实施办法（南党发〔2021〕63号）规定各学院不得将报考本校作为遴选推免生的限制条件，也不得以其他任何形式限制推免生自主报考 \nref{S0561}。

政策允许的例外：上海财经大学2027年操作说明写明通过支教团专项计划遴选的推免生只能选择报考本校专业志愿 \nref{S0183}。
\nref{S0048}
\nref{S0205}
\nref{S0561}
\nref{S0183}
\end{qa}

\iin{中华人民共和国教育部}
\ith{保内保外}
\ist{院校与项目差异}
\begin{qa}{E0383}{有的学校要求九推复试必须到校参加，路远怎么办？}
\textbf{结论：要求到校是院系自定事项，同一时期不同院系差异极大：2027年北京大学写现场复试，同年中国社会科学院大学写网络远程。路远的解法是提前读复试形式条款，而不是事后协调。}

\par\vspace{2pt}\nobreak\begin{center}\footnotesize
\begin{tabularx}{\textwidth}{>{\raggedright\arraybackslash}X>{\raggedright\arraybackslash}X}
\toprule
\textbf{单位（年份）} & \textbf{复试形式原文} \\ \midrule
北京大学（2027办法，校本部） & 推免接收复试工作采取现场复试形式进行，学生签署诚信承诺书 \nref{S0682} \\
上海交通大学船舶海洋与建筑工程学院（2027预推免办法） & 考核时间拟定9月17至18日，地点上海闵行；考生往返交通及食宿自理 \nref{S0575} \\
中国科学院紫金山天文台（2027） & 面试时间初定9月中下旬，采取现场面试方式进行（地点南京）\nref{S0294} \\
中国社会科学院大学经济学院（2027） & 复试考核采取网络远程方式进行 \nref{S0308} \\
中国社会科学院大学法学院（2027） & 拟采用网络远程复试的方式进行，采取面试考核形式 \nref{S0243} \\
中国科学院数学与系统科学研究院（2026） & 拟定京外高校学生线上，京内高校学生线下参加 \nref{S0074} \\
\bottomrule
\end{tabularx}
\end{center}\vspace{2pt}\par

能不能远程，看该院当年一句原文；同期撞车时按现场的优先排行程。
\nref{S0682}
\nref{S0575}
\nref{S0294}
\nref{S0308}
\nref{S0243}
\nref{S0074}
\end{qa}

\iin{北京大学}
\ith{复试地点}
\ist{院校与项目差异}
\begin{qa}{E0384}{各校待录取通知的确认时限不一样吗？最短的是多少？}
\textbf{结论：各校时限差别很大，从30分钟到24小时都有；本轮所查资料中见到的最短是30分钟（南京大学商学院2026年）。}

\par\vspace{2pt}\nobreak\begin{center}\footnotesize
\begin{tabularx}{\textwidth}{>{\raggedright\arraybackslash}X>{\raggedright\arraybackslash}X}
\toprule
\textbf{单位（年份）} & \textbf{待录取确认时限原文} \\ \midrule
南京大学商学院（2026） & 请收到待录取通知的申请人于30分钟内完成确认（9月25日14:00前收到通知的以14:00为起始时间，14:00后收到的以收到通知时间为起始时间），未按时确认视为自动放弃并启动补录 \nref{S0282} \\
大连理工大学计算机科学与技术学院（2027） & 申请者请在1小时内完成待录取确认，超过规定时间取消待录取资格 \nref{S0106} \\
大连理工大学商学院（2027） & 9月24日14:00前发出的待录取通知须在15:00前确认；14:00后发出的须在待录取发出后1小时内确认 \nref{S0550} \\
中国社会科学院大学经济学院（2027） & 应在待录取通知发出后2小时内确认 \nref{S0308} \\
西安电子科技大学计算机学院（2027） & 在接到待录取通知且确认功能开通2小时内完成待录取确认 \nref{S0120} \\
哈尔滨师范大学（2027） & 须在12小时内确认接受我校拟录取，逾期未确认者取消其拟录取资格 \nref{S0058} \\
华南理工大学（2027） & 必须在学院（系）规定的时间内（最长不超过24小时）通过推免服务系统予以确认 \nref{S0479} \\
西南财经大学（2027） & 须在24小时内通过系统确认接受待录取通知 \nref{S0483} \\
\bottomrule
\end{tabularx}
\end{center}\vspace{2pt}\par

全国层面：志愿提交后48小时内不允许修改，48小时后仍未接到复试通知可继续填报其他志愿 \nref{S0388}\nref{S0444}。
\nref{S0282}
\nref{S0106}
\nref{S0550}
\nref{S0308}
\nref{S0120}
\nref{S0058}
\nref{S0479}
\nref{S0483}
\nref{S0388}
\nref{S0444}
\end{qa}

\iin{南京大学}
\ith{时限差异}
\ist{院校与项目差异}
\begin{qa}{E0385}{有的学校要求提前签承诺书才发复试通知，正常吗？}
\textbf{结论：签诚信承诺书是常见且有依据的环节，通常出现在报名材料或复试阶段，内容是诚信与材料真实性；如果承诺书要求你放弃其他学校录取，那就超出了常规诚信承诺的范围。}

\par\vspace{2pt}\nobreak\begin{center}\footnotesize
\begin{tabularx}{\textwidth}{>{\raggedright\arraybackslash}X>{\raggedright\arraybackslash}X}
\toprule
\textbf{单位（年份）} & \textbf{承诺书相关原文} \\ \midrule
北京大学（2027办法） & 推免接收复试工作采取现场复试形式进行，学生签署诚信承诺书 \nref{S0682} \\
南京大学（2027办法） & 附件2为南京大学考生诚信复试承诺书 \nref{S0505} \\
四川大学（2027章程） & 复试期间应自觉遵守考核规则，签订诚信考试承诺书，违反规定或破坏考核秩序的将取消相关成绩并按国家相关规定进行处理 \nref{S0476} \\
中国社会科学院大学法学院（2027） & 申请材料清单含中国社会科学院大学2027年推免招生考试考生诚信承诺书 \nref{S0243}；经济学院为同类要求 \nref{S0308} \\
江南大学（2027预报名公告） & 申请材料含诚信承诺书（附件2，本人签名）\nref{S0655} \\
华东师范大学（2027预报名通知） & 报名材料含诚信承诺书（报名系统下载打印、签字后扫描上传）\nref{S0610} \\
\bottomrule
\end{tabularx}
\end{center}\vspace{2pt}\par

另需注意：确认接受待录取后不得随意撤销，如上海财经大学2027年操作说明写一经确认接受待录取通知，我校不再接受撤销待录取申请 \nref{S0183}。
\nref{S0682}
\nref{S0505}
\nref{S0476}
\nref{S0243}
\nref{S0308}
\nref{S0655}
\nref{S0610}
\nref{S0183}
\end{qa}

\iin{北京大学}
\ith{额外要求}
\ist{院校与项目差异}
\begin{qa}{E0388}{学科评估 A 类的学校一定比 B 类难考吗？}
\textbf{结论：学科评估等级确实被写进部分单位的申请条件，但它是资格线，不是难度排序；A类学科的院系也可能因为招生规模大而更好进。}

\par\vspace{2pt}\nobreak\begin{center}\footnotesize
\begin{tabularx}{\textwidth}{>{\raggedright\arraybackslash}X>{\raggedright\arraybackslash}X}
\toprule
\textbf{单位（年份）} & \textbf{学科评估在文件中的用法} \\ \midrule
西安交通大学网络空间安全学院（2026） & 把它当成放宽排名的条件：所在学校为双一流大学建设高校的，综合成绩排名须在专业前30\%；所在学校专业为双一流建设学科或最新一轮学科评估A-及以上的，排名放宽至前15\% \nref{S0132} \\
山东大学控制科学与工程学院（2026） & 在学院介绍中写控制科学与工程学科被评为A类学科，作为学科实力的陈述 \nref{S0102} \\
厦门大学经济学科（2026） & 在报名指南中写理论经济学、应用经济学、统计学三个一级学科获评教育部最新一轮学科评估A类学科 \nref{S0314} \\
\bottomrule
\end{tabularx}
\end{center}\vspace{2pt}\par

注意这三处用法完全不一样——第一处是硬条件，后两处是宣传陈述。同一句话出现在申请条件段落还是学院简介段落，含义天差地别。
\nref{S0132}
\nref{S0102}
\nref{S0314}
\end{qa}

\iin{西安交通大学}
\ith{学科评估}
\ist{院校与项目差异}
\section{信息核实}
\begin{qa}{E0351}{清华北大的预推免是不是基本只招本校生？}
\textbf{结论：两校的校级文件都明文包含外校申请人的申请办法，不存在只招本校生的表述；差异只体现在材料提交范围上。}

\par\vspace{2pt}\nobreak\begin{center}\footnotesize
\begin{tabularx}{\textwidth}{>{\raggedright\arraybackslash}X>{\raggedright\arraybackslash}X}
\toprule
\textbf{学校（年份）} & \textbf{原文口径} \\ \midrule
清华大学（2027年接收推免生报名通知） & 符合申请条件的申请人（包括本校应届生）登录清华大学研究生申请服务系统填报信息；申请条件为在目前就读学校取得教育部推荐免试研究生资格的优秀应届本科毕业生；接收类型为推荐免试攻读硕士学位、博士学位研究生；申请材料第2、5、6项（在学证明、专家推荐信两封、本科学业成绩单含排名证明）本校应届本科毕业生不需提供 \nref{S0491} \\
北京大学（2027年接收推荐免试研究生办法（校本部）） & 申请人基本条件仅为拥护中国共产党的领导、品德良好、遵纪守法，以及获得所在学校推荐免试资格；且北京大学本校推免生申请攻读本院系研究生可提交申请表、个人陈述、外语水平证明三项各一份，外校及本校跨院系申请的推免生须向招生院系研究生教务办公室提交以上全部申请材料 \nref{S0682} \\
\bottomrule
\end{tabularx}
\end{center}\vspace{2pt}\par

真正该关注的是院系里放多少名额给外校。
\nref{S0491}
\nref{S0682}
\end{qa}

\iin{清华大学}
\ith{院校特点}
\ist{院校与项目差异}
\begin{qa}{E0352}{某校某院是不是只招直博？怎么确认？}
\textbf{判断某院是不是只招直博，最直接的办法是看该院通知里招生类型与招生专业那一段，再看标题和教育部系统里的学位层次。}

\par\vspace{2pt}\nobreak\begin{center}\footnotesize
\begin{tabularx}{\textwidth}{>{\raggedright\arraybackslash}X>{\raggedright\arraybackslash}X}
\toprule
\textbf{单位（年份）} & \textbf{原文口径} \\ \midrule
中国人民大学财政金融学院（2027） & 把直博项目与专业学位硕士项目拆成两份报名通知分别发布，报名前必须看清自己报的是哪一份 \nref{S0018}\nref{S0019} \\
中国社会科学院大学法学院（2027） & 明写我院2027年拟接收优秀应届本科生免试攻读法学硕士学位研究生11人、免试攻读博士学位研究生（直博生）2人，直博另设条件 \nref{S0243} \\
中国社会科学院大学经济学院（2027） & 经济学院各二级学科均接收推荐免试攻读硕士学位研究生；直博只在经济史、西方经济学、发展经济学和数字经济学专业接收 \nref{S0308} \\
长安大学建筑学院（2027） & 标题即2027年推荐免试研究生（含直博生）招生工作细则，标题带含直博生说明硕士也招 \nref{S0024} \\
中南大学（2027直博通知） & 直博生指导性招生计划400名，其中含强基转段等；二级招生单位各专业直博生招生人数一般不超过本专业招生人数的25\% \nref{S0625} \\
\bottomrule
\end{tabularx}
\end{center}\vspace{2pt}\par

核实四步：读招生类型段，看标题，查教育部推免系统里的学位层次字段，电话问学院教务并留书面答复。
\nref{S0018}
\nref{S0019}
\nref{S0243}
\nref{S0308}
\nref{S0024}
\nref{S0625}
\end{qa}

\iin{中国人民大学}
\ith{项目规则}
\ist{院校与项目差异}
\begin{qa}{E0353}{某校的预推免是不是走过场，还是真刷人？}
\textbf{判断预推免是不是走过场，看两样东西：该院往年9月下旬有没有出考核结果公示，以及考核结果里有没有递补档位。有，就是真刷人。}

\par\vspace{2pt}\nobreak\begin{center}\footnotesize
\begin{tabularx}{\textwidth}{>{\raggedright\arraybackslash}X>{\raggedright\arraybackslash}X}
\toprule
\textbf{单位（年份）} & \textbf{可查的原文材料} \\ \midrule
大连理工大学商学院（2027） & 发布接收推荐免试攻读硕士学位研究生（预推免）复试考核结果公示，2026年9月20日 \nref{S0551} \\
东南大学计算机科学与工程学院、软件学院（2027） & 分别发布接收推荐免试研究生考核结果、考核入选名单 \nref{S0116}\nref{S0117}\nref{S0127} \\
中南大学计算机学院（2027） & 发布优秀本科生拟预推免攻读硕士学位研究生名单公示 \nref{S0124} \\
吉林大学计算机科学与技术学院（2027） & 公示原文写经学院组织复试，现对2027年我院接收推荐免试研究生拟录取名单进行公示，公示时间2026年9月28至30日 \nref{S0086} \\
西安电子科技大学计算机学院（2027） & 方案写明我院推免生拟录取结果均向社会公示，公示期不少于7个日，最终录取名单以推免服务系统为准 \nref{S0120} \\
大连理工大学商学院（2027） & 章程中同样出现待录取、递补录取档 \nref{S0550} \\
\bottomrule
\end{tabularx}
\end{center}\vspace{2pt}\par

只出入营名单、不出考核结果的院，才更接近走过场。
\nref{S0551}
\nref{S0116}
\nref{S0117}
\nref{S0127}
\nref{S0124}
\nref{S0086}
\nref{S0120}
\nref{S0550}
\end{qa}

\iin{大连理工大学}
\ith{复试性质}
\ist{院校与项目差异}
\begin{qa}{S0329}{清华大学计算机推免是不是基本只招直博？}
\textbf{判断：未找到清华大学计算机系（或相关院系）当年推免与夏令营简章原文，无法确认该系是否只招直博。}

能核实的校级事实：清华大学2027年接收推免生报名通知面向在目前就读学校取得教育部推荐免试研究生资格的优秀应届本科毕业生，包括本校应届生；接收类型明确写为推荐免试攻读硕士学位、博士学位研究生；我校对通过复试并同意接收为推荐免试攻读硕士学位、博士学位研究生的优秀应届本科毕业生不再发放纸质版初步录取通知（接收函），初步录取结果请申请人查询我校研究生申请服务系统 \nref{S0491}。

清华其他院系确实招推免硕士：清华大学经济管理学院2027年发布接收获推免资格学生攻读硕士学位研究生的报名通知（金融硕士、管理硕士等方向）\nref{S0254}。

所以清华全校只招直博不成立；计算机系是否以直博为主，必须查该系当年通知。本轮未找到该系一手文件。
\nref{S0491}
\nref{S0254}
\end{qa}

\iin{清华大学}
\ith{院校特定-清华}
\ist{学科与院校差异}
\begin{qa}{S0350}{「某校预推免只是走过场」这类说法该怎么核实？}
\textbf{核实预推免只是走过场这类说法，有四步可操作的查证路径，全部基于公开材料，不需要打听。}

\begin{enumerate}
  \item 查该院有没有预推免考核结果公示。例：大连理工大学商学院2027年（预推免）复试考核结果公示（2026年9月20日）\nref{S0551}；东南大学计算机学院、软件学院2027年考核结果与考核入选名单 \nref{S0116}\nref{S0117}\nref{S0127}；中南大学计算机学院2027年拟预推免名单公示 \nref{S0124}。
  \item 查该院预推免办法里有没有差额与递补条款。例：大连理工大学商学院2027年章程出现待录取、递补录取档位 \nref{S0550}；中国社会科学院大学经济学院2027年原则上按不低于120\%的比例差额组织复试考核 \nref{S0308}。
  \item 用公示里的两个人数算通过率。进入复试人数与最终拟录取人数都在公示里，自己除一下就有数。
  \item 查不到就直接问学院研究生教务。教育部教学厅〔2014〕5号要求推免名额、推荐办法、推免生名单、复试录取办法、拟录取推免生名单、咨询申诉渠道全部公开 \nref{S0048}\nref{S0205}，公开渠道本就该能问到。
\end{enumerate}

不要采信某某说这类无出处的口碑。
\nref{S0551}
\nref{S0116}
\nref{S0117}
\nref{S0127}
\nref{S0124}
\nref{S0120}
\nref{S0550}
\nref{S0308}
\nref{S0048}
\nref{S0205}
\end{qa}

\iin{大连理工大学}
\ith{信息核实}
\ist{学科与院校差异}
\begin{qa}{S0351}{怎么判断某校某学院「是否只招直博」？}
\textbf{判断某院是否只招直博，按这四步走，每一步都有原文可依。}

\begin{enumerate}
  \item 读该院当年接收推免生通知或实施细则的招生类型、招生专业段。例：中国社会科学院大学经济学院2027年细则明写各二级学科均接收推荐免试攻读硕士学位研究生，直博只在经济史、西方经济学、发展经济学和数字经济学四个专业接收 \nref{S0308}；法学院2027年明写学硕11人、直博2人 \nref{S0243}。
  \item 看通知标题的括号。标题写含直博生说明硕士也招，例如中国矿业大学2027年《关于做好2027年接收优秀应届本科毕业生免试攻读研究生（含直博生）工作的通知》\nref{S0626}、湖南大学计算机学院2027年预报名通知 \nref{S0129}、长安大学建筑学院2027年招生工作细则 \nref{S0024}；标题只写直博的才可能只招直博。
  \item 在教育部推免服务系统里查该专业的学位层次字段。
  \item 电话或邮件问学院研究生教务，并请对方给书面答复。教育部教学厅〔2014〕5号要求复试录取办法与拟录取名单公开 \nref{S0048}，学院有义务说清。
\end{enumerate}
\nref{S0308}
\nref{S0243}
\nref{S0626}
\nref{S0129}
\nref{S0024}
\nref{S0048}
\end{qa}

\iin{中国社会科学院大学经济学院}
\ith{信息核实}
\ist{学科与院校差异}
\begin{qa}{S0352}{怎么查某校某专业往年夏令营的优营率？}
\textbf{结论：官方不会直接发布优营率这个指标，但你可以用三份公开材料自己拼出分子分母；而且很多营根本不评选优秀营员，这时该指标不存在。}

\par\vspace{2pt}\nobreak\begin{center}\footnotesize
\begin{tabularx}{\textwidth}{>{\raggedright\arraybackslash}X>{\raggedright\arraybackslash}X}
\toprule
\textbf{要用的材料} & \textbf{可查实例（单位与年份）} \\ \midrule
计划营员数（分母参考） & 中国科学院数学与系统科学研究院2026年数学夏令营招生计划为130人左右 \nref{S0074}；天津大学化工学院2026年外校营员招生计划120人左右 \nref{S0304}；华东师范大学信息与电子工程学院2026年计划夏令营营员人数120名 \nref{S0217} \\
优秀营员名单公示（分子） & 山西大学哲学学院2026年通知写明夏令营活动结束后我院将选拔出优秀营员并在学院网站予以公示 \nref{S0534}；江南大学2026年对表现优秀的营员颁发优秀营员证书 \nref{S0258} \\
拟录取名单公示（对照） & 吉林大学计算机学院2027年拟录取名单公示 \nref{S0086}；东南大学计算机学院2027年考核结果 \nref{S0116} \\
\bottomrule
\end{tabularx}
\end{center}\vspace{2pt}\par

注意：有的院明写本次夏令营不评选优秀营员，例如华东师范大学信息与电子工程学院2026年 \nref{S0217}、北京师范大学教育学部2026年 \nref{S0401}；也有院明写不设选拔考核与优秀营员评选 \nref{S0104}。
\nref{S0074}
\nref{S0304}
\nref{S0217}
\nref{S0534}
\nref{S0258}
\nref{S0086}
\nref{S0116}
\nref{S0401}
\nref{S0104}
\end{qa}

\iin{中国科学院}
\ith{信息核实}
\ist{学科与院校差异}
\begin{qa}{S0353}{某校夏令营「入营即优营」的说法可信吗？}
\textbf{结论：这类说法要逐院验证，因为不同单位的原文口径正好相反。}

\par\vspace{2pt}\nobreak\begin{center}\footnotesize
\begin{tabularx}{\textwidth}{>{\raggedright\arraybackslash}X>{\raggedright\arraybackslash}X>{\raggedright\arraybackslash}X}
\toprule
\textbf{口径} & \textbf{单位（年份）} & \textbf{原文} \\ \midrule
明写不设评选与录取承诺 & 北京交通大学计算机科学与技术学院（2026） & 夏令营为招生宣传活动，不设选拔考核、优秀营员评选及任何形式的录取承诺 \nref{S0104} \\
明写不评选优秀营员 & 华东师范大学信息与电子工程学院（2026） & 本次夏令营不评选优秀营员，计划夏令营营员人数120名 \nref{S0217} \\
明写不评选优秀营员 & 北京师范大学教育学部（2026） & 本次暑期夏令营活动不评选优秀营员 \nref{S0401} \\
明写不与名单挂钩 & 湘潭大学机械工程与力学学院（2026） & 夏令营活动不涉及全国硕士研究生招生统考的报名资格预审、预选拔、提前复试等考核活动，不选拔优秀营员，夏令营活动不与招录相关的资格名单或拟录取名单挂钩 \nref{S0235} \\
确实在营期内组织面试 & 中国科学院数学与系统科学研究院（2026） & 夏令营期间将组织推免生面试，优秀营员有机会被优先录取为免试硕士研究生或直博生 \nref{S0074} \\
确实评优营但不等于录取 & 上海交通大学人工智能学院（2026） & 选拔活动中表现优秀者将获得优秀营员身份，优秀营员如获得所在学校推荐免试研究生资格、申请我院并符合我校推免生招生条件，学院将根据接收规模从高分到低分择优录取，额满为止 \nref{S0514} \\
\bottomrule
\end{tabularx}
\end{center}\vspace{2pt}\par

入营即优营多半是把不评选说反了，也可能把结营证明当成了优营。读原文就能定性。
\nref{S0104}
\nref{S0217}
\nref{S0401}
\nref{S0235}
\nref{S0074}
\nref{S0514}
\end{qa}

\iin{计算机科学与技术学院}
\ith{信息核实}
\ist{学科与院校差异}
\section{学科差异-医学}
\begin{qa}{S0280}{机械类保研和车辆、航空等方向保研的差异在哪里？}
“机械”与“车辆、航空”在推免里不是并列学科，而是同一学院下的不同方向，差异写在各院自己发布的细则里，不是校级办法里。

\par\vspace{2pt}\nobreak\begin{center}\footnotesize
\begin{tabularx}{\textwidth}{>{\raggedright\arraybackslash}X>{\raggedright\arraybackslash}X}
\toprule
\textbf{学院（年份）} & \textbf{自设门槛} \\ \midrule
西安交通大学机械工程学院 2026 年夏令营 & 所在学校为 A 类或 B 类“双一流”建设高校且专业前 35\%；或所在专业在最新一轮学科评估中排 B+ 及以上且专业前 5\% \\
浙江大学航空航天学院 2026 年夏令营 & 力学、航空宇航科学与技术、电子科学与技术等方向；申请直博者原则上六级 460 分及以上 \\
上海交通大学机械与动力工程学院 2027 年预推免 & 学院自设申请条件，细则单独挂网，学校只出校级办法 \\
西安交通大学 2027 年 & 机械工程学院与航天航空学院并列挂在“各学院（部、中心）招收推免生实施细则”页，各发一份 \\
\bottomrule
\end{tabularx}
\end{center}\vspace{2pt}\par

所以“机械门槛高还是航空门槛高”问错了对象，该问的是哪个学院的哪一年细则，跨校不可套用。
\nref{S0263}
\nref{S0017}
\nref{S0262}
\nref{S0545}
\end{qa}

\iin{西安交通大学}
\ith{学科差异-机械}
\ist{院校与项目差异}
\begin{qa}{S0295}{临床医学保研中规培基地的选择有多重要？}
结论：在临床、口腔专硕里，选“医院”和选“专业”是同一个动作，二者被绑定在一起。

依据一：中山大学 2027 年接收推免生报名的通知第（八）条第 4 项写明，被录取为临床医学、口腔医学专业学位的学生均须参加住院医师规范化培训；各附属医院均不接受入学前已获得住院医师规范化培训合格证书的学生报考临床（口腔）医学硕士专业学位。

依据二：南开大学医学院 2026 年夏令营通知把“天津地区附属医院可报名专业范围”逐院列表——第一附属医院可报内科学、急诊医学、全科医学、外科学（含骨科方向）、麻醉学、肿瘤学、放射影像学、临床检验诊断学；天津市眼科医院仅眼科学；天津市口腔医院仅口腔医学。

也就是说，同一所医学院内不同附属医院开放的二级学科范围并不相同，报名表里的“报名医院”一栏，实际上就是你未来三年的规培基地。
\nref{S0167}
\nref{S0266}
\end{qa}

\iin{中山大学}
\ith{学科差异-医学}
\ist{院校与项目差异}
\begin{qa}{S0296}{医学影像、麻醉、口腔等小专业保研的特殊性在哪？}
医学小专业的第一道关卡是“本科专业白名单”，而且各校口径并不一致。

中山大学 2027 年接收推免生通知：临床医学硕士专业学位只接受本科毕业于 5 年及以上学制的临床医学、麻醉学、精神医学、医学影像学、放射医学、眼视光医学、儿科学等专业的学生；口腔医学硕士专业学位只接受 5 年制口腔医学专业的应届本科生；1002 临床医学、1003 口腔医学、1006 中西医结合三个学术学位则只接受授医学学位的应届本科生。

南开大学医学院 2026 年夏令营通知的写法是“临床医学和口腔医学专业学位可报名的本科专业必须为临床医学和口腔医学类的五年制学生”，麻醉学、医学影像学没有被单列出来。

两份都是当届文件，但名单一宽一窄。同一所学校的口径尚且会变，跨校更不能套用。
\nref{S0167}
\nref{S0266}
\end{qa}

\iin{中山大学}
\ith{学科差异-医学}
\ist{院校与项目差异}
\begin{qa}{S0297}{中医类院校保研和西医院校的差异是什么？}
这一问目前证据最不均衡：找到的中医类院校文件属于“推荐端”，西医院校文件属于“接收端”，两者不能直接对照评分。

可用的有两点。一是山东中医药大学中医学院 2027 年推免遴选实施细则（2026 年 9 月 4 日发布），依据为教育部教学厅〔2026〕1 号与校学字〔2026〕12 号、13 号，其选拔对象明确不含专升本、第二学士学位专业学生——与综合大学口径一致。

二是接收端，中山大学 2027 年接收推免生通知把 1006 中西医结合与 1002 临床医学、1003 口腔医学并列，学术学位各专业只接受授医学学位的应届本科毕业生。

至于中医专硕与住院医师规范化培训的衔接、中医师承制在推免中的体现，本次未找到可直接对照的一手文件。
\nref{S0639}
\nref{S0167}
\end{qa}

\iin{中医学院}
\ith{学科差异-医学}
\ist{院校与项目差异}
\begin{qa}{S0298}{护理学保研的路径和临床医学有什么不同？}
最硬的差别在学位类别设置。浙江大学医学院 2026 年夏令营通知写明：该院博士专业学位类别为临床医学、口腔医学、公共卫生 3 种，硕士专业学位类别才是临床医学、口腔医学、公共卫生、护理 4 种——护理在该院没有对应的博士专业学位通道。

第二处差别在规培。中山大学 2027 年接收推免生通知第（八）条第 4 项只规定了“被录取为临床医学、口腔医学专业学位的学生均须参加住院医师规范化培训”，护理不在其列。

第三处是招生专业目录。深圳大学医学部 2027 级接收推免公告列出的专业学位有 105100 临床医学、105200 口腔医学、105300 公共卫生、085409 生物医学工程等，护理是否单列要另看当年目录。

以上都是当届文件，不代表学科优劣。
\nref{S0028}
\nref{S0167}
\nref{S0265}
\end{qa}

\iin{浙江大学}
\ith{学科差异-医学}
\ist{院校与项目差异}
\begin{qa}{S0299}{农学、林学、动物医学保研的特点和优势在哪？}
未找到把农学、林学、动物医学放在同一所接收单位下做可对照规定的推免文件，这一问证据薄弱，只能给出三条可查证事实。

一、推荐端有总量也有分配逻辑。甘肃农业大学 2027 年遴选推荐管理办法写明全校推荐名额 433 人，按学科建设、人才培养、应届本科毕业生规模等要素分配。

二、接收端有排名门槛的样本。中国农业大学经济管理学院 2027 年接收推免实施细则规定：国家“双一流”建设高校应届生前 3 学年 GPA 专业年级排名前 20\%，其他取得推免资格高校应届生前 10\%，且不含分校及专升本毕业生。

三、专业身体条件方面，教育部《普通高等学校招生体检工作指导意见》对林学、园林、农学、水产养殖学等列有轻度色觉异常受限条款，动物医学另有听力相关条款。

“农林保研率高”这类说法没有文件依据，本手册不作断言。
\nref{S0059}
\nref{S0684}
\nref{S0039}
\end{qa}

\iin{甘肃农业大学}
\ith{学科差异-医学}
\ist{院校与项目差异}
\begin{qa}{S0319}{独立医科大学保研和综合大学医学院保研的差异在哪？}
差异在“招生指标挂在谁身上”。

综合大学医学院：南开大学医学院 2026 年夏令营通知把招生指标拆到天津地区多家附属医院，逐院列出可报名专业范围（第一附属医院、第一中心医院、第三中心医院、中心妇产科医院、眼科医院、口腔医院），并要求报名专业学位的学生在申请表中注明报名医院。

独立医科大学：大连医科大学 2027 年以“学校”为单位统一发布接收推免工作安排（2026 年 9 月 21 日），网报与复试的时间节点由学校一张表排出，不经学院二次分解。

第三种结构：浙江大学医学院 2027 年招收推免生综合考核细则显示，学院下设 15 个二级招生单位，各单位成立研究生招生委员会，自行负责报考咨询、资格审查及复试录取。

三种结构并存，你报的是哪一种，决定了你的材料由谁初审。
\nref{S0266}
\nref{S0408}
\nref{S0029}
\end{qa}

\iin{南开大学}
\ith{学科差异-医学}
\ist{院校与项目差异}
\begin{qa}{S0357}{医学保研需要提前联系导师吗，还是统一招录？}
医学里不少单位明文写“统一招录、不接受私下联系”，这与理工科导师先行的写法正好相反。

北京大学口腔医学院 2027 年接收推免生复试录取工作通知：各学科组核心成员及招生导师共同组成专家组，根据考生提交的电子版材料集中审核、评分并确定复试名单；同时写明“专业间调剂由教育处统一安排，考生不能私自联系科室参加面试”。

南开大学医学院 2026 年夏令营通知同样把医院分配写进学院统一报名流程。

对照组：西安交通大学网络空间安全学院 2026 年夏令营通知写“申请者通过课题组提供的邮箱联系导师”，这是导师先行的模式。

判断方法只有一个——看目标学院当年通知里，联系导师是写在报名流程里，还是被明文禁止。
\nref{S0340}
\nref{S0266}
\nref{S0132}
\end{qa}

\iin{北京大学口腔医学院}
\ith{学科差异-医学}
\ist{院校与项目差异}
\section{学科差异-数学}
\begin{qa}{S0274}{数学类保研主要看什么，为什么数学专业的保研率普遍较高？}
数学类院系的接收端结构本身就偏向学术型，这一点有硬数据。北京师范大学数学科学学院 2026 年夏令营通知写明：拟招收学术型硕士和直博生营员共 200 人左右，专业型硕士（珠海校区培养）营员 30 人左右。

考核也偏笔试：该营 7 月 7 日下午笔试，科目写明是分析、代数、几何；7 月 8 日面试。

\par\vspace{2pt}\nobreak\begin{center}\footnotesize
\begin{tabularx}{\textwidth}{>{\raggedright\arraybackslash}X>{\raggedright\arraybackslash}X}
\toprule
\textbf{学院（年份）} & \textbf{排名门槛} \\ \midrule
西安交通大学数学与统计学院 2026 年夏令营 & “一流大学建设高校”或“一流学科建设高校”前 20\%；其他高校前 5\% \\
兰州大学数学与统计学院 2026 年夏令营 & 前五个学期成绩排名年级前 40\% \\
北京师范大学数学科学学院 2026 年夏令营 & 一流学科、211、985 及省属重点高校；全国大学生数学竞赛获奖者优先考虑 \\
\bottomrule
\end{tabularx}
\end{center}\vspace{2pt}\par

“保研率高”属于经验判断，本章不作断言。
\nref{S0257}
\nref{S0260}
\nref{S0259}
\end{qa}

\iin{北京师范大学}
\ith{学科差异-数学}
\ist{院校与项目差异}
\begin{qa}{S0275}{数学、统计专业跨保计算机或金融的成功率如何？}
没有官方“成功率”数据，这里只说可查证的通道。

跨计算机方向：西安交通大学数学与统计学院 2026 年夏令营明确招收“数学、统计学及相关理工专业”的三年级本科生，说明数学背景在校内被主动接受；北京大学计算机学院 2026 年夏令营另有自己的申请条件，需单独核对。

跨金融、经济方向：厦门大学经济学科 2026 年系列夏令营的申请条件写“经济类及相关专业”，并额外要求“具备良好的经济学基础和数理基础”，同时限定生源为“第四轮学科评估中理论经济学或应用经济学或统计学排名 A- 及以上高校对应专业”，或“第二轮双一流建设高校和建设学科中相关学科”。

也就是说，数学进金融的通道是开着的，但第一关是本科院校那一关；跨计算机则以各院夏令营列出的“拟招收专业”清单为准。
\nref{S0260}
\nref{S0329}
\nref{S0349}
\end{qa}

\iin{西安交通大学}
\ith{学科差异-数学}
\ist{院校与项目差异}
\begin{qa}{S0281}{材料类保研为什么被认为是「保研率最高的天坑」？}
“保研率最高的天坑”没有文件依据，本手册不作断言。可查证的是材料院系推荐端的流程细节。

中南大学材料科学与工程学院 2027 届预推荐实施细则规定：加分只认学校与学院重点赛事、排行榜重点赛事，或教育部等中央部委、国家一级学会、省教育厅、省科技厅组织的赛事，其他赛事不予加分；个人加分累计不超过 2 分。

该院 2027 届预推荐名单公示后，又发过一次增补公示（2026 年 9 月 17 日）——新增一个工程硕博指标，由候补学生自愿报名后遴选。这说明名单并非一次定死。

接收端样本：上海大学材料科学与工程学院 2027 年出其接收推免（含直博生）实施办法；中国科学技术大学 2027 年工程硕博士培养改革专项在 0856 材料与化工等六个类别拟招 146 人。

“天坑”感来自去向讨论，不是规则本身。
\nref{S0273}
\nref{S0275}
\nref{S0256}
\nref{S0673}
\end{qa}

\iin{中南大学}
\ith{学科差异-数学}
\ist{院校与项目差异}
\begin{qa}{S0302}{心理学保研跨专业申请的成功率如何？}
中国科学院心理研究所 2026 年夏令营的报名条件里直接给出了跨专业优先名单——“优先考虑心理学、生理学、生物化学、医学、计算机科学、人工智能等专业的学生”。这是接收方自己写明的接受范围，比“心理学相关专业”这种模糊表述具体得多。

考核方式上，华东师范大学心理与认知科学学院 2027 年推免生综合考核实施细则规定：综合考核总成绩等于专业素质及综合能力面试成绩的 70\% 加英语能力成绩的 30\%，全部为面试，不含机试与笔试。

可作跨校对照的还有北京师范大学心理学部 2026 年、山东师范大学心理学部 2026 年的夏令营通知。

按专业名单去卡，比按“成功率”估算更可靠。
\nref{S0298}
\nref{S0295}
\nref{S0297}
\nref{S0296}
\end{qa}

\iin{中国科学院}
\ith{学科差异-数学}
\ist{院校与项目差异}
\begin{qa}{S0303}{环境、地学类保研的院校梯度和就业出口如何？}
这一问证据偏弱，只能列出可查证的接收端文件。

中国科学院西北生态环境资源研究院 2026 年举办“相聚丝路兰州 共筑科学梦想”优秀大学生夏令营；兰州大学大气科学学院 2026 年发布优秀大学生夏令营活动通知；浙江大学 2026 年各学院夏令营汇总页中，地球科学学院与环境与资源学院是两个各自独立的报名入口。

这三条说明：地学类的接收单位既包括高校院系，也包括中科院系统的独立研究所，两套体系的报名入口、单位名称和考核口径都不一样。

至于“院校梯度”和“就业出口”，推免文件本身不会写，本次也未找到可用的一手依据，本手册不作断言。
\nref{S0054}
\nref{S0075}
\nref{S0031}
\end{qa}

\iin{中国科学院}
\ith{学科差异-数学}
\ist{院校与项目差异}
\begin{qa}{S0346}{兰州大学、东北大学等「没落名校」的保研性价比如何？}
先给结论：这两类学校的差异主要写在“复试环节”和“排名门槛”上，不在招牌上。

东北大学 2027 年接收推荐免试研究生工作办法：复试以面试考核为主，具体考核内容、形式及安排“关注学院通知”；申请条件只要求原则上通过国家大学外语四级考试（四级成绩 425 分及以上）。

兰州大学的门槛分档很细，且各院系单独发营：数学与统计学院 2026 年夏令营要求前五个学期成绩排名年级前 40\%；物理学院、大气科学学院、政治与国际关系学院、文学院各自另发 2026 年夏令营通知。

对照西安交通大学数学与统计学院 2026 年夏令营的“一流大学或一流学科建设高校前 20\%、其他高校前 5\%”，兰州大学数学与统计学院 2026 年的 40\% 明显更宽。

“性价比”是主观判断，本手册不作断言。
\nref{S0470}
\nref{S0259}
\nref{S0260}
\nref{S0290}
\end{qa}

\iin{东北大学}
\ith{学科差异-数学}
\ist{院校与项目差异}
\section{学科差异-法学}
\begin{qa}{S0283}{建筑学（五年制）保研的时间节点和其他专业有什么不同？}
时间节点的差别主要出在“推荐端算哪几个学年”，而不是九推日历。

甘肃农业大学 2027 年遴选推荐管理办法：前三学年（五年制专业为前四学年）纳入课程平均成绩计算范围的课程考核合格。

西南交通大学 2027 年研究生支教团组建通知同样写：须完成并通过专业免研资格课程，即专业培养方案要求的前三学年（五年制专业为前四学年）开设的全部必修课程及指定的限修课程。

也就是说，五年制建筑学要到第四学年结束才凑齐前四学年成绩进入排名池，比四年制晚一整年，夏令营时间点却与四年制一致。

接收端可对照：长安大学建筑学院 2027 年推免招生细则（2026 年 9 月 9 日）列出学术学位 081300 建筑学、083300 城乡规划学，专业学位 085100 建筑等；浙江大学建筑工程学院 2027 年通知另对直博设六级 460 分门槛。
\nref{S0059}
\nref{S0061}
\nref{S0024}
\nref{S0025}
\end{qa}

\iin{甘肃农业大学}
\ith{学科差异-法学}
\ist{院校与项目差异}
\begin{qa}{S0289}{法学保研中「五院四系」的偏好和普通985有什么不同？}
先说明一点：本手册所引推免文件中，没有任何一份使用“五院四系”这个口径设置条件，所以不能按这个标签比较。可比的是各校自己写的排名门槛与专业目录。

\par\vspace{2pt}\nobreak\begin{center}\footnotesize
\begin{tabularx}{\textwidth}{>{\raggedright\arraybackslash}X>{\raggedright\arraybackslash}X}
\toprule
\textbf{学院（年份）} & \textbf{自设门槛与专业} \\ \midrule
西安交通大学法学院 2026 年夏令营 & 成绩排名本专业前 35\%，有论文、国际级或国家级竞赛获奖、双学位者排名可适当放宽；择优录取约 100 人；含学术学位法学与法律硕士（法学）、（非法学）、（国际仲裁）、（国际法律事务）、知识产权硕士 \\
华中师范大学法学院 2027 年预报名 & 含学术型法学与法律（法学）、法律（非法学）；2026 年 9 月 3 日至 11 日预报名 \\
同济大学法学院 2027 年 & 学院单独发布接收推免预报名通知 \\
中南大学法学院 2027 年 & 发布考核录取工作通知，含复试方式说明 \\
\bottomrule
\end{tabularx}
\end{center}\vspace{2pt}\par

“偏好”是招生结果层面的印象，不能替代逐校查当年通知。
\nref{S0154}
\nref{S0011}
\nref{S0242}
\nref{S0237}
\end{qa}

\iin{西安交通大学}
\ith{学科差异-法学}
\ist{院校与项目差异}
\begin{qa}{S0290}{法硕（非法学）保研对本科专业有偏好吗？}
这一条是全国统一规则，不是院校自定，所以“偏好”这个词并不准确。

教育部《2026 年全国硕士研究生招生工作管理规定》第十六条：报考法律（法学）专业学位的，须符合第十五条各项要求，且报名前所学专业为法学专业，获得法学第二学士学位的人员可报考；报考法律（非法学）专业学位的，须报考前所学专业为非法学专业。

山西大学 2027 年推免接收文件把这条直接落到推免环节：报名参加法律（非法学）硕士复试的推免生，本科阶段所学专业须为非法学专业；参加法律（法学）硕士复试的，本科阶段所学专业须为法学专业，获得法学第二学士学位的可报考。

中山大学 2027 年接收推免生通知第（九）条同款：申请法律（非法学）专业学位，本科专业应为非法学专业。

它是硬性前置条件，不分校。
\nref{S0051}
\nref{S0667}
\nref{S0167}
\end{qa}

\iin{中华人民共和国教育部}
\ith{学科差异-法学}
\ist{院校与项目差异}
\begin{qa}{S0300}{艺术类（美术、音乐、设计）保研的评审标准和普通学科差别在哪？}
在综合大学里，艺术类通常走“专项加分”通道而非独立的评审标准，而且分值可查。

湖南大学 2027 年推免工作通知明确：专项评价成绩包含科研成果、省部级竞赛获奖、荣誉表彰、志愿服务、国际组织实习和参军入伍服兵役、双创三大赛、体育艺术类比赛国家级及以上奖励等。

该校 2027 年加分名单公示给出了实例：教育部全国大学生艺术展演（集体）国家级一等奖、排名第 3，加分 1.67 分。

同一所学校的数学学院 2027 年推免工作细则进一步限定：专项评价成绩（一）不超过 3 分，加分范围与标准见学院附件。

另外，教育部办公厅新增推免资格高校备案名单中列有天津美术学院、广州美术学院、四川音乐学院、四川美术学院等——艺术类院校取得推免资格本身是近年才铺开的事。专业艺术院校独立设置的评审标准，本次未找到。
\nref{S0540}
\nref{S0541}
\nref{S0661}
\nref{S0640}
\end{qa}

\iin{湖南大学}
\ith{学科差异-法学}
\ist{院校与项目差异}
\begin{qa}{S0301}{体育类保研的考核方式和普通学科有什么不同？}
体育类在综合大学里往往有独立的接收单位建制，这是与普通学科最直观的不同。

北京航空航天大学 2027 年各学院接收推免通知汇总页中，“体育部”作为一个独立招生单位，与计算机学院、经济管理学院、集成电路科学与工程学院等并列发布接收通知（2026 年 9 月至 10 月陆续更新）。

西安交通大学研招网的“各学院（部、中心）2027 年招收推免生实施细则”页中同样单列“体育中心”或“体育学院”。

加分侧：湖南大学 2027 年推免生体育类竞赛国家级及以上奖励加分名单中，一名学生以 2024 年第 22 届全国大学生田径锦标赛 4x400 米混合接力第一名加女子 400 米栏第二名，合计加分 8 分。

需要区分的是：教育部《2025 年全国硕士研究生招生工作管理规定》把体育学、体育硕士列入照顾专业单独划线，那是统考规则，不适用于推免。
\nref{S0065}
\nref{S0545}
\nref{S0541}
\nref{S0049}
\end{qa}

\iin{北京航空航天大学}
\ith{学科差异-法学}
\ist{院校与项目差异}
\begin{qa}{S0358}{文科保研中「导师制」和「委员会制」的差别怎么判断？}
判断方法有三步，都能从公开文件看出来。

第一步，看报名系统里有没有“意向导师”字段。北京师范大学数学科学学院 2026 年夏令营通知要求申请人明确填写志愿导师，并在通知中给出导师方向教研室网址——这是导师制信号。

第二步，看复试是集体筛选还是导师个人决定。北京大学历史学系 2027 年研究生推免复试通知写的是“经过历史学系 2027 年研究生推免专家审核小组的严格评审、筛选，确定以下申请者参加推免复试”——这是委员会制信号。

第三步，看入学后的博士招收方式。北京大学历史学系 2027 年研究生招生指南写明该系从 2021 年起采取直博、硕博连读（仅限本系学生）和申请—考核制三种方式招收博士研究生。

把这三步套到中国人民大学国学院 2027 年推免生接收工作报名通知上，同样可以判断。
\nref{S0257}
\nref{S0201}
\nref{S0200}
\nref{S0006}
\end{qa}

\iin{北京师范大学}
\ith{学科差异-法学}
\ist{院校与项目差异}
\section{学科差异-生物}
\begin{qa}{E0363}{医学“四证合一”的项目有哪些？怎么报？}
本手册所引资料全文检索“四证合一”，零命中——未找到任何一个使用该表述的一手招生文件，因此不能给出项目清单，也不能从名称反推报名方式。

只能给出相邻的可查证事实。

一、中山大学 2027 年接收推免生通知第（八）条第 4 项：被录取为临床医学、口腔医学专业学位的学生均须参加住院医师规范化培训；各附属医院均不接受入学前已获得规培合格证书的学生报考临床（口腔）医学硕士专业学位。

二、天津医科大学 2026 年全日制博士研究生招生章程：报考临床医学专业学位博士者，本科所学专业须为临床医学，硕士所学专业或规培专业须与报考专业一致，并已通过国家执业医师资格考试取得临床《医师资格证书》。

这两条说明“学位证加执业医师资格加规培”的绑定关系确实存在于文件里，但“四证合一”这一项目名称未见于本手册收录范围。
\nref{S0167}
\nref{S0184}
\end{qa}

\iin{中山大学}
\ith{学科差异-生物}
\ist{院校与项目差异}
\begin{qa}{S0278}{生物类保研中直博比例为什么特别高，不想直博怎么办？}
生物类直博名额多，是因为基础学科在制度上就被安排成直博的主场。

教育部办公厅 2013 年的推免工作通知给直博设了比例框架：已备案高校的直博招生专业一般为基础学科，人数原则上不超过本校博士生招生规模的 10\%；其他高校可在理工农医门类的国家重点一级学科专业内招收，上限为本专业招生人数的 20\%。这条框架至今仍是理解学科间直博容量差异的起点。

落到操作层：吉林大学生命科学学院 2027 年接收预推免工作通知写明直博生基础学制 5 年，并且申请直博的学生需在报名前先与该院研办老师联系，以确定直博生招生的条件、范围及相关要求——名额不透明是这类方向的特点，靠公告很难推算出竞争强度。

不想直博：直博与硕士推免是两个名额池，上述规定只限定直博的来源人群，并未要求学术型推免生必须转博，报名时只选硕士类别即可。
\nref{S0047}
\nref{S0012}
\end{qa}

\iin{中华人民共和国教育部}
\ith{学科差异-生物}
\ist{院校与项目差异}
\begin{qa}{S0279}{生物、医学类保研中「实验室轮转」经历重要吗？}
对本手册所引材料检索“轮转”，零命中——未找到任何一所学校在推免文件中使用“实验室轮转”这一表述，也未找到把它设为一个选拔环节的规定。

能提供的相邻依据只限于夏令营的内容描述。中山大学生物医学工程学院 2026 年夏令营报名通知、西安交通大学生命科学与技术学院 2026 年夏令营通知，都是本年度生物医学方向的一手文件，可作为你核实当年安排的起点。

但“轮转经历是否重要”在推免规则层面没有依据可查。它更可能出现在入营后的师生交流环节，或入学后的培养方案中，而不在选拔条件里。本手册不对无文件支撑的说法作断言。
\nref{S0078}
\nref{S0323}
\end{qa}

\iin{中山大学}
\ith{学科差异-生物}
\ist{院校与项目差异}
\begin{qa}{S0282}{土木、建筑类保研受行业下行影响，策略要不要调整？}
行业走势不是推免规则，本手册不对策略作断言。可查证的是接收端当年并未停摆。

同济大学土木工程学院 2027 年接收推免（含直接攻博）预报名通知；上海交通大学船舶海洋与建筑工程学院 2027 年接收推免预推免办法；浙江大学建筑工程学院 2027 年接收推免研究生通知；西安建筑科技大学 2027 年接收推免（含直博生）通知——这些都是 2026 年下半年发布的当届文件。

浙江大学建筑工程学院 2027 年通知还给出直博的英语门槛：全国大学英语六级 460 分及以上，或 WSK（PETS 5）合格或英语专业八级合格，或雅思 5.5 及以上或新托福 80 及以上，成绩有效期五年。

要不要调整策略，取决你对出口的判断，而不是文件；文件层面能确认的只是通道仍然开着。
\nref{S0025}
\nref{S0575}
\nref{S0522}
\end{qa}

\iin{浙江大学}
\ith{学科差异-生物}
\ist{院校与项目差异}
\begin{qa}{S0292}{人文类（文史哲）保研中，导师推荐和学术论文怎么权衡？}
没有一份文件给出权重表，可查证的是两条并存的事实。

论文侧：院系把论文当作可选增强材料。北京大学历史学系 2027 年研究生招生指南列出的申请材料为盖章成绩单、获奖证书、外语水平证明等，外语门槛明写 CET6 合格及以上、TOEFL 90 分及以上、IELTS（A 类）6.0 及以上。

推荐侧：同年 9 月 10 日该系发布的 2027 年研究生推免复试通知写的是“经过历史学系 2027 年研究生推免专家审核小组的严格评审、筛选，确定以下申请者参加推免复试”——说明初筛是委员会集体评审，不是某一封推荐信单独决定。

另有中国人民大学国学院 2027 年推免生接收工作报名通知可作对照。

两者不是二选一关系，别用“谁更重”的思路准备。
\nref{S0200}
\nref{S0201}
\nref{S0006}
\end{qa}

\iin{北京大学}
\ith{学科差异-生物}
\ist{院校与项目差异}
\begin{qa}{S0293}{外语类保研中，专四专八成绩和雅思托福哪个更被认可？}
这两类成绩在文件里常常是“或”的关系，可以互相替代。

北京语言大学语言科学与资源学院 2026 年夏令营通知：外语水平良好，达到以下条件之一者优先考虑——大学英语六级 425 分及以上；英语专业四级合格；托福 90 分及以上；雅思 6.5 分及以上。四种并列，专四与托福雅思被放在同一档。

北京林业大学外语学院 2026 年夏令营要求提交“全国英语专业四级（TEM-4）证书或其他英语水平证书扫描件”，“或其他”即承认替代。

工科方向对照：浙江大学信息与电子工程学院 2026 年夏令营要求国家英语六级 480 分及以上（460 至 480 分之间在其他方面有突出表现者也可考虑），或托福 80 分及以上或雅思 5.5 分及以上。

外语类院系把专四与托福雅思并列；工科院系更常只认六级分数线。以目标院校当年通知为准。
\nref{S0021}
\nref{S0023}
\nref{S0036}
\end{qa}

\iin{语言科学与资源学院}
\ith{学科差异-生物}
\ist{院校与项目差异}
\section{名额结构}
\begin{qa}{E0356}{C9 高校里哪些学校九推还会放出名额？}
“哪些 C9 会放名额”没有可穷举的一手依据，但“补录、补充报名”这个机制有明文可查。

哈尔滨工业大学计算学部发布了《关于接收 2027 年优秀应届本科毕业生免试攻读研究生补充报名的通知》——这就是九推阶段仍有名额时的标准动作：由学部单独发一份补充报名通知，而不是重发校级办法。

上海财经大学 2027 年推免生使用推免服务系统的操作说明也写明：本次推免选拔工作结束后，如有个别专业仍有推免招生计划需再次接收学生申请，我校将在推免服务系统中公布并开放相关专业，按批次先申请先审核、额满为止。

所以可行做法是盯“学院、学部级”的补充报名公告，而不是盯学校级别的名单。本次未找到一份可比对的 C9 补录清单，不作断言。
\nref{S0552}
\nref{S0183}
\end{qa}

\iin{哈尔滨工业大学}
\ith{名额结构}
\ist{院校与项目差异}
\begin{qa}{E0364}{金融专硕和学硕在推免里怎么选？}
厦门大学经济学科 2026 年给了最清楚的例子：它把学硕和专硕办成不同夏令营、分别报名。其中之二是王亚南经济研究院第十八届全国优秀大学生经济学（学硕）夏令营，之一是经济学院第十八届全国优秀大学生经济学夏令营，之四是第十二届全国优秀大学生金融硕士（含人工智能金融方向）夏令营。三条通知各自列申请条件，你选错营等于报错类别。

其他学校则合并处理：清华大学五道口金融学院 2025 级金融专业硕士研究生夏令营单独发报名说明；中国人民大学财政金融学院 2027 年把“专业学位硕士项目”和“直博项目”拆成两份报名通知；清华大学经济管理学院 2027 年一份通知同时覆盖金融硕士、管理硕士和清华—哥大双学位项目。

先看清目标院校是分开招还是合并招，再决定策略。
\nref{S0314}
\nref{S0349}
\nref{S0329}
\nref{S0382}
\nref{S0254}
\nref{S0019}
\end{qa}

\iin{厦门大学}
\ith{名额结构}
\ist{院校与项目差异}
\begin{qa}{E0368}{学硕和专硕的推免名额是分开的吗？}
分三层看。

全国层（推荐端）：教育部办公厅 2014 年《关于进一步完善推荐优秀应届本科毕业生免试攻读研究生工作办法的通知》明确，2014 年起教育部下达推免名额时不再区分学术学位和专业学位，不再设置留校限额；推荐高校也不得对本校推免名额限制学术学位与专业学位报考类型。所以你拿到的推荐资格本身不分学硕专硕。

接收端（学院层）：会分。武汉大学计算机学院 2027 年接收推免实施细则写明，拟招收硕士推免生约 101 人（工程硕博国家专项除外），其中学术学位硕士推免生约 53 人、专业学位硕士推免生约 48 人；直博生约 50 人另列。

计划调整规则：教育部《2026 年全国硕士研究生招生工作管理规定》第六十六条——同一招生单位内，学术学位招生计划可调整到专业学位使用，专业学位招生计划不得调整到学术学位专业使用。方向是单向的。
\nref{S0048}
\nref{S0119}
\nref{S0051}
\end{qa}

\iin{中华人民共和国教育部}
\ith{名额结构}
\ist{院校与项目差异}
\begin{qa}{E0369}{直博的名额是另外单列的吗？}
是单列的，但单列在哪一层，不同学校给了不同答案。

推荐端单列：云南师范大学 2027 年推免工作通知把名额分块——“直博生”推免名额为 8 个，原则上给有学术博士学位授权点的单位各分配 1 个名额；同校另有卓越人才培养计划 156 个、专业建设奖励 226 个、研究生支教团 15 个。

接收端单列：武汉大学计算机学院 2027 年实施细则在“硕士推免生约 101 人”之外，单独写“学院选拔招收直博生约 50 人”。

占用规则：华中科技大学 2026 年强基计划学生本研衔接转段接收工作通知写明，转段进入研究生阶段培养的强基生“占用转段接收单位 2027 年研究生招生计划和推免生（直博或硕士）接收计划”。

另有一条全国上限：教育部办公厅 2013 年文件规定，已备案高校直博招生专业一般为基础学科，人数原则上不超过本校博士生招生规模的 10\%。
\nref{S0568}
\nref{S0119}
\nref{S0536}
\nref{S0047}
\end{qa}

\iin{云南师范大学}
\ith{名额结构}
\ist{院校与项目差异}
\begin{qa}{E0390}{教育部直属高校和省属高校的推免名额差多少？}
未找到按“教育部直属高校”与“省属高校”分类的推免名额下达数据，无法给出差额，本手册不编数字。可查证的是名额怎么被切块。

全国规则层面：教育部办公厅 2013 年通知要求，所有推免生名额均可用于向其他招生单位推荐，推免生的留校数量不得高于教育部下达的留校限额（该留校限额自 2014 年起取消）。

学校层面：云南师范大学 2027 年通知给出省属高校的名额结构——研究生支教团 15 个、卓越人才培养计划 156 个、直博生 8 个、专业建设奖励 226 个，并注明各单位分配名额未使用完的由学校收回再分配。甘肃农业大学 2027 年遴选推荐名额为 433 人，按学科建设、人才培养、应届本科毕业生规模等要素分配。

前置门槛方面，教育部办公厅近年备案公布了 67 所新增推免资格高校，其中省属与地方院校占绝大多数。
\nref{S0047}
\nref{S0568}
\nref{S0059}
\nref{S0640}
\end{qa}

\iin{中华人民共和国教育部}
\ith{名额结构}
\ist{院校与项目差异}
\section{学科差异-经管}
\begin{qa}{S0277}{化学、化工保研中本科进组科研经历是不是决定性因素？}
文件写的是“优先考虑”，不是决定性。

天津大学化工学院 2027 年推荐免试攻读硕士学位研究生和直接攻读博士学位研究生招生录取工作办法第（七）条：申请人在本科阶段公开发表过学术论文，或取得科研成果，或在国内外重大竞赛中获奖，可优先考虑。

同一学院 2026 年全国优秀大学生暑期夏令营活动通知是另一份当届文件，可作对照。

也就是说，进组经历在评分表里是加分项而不是准入项。同期文件里更硬的是成绩排名——郑州大学化工学院 2027 年接收推免研究生（含直博生）实施细则与天津大学化工学院 2027 年办法都以排名形式设出前置条件。

把“优先考虑”读成“必须有”是过度解读；但完全不进组，在面试问答环节会明显吃亏。
\nref{S0303}
\nref{S0304}
\nref{S0396}
\end{qa}

\iin{天津大学}
\ith{学科差异-经管}
\ist{院校与项目差异}
\begin{qa}{S0284}{经管类保研中实习经历和科研经历哪个权重更大？}
在找到的经管类文件里，实习出现在“培养环节、就业去向”段落，科研与推荐信出现在“申请条件”段落。这个位置差异是可查证的，但它不等于权重排序。

例一：厦门大学经济学科 2026 年金融硕士夏令营报名指南介绍学院近两年累计选派 500 余名学生赴天津港、白鸽在线、苏州农行等 40 余家单位实习——这是培养描述，不是申请条件。

例二：同一批厦大经济学（学硕）夏令营通知的申请条件里写“具备良好的经济学基础和数理基础”，并要求两位副教授及以上专家的推荐信。

例三：上海交通大学上海高级金融学院 2027 级全日制金融硕士夏令营通知单独发布，申请要求以该院当年通知为准。

建议直接查目标学院当年通知的“申请条件”与“考核办法”两段，权重就写在那里。
\nref{S0329}
\nref{S0349}
\nref{S0272}
\end{qa}

\iin{厦门大学}
\ith{学科差异-经管}
\ist{院校与项目差异}
\begin{qa}{S0286}{会计、审计保研中CPA等证书有没有加分作用？}
对本手册所引材料检索“CPA”与“注册会计师”，未命中任何一所学校把 CPA 列入推免加分项的推免文件，因此本条无正面依据。

反面证据可查：湖南大学数学学院 2027 年推免工作细则把专项评价成绩的加分范围列死为科研成果、竞赛获奖、荣誉表彰、志愿服务、国际组织实习等，专项评价成绩（一）不超过 3 分，具体标准见学院附件；同一学校 2027 年加分名单公示的类别也只有创新创业赛事（挑战杯）、体育类、艺术类、志愿服务四类，没有职业资格证书。

中国农业大学经济管理学院 2027 年接收推免实施细则的申请条件同样只写 GPA 排名——国家“双一流”建设高校应届生前 20\%、其他推免资格高校应届生前 10\%——以及学术研究潜力或实践创新能力，未提任何证书。

结论：在推免环节，CPA 目前找不到文件依据。
\nref{S0661}
\nref{S0541}
\nref{S0684}
\end{qa}

\iin{湖南大学}
\ith{学科差异-经管}
\ist{院校与项目差异}
\begin{qa}{S0287}{经济学学硕保研为什么更看重数理基础和论文？}
因为“数理基础”写在申请条件里，而不是留给人猜。

厦门大学经济学科 2026 年经济学（学硕）夏令营申请条件第 2 条：生源限于国家高水平大学、国家重点大学或国家基础学科人才培养基地班的经济类及相关专业；或第四轮学科评估中理论经济学、应用经济学或统计学排名 A- 及以上高校的对应专业；或第二轮双一流建设高校和建设学科中的相关学科。该条末尾紧跟一句“具备良好的经济学基础和数理基础”。同批的金融硕士夏令营通知（之四）用同一套条件。

另一处硬证据在跨专业加试：南京大学商学院 2026 年接收推免实施细则规定，跨专业考生（本科专业与申请专业非同一一级学科且专业课程重合度低于 60\%）复试时加试本科主干，形式为笔试、时长 2 小时、总分 100 分；学术研究生（含推免硕士和直博生）考西方经济学和政治经济学，各 50 分。
\nref{S0314}
\nref{S0349}
\nref{S0282}
\end{qa}

\iin{厦门大学}
\ith{学科差异-经管}
\ist{院校与项目差异}
\begin{qa}{S0288}{管理科学与工程、信管保研的跨专业门槛高吗？}
门槛高低写在定义里，可以直接查。

南京大学商学院 2026 年接收推免工作实施细则给出两条明文界定：一是“工商管理学、管理科学与工程不视为跨专业”；二是“对金融硕士和金融学硕士而言，金融工程（工学学士或理学学士）不视为跨专业”。你本科是工商管理类或管理科学与工程类去申管科，就不触发加试。被认定为跨专业的，复试要加试本科主干课程，形式为笔试、时长 2 小时、总分 100 分。

另一端是很宽的白名单：北京大学信息管理系 2026 年夏令营（出版硕士、图书情报硕士项目）的申请条件是“全国重点高校编辑出版、信息管理、新闻传播、人文（中文、历史、哲学、考古）、理工（数学、计算机、物理、化学、生命科学等）或其他相关专业的大三在读本科生”。

吉林大学商学与管理学院 2027 年预报名通知则把管理科学与工程、信息资源管理、图书情报分设独立邮箱报名。
\nref{S0282}
\nref{S0035}
\nref{S0020}
\end{qa}

\iin{南京大学}
\ith{学科差异-经管}
\ist{院校与项目差异}
\section{学科差异-电子}
\begin{qa}{S0271}{电子信息类保研和计算机类保研在考核方式上有什么差别？}
差别不在“电子还是计算机”，而在具体学院有没有把机试写进必考环节。同为计算机学院，各校写法就不同：

\par\vspace{2pt}\nobreak\begin{center}\footnotesize
\begin{tabularx}{\textwidth}{>{\raggedright\arraybackslash}X>{\raggedright\arraybackslash}X}
\toprule
\textbf{学院（年份）} & \textbf{考核构成} \\ \midrule
武汉大学计算机学院 2027 年 & 复试综合成绩＝素质评估×20\%＋上机测试×40\%＋综合面试×40\% \\
贵州大学计算机科学与技术学院 2027 年 & 复试总成绩 200 分：上机考试 100 分、英语面试 20 分、综合面试 80 分 \\
北京师范大学人工智能学院 2027 年 & 综合考核成绩＝机试成绩＋面试成绩；机试考程序设计、数据结构与算法 \\
华中师范大学计算机学院 2026 年 & 线下面试，含综合能力面试与英语能力面试——不含上机 \\
\bottomrule
\end{tabularx}
\end{center}\vspace{2pt}\par

电子信息侧：浙江大学信息与电子工程学院 2026 年夏令营更强调直博意向与英语门槛（六级 480 分及以上、托福 80 或雅思 5.5），通知中未列上机环节。所以“电子信息一定不机试”是错的。
\nref{S0119}
\nref{S0002}
\nref{S0072}
\nref{S0105}
\nref{S0036}
\end{qa}

\iin{武汉大学}
\ith{学科差异-电子}
\ist{院校与项目差异}
\begin{qa}{S0272}{电子、微电子保研中，去中科院微电子所和去高校怎么比较？}
可比的是两家的申请门槛，不是平台好坏。

\par\vspace{2pt}\nobreak\begin{center}\footnotesize
\begin{tabularx}{\textwidth}{>{\raggedright\arraybackslash}X>{\raggedright\arraybackslash}X}
\toprule
\textbf{接收单位（年份）} & \textbf{门槛} \\ \midrule
中国科学院微电子研究所 2026 年“王守武”夏令营 & 微电子、集成电路、电子科学与技术、材料、光电信息、计算机、自动化、物理、数学等专业 2023 级本科生；专业成绩排名一般应为年级或专业前 20\%；英语通过四级；线上形式 \\
西安交通大学电信学部微电子学院 2026 年夏令营 & 所在学校为“双一流”大学建设高校，或所在专业在最新一轮学科评估中排 B+ 及以上；综合成绩排名专业前 20\% \\
浙江大学信息与电子工程学院 2026 年夏令营 & 优先选拔有直博意愿者；六级 480 分及以上（460 至 480 分有突出表现可考虑），或托福 80、雅思 5.5 及以上 \\
\bottomrule
\end{tabularx}
\end{center}\vspace{2pt}\par

微电子所英语只要求四级、但排名同样卡前 20\%；高校端普遍叠加学科评估或六级门槛。另外，北京大学软件与微电子学院 2026 年、山东大学集成电路学院 2026 年的夏令营通知可作进一步对照。
\nref{S0373}
\nref{S0144}
\nref{S0036}
\nref{S0214}
\nref{S0341}
\end{qa}

\iin{www.ime.ac.cn}
\ith{学科差异-电子}
\ist{院校与项目差异}
\begin{qa}{S0273}{通信、电子类保研中电子设计竞赛的加分权重有多大？}
先给结论：加分没有全国标准，而且各校普遍设上限。

中南大学材料科学与工程学院 2027 届预推荐实施细则规定：加分只认学校与学院重点赛事、排行榜重点赛事，或教育部等中央部委、国家一级学会、省教育厅、省科技厅组织的赛事，其他赛事不予加分；个人加分累计不超过 2 分。

湖南大学 2027 年推免工作通知把加分归入“专项评价成绩”，范围包含科研成果、省部级竞赛获奖、荣誉表彰、双创三大赛、体育艺术类比赛国家级及以上奖励等；该校数学学院 2027 年细则进一步限定专项评价成绩（一）不超过 3 分，具体标准见学院附件。

实际分值可从同校加分公示反推：挑战杯国家级特等奖排名第 1 记 6 分。

“电子设计竞赛值几分”要去查目标学院当年的重点赛事目录与加分上限；本次未找到直接写明电子设计竞赛分值的文件。
\nref{S0273}
\nref{S0540}
\nref{S0661}
\nref{S0541}
\end{qa}

\iin{中南大学}
\ith{学科差异-电子}
\ist{院校与项目差异}
\begin{qa}{S0291}{新闻传播保研中作品集和实习经历的权重有多大？}
找到的 2025 至 2026 年新闻传播院系文件中，申请材料清单里写的是“学术论文”，不是“作品集”。

山东大学新闻传播学院 2026 年“AI+”暑期学校暨夏令营公告的材料清单为：本科成绩单、外语成绩证明，以及“其它证明材料，包括能体现自身学术水平的代表性学术论文（最多提交 3 篇）”；同时对申请人设了专业或年级排名前 30\% 的门槛，学术硕士须通过英语六级 425 分或雅思 6.5 以上。

可对照的当届文件还有上海理工大学 2026 年新闻传播学（学硕）、新闻与传播（专硕）云端夏令营通知，浙江大学传媒与国际文化学院 2026 年夏令营通知，西安交通大学新闻与新媒体学院 2026 年夏令营通知。

上述文件中未出现“作品集”这一材料项，实习经历也未进入申请条件。若你所在院系有例外，以该院当年通知为准。
\nref{S0376}
\nref{S0081}
\nref{S0027}
\nref{S0403}
\end{qa}

\iin{山东大学}
\ith{学科差异-电子}
\ist{院校与项目差异}
\section{学科差异}
\begin{qa}{E0365}{文科保研和理工科保研在流程上有什么差别？}
流程骨架相同——预报名、资格审查、复试、推免服务系统确认——差别出在复试环节的名目上。

理工科院系常把机试写成必考：武汉大学计算机学院 2027 年细则的复试构成是素质评估 20\% 加上机测试 40\% 加综合面试 40\%；贵州大学计算机科学与技术学院 2027 年方案是上机考试 100 分加英语面试 20 分加综合面试 80 分。

文科院系写的是材料评估与面试：武汉大学文学院 2026 年推免生接收预报名通知写“考核包括背景评估、综合面试等，采取现场考核方式”；北京大学历史学系 2027 年推免复试由专家审核小组评审筛选后公布复试名单；南昌大学人文学院 2027 年复试细则把名额按专业细分（哲学 2 名、国际中文教育 6 名、中国史 5 名）。

同一所学校内部也会不同，武汉大学文学院与计算机学院就是两套。
\nref{S0119}
\nref{S0002}
\nref{S0097}
\nref{S0201}
\nref{S0016}
\end{qa}

\iin{武汉大学}
\ith{学科差异}
\ist{院校与项目差异}
\begin{qa}{S0359}{直博名额在不同学科中的分布差异有多大？}
理解分布差异的起点是一条全国统一的比例天花板。

教育部办公厅 2013 年《关于进一步加强推荐优秀应届本科毕业生免试攻读研究生工作的通知》：推免攻读直博生必须从已取得学术型推免生资格且获得所在高校推免生名额的优秀应届本科毕业生中遴选；2010 年以前已在教育部备案开展招收直博生工作的高校，招生专业一般为基础学科，招生人数原则上不超过本校博士生招生规模的 10\%；其他高校可在具有理、工、农、医学科门类国家重点一级学科的专业内招收直博生，招生人数一般不超过本专业招生人数的 20\%。

具体量级样本：武汉大学计算机学院 2027 年直博生约 50 人，硕士推免约 101 人；云南师范大学 2027 年直博推免名额 8 个；中国科学技术大学 2027 年工程硕博士培养改革专项在电子信息、机械、材料与化工等六个工程专业学位类别拟招 146 人（含工程博士）。
\nref{S0047}
\nref{S0119}
\nref{S0568}
\nref{S0673}
\end{qa}

\iin{中华人民共和国教育部}
\ith{学科差异}
\ist{院校与项目差异}
\begin{qa}{S0360}{专硕名额在不同学科中的占比差异有多大？}
有两个可以直接相除的样本，差距很大。

计算机方向：武汉大学计算机学院 2027 年接收推免实施细则写明，拟招收硕士推免生约 101 人（工程硕博国家专项除外），其中学术学位硕士推免生约 53 人、专业学位硕士推免生约 48 人——专硕约占该院硕士推免名额的一半（约 47.5\%）。

数学方向：北京师范大学数学科学学院 2026 年夏令营通知写明，拟招收报名学术型硕士和直博生营员共 200 人左右，专业型硕士（珠海校区培养）营员 30 人左右——专硕约占 13\%。

再有全国规则：教育部《2026 年全国硕士研究生招生工作管理规定》第六十六条，同一招生单位内学术学位招生计划可调整到专业学位使用，专业学位招生计划不得调整到学术学位专业使用，调整方向是单向的。

所以“专硕占比”跟着院系当年计划走，不能跨学科推。
\nref{S0119}
\nref{S0257}
\nref{S0051}
\end{qa}

\iin{武汉大学}
\ith{学科差异}
\ist{院校与项目差异}
\section{特殊身份}
\begin{qa}{S0307}{C9高校之间保研的相互认可度如何，跨C9难度大吗？}
能查到的唯一一处明文“C9 优待”在上海交通大学。

上海交通大学人工智能学院 2026 年优秀大学生夏令营报名通知：报名资格要求大学前三年成绩专业排名前 20\%，紧接着写“C9 高校或属于教育部基础学科拔尖学生培养计划 2.0 基地拔尖班的同学或在研究方面表现特别突出者（如发表高水平论文、竞赛获奖等），可适当放宽条件”。

这是“C9 可以放宽排名”的书面表述，不是“C9 之间互认”的制度安排。同校的上海交通大学计算机学院 2026 年夏令营活动通知另发一套条件，写法与人工智能学院不同，需要分别核对。

至于“跨 C9 难度大不大”，本次未找到任何可比对的录取数据，本手册不作断言。
\nref{S0514}
\nref{S0358}
\end{qa}

\iin{上海交通大学}
\ith{特殊身份}
\ist{院校与项目差异}
\begin{qa}{S0327}{专升本学生保研的可能性和路径是什么？}
这条在多校文件里是明文排除，不是“难度大”。

北京大学 2027 年关于推荐 2027 届优秀本科毕业生免试攻读研究生的办法（校本部）：申请推免的学生须为纳入国家普通本科招生计划录取的应届毕业生，并注明不含专升本、第二学士学位学生。

广东外语外贸大学 2026 年推免实施办法第二条：适用于学校全日制应届本科毕业生，不含辅修学位、专升本、第二学士学位学生、本硕项目的国际班学生和独立学院学生。

甘肃农业大学 2027 年遴选推荐管理办法：推荐对象为本校通过普通高考招生录取的全日制本科生，不含专升本、第二学士学位学生。

中国农业大学经济管理学院 2027 年接收推免细则：国家“双一流”建设高校应届本科生不含分校及专升本毕业生，其他推免资格高校同样不含。四校口径一致，本次未找到专升本成功推免的一手案例。
\nref{S0553}
\nref{S0134}
\nref{S0059}
\nref{S0684}
\end{qa}

\iin{北京大学}
\ith{特殊身份}
\ist{院校与项目差异}
\begin{qa}{S0328}{第二学士学位在读期间可以保研吗？}
在“推免推荐”这个环节，第二学士学位在读生同样被多校明文排除。

以下五份当届文件的适用对象范围里都写明不含第二学士学位学生：北京大学 2027 年推荐免试研究生办法（校本部）、广东外语外贸大学 2026 年推免实施办法第二条、甘肃农业大学 2027 年遴选推荐管理办法、山东中医药大学中医学院 2027 年遴选工作实施细则第二条、湖南大学数学学院 2027 年推免工作细则（注明适用对象不含第二学士学位学生、不含理科试验班）。

要区分的是另一条路径。教育部《2026 年全国硕士研究生招生工作管理规定》第十六条规定，报考法律（法学）专业学位的，须报名前所学专业为法学专业，获得法学第二学士学位的人员可报考——那是统考报考资格，不是推免推荐资格。

所以在读期间不能推免；毕业后走统考是另一套规则。
\nref{S0553}
\nref{S0134}
\nref{S0059}
\nref{S0639}
\nref{S0661}
\nref{S0051}
\end{qa}

\iin{北京大学}
\ith{特殊身份}
\ist{院校与项目差异}
\section{校区文凭}
\begin{qa}{E0382}{分校区（威海、深圳、苏州）的毕业证和本部一样吗？}
本次所查资料未找到任何说明分校区与本部毕业证、学位证关系的推免文件，检索“校区加毕业证”零命中，因此本手册不对“是否一样”作判断。

可查证的相邻事实：哈尔滨工业大学 2027 年接收优秀应届本科毕业生免试攻读研究生工作办法由校本部研究生院统一发布（2026 年 9 月 13 日），正文未对校区作分级说明，只在研招网页侧栏单列“威海校区”“深圳校区”两个入口链接；同校的报名通知（2026 年 7 月 19 日）另附各招生单位联系方式、推免生信息采集系统填报说明和 2027 年专项计划专业汇总表。

校区差异如果存在，会体现在各学院、各校区自己发布的接收细则里，而不在校级办法里。建议直接查目标校区当年发布的学院级细则，并要求招生办就证书样式作书面说明。
\nref{S0066}
\nref{S0552}
\nref{S0678}
\end{qa}

\iin{哈尔滨工业大学}
\ith{校区文凭}
\ist{院校与项目差异}
\begin{qa}{S0336}{哈尔滨工业大学（含深圳、威海校区）推免的校区差异在哪？}
哈尔滨工业大学 2027 年的文件体系是“校级办法加学院级细则”两层，校区差异不在这两份校级文件里。

校级：《哈尔滨工业大学 2027 年接收优秀应届本科毕业生免试攻读研究生工作办法》（2026 年 9 月 13 日）统一规定申请条件（定向生须提交合同单位书面同意书）、材料审核专家组每组不少于 3 人、面试专家组一般不少于 5 人、面试采用三随机原则、每名考生面试时间一般不少于 20 分钟；业务考核以面试为主，各专业还可增加笔试、实验考核、实践能力考核等环节。

报名侧：《哈尔滨工业大学关于 2027 年接收优秀应届本科毕业生免试攻读研究生报名的通知》（2026 年 7 月 19 日）另发，附各招生单位联系方式、信息采集系统填报说明与专项计划专业汇总表。要判断校区差异，需逐院查当届细则。
\nref{S0066}
\nref{S0678}
\end{qa}

\iin{哈尔滨工业大学}
\ith{校区文凭}
\ist{院校与项目差异}
\section{名额分配}
\begin{qa}{E0371}{本校保研和外校保研的名额怎么分？}
这条有明确的政策沿革，而且限额已经被取消。

现行（2014 年起）：教育部办公厅《关于进一步完善推荐优秀应届本科毕业生免试攻读研究生工作办法的通知》规定，所有推免名额（除有特殊政策要求的专项计划外）均可向其他招生单位推荐，推荐高校不得将报考本校作为遴选推免生的条件；2014 年起教育部下达推免名额时不再区分学术学位和专业学位，不再设置留校限额。

历史（2013 年及以前）：《全国普通高等学校推荐优秀应届本科毕业生免试攻读硕士学位研究生工作管理办法（试行）》第十七条规定，设有研究生院的高等学校接收本校推免生的人数不得超过本校推免生总数的 65\%，地处西部省份或军工、矿业、石油、地质、农林等特殊类型的高等学校可适当放宽，但不得超过 75\%。

所以今天这个问题不由政策配额决定，由你填志愿和接收单位录取结果决定。
\nref{S0048}
\nref{S0454}
\nref{S0047}
\end{qa}

\iin{中华人民共和国教育部}
\ith{名额分配}
\ist{院校与项目差异}
